\section{Differential invariants. Differentially smooth morphisms}
\label{section:IV.16}

\oldpage[IV-4]{5}
In this section we present, in global form, some notions of
differential calculus that are particularly useful in Algebraic Geometry. We pass
over in silence many developments that are classical in Differential Geometry (connections,
infinitesimal transformations associated to a vector field, jets, etc.), although these
notions can be expressed particularly naturally in the framework of schemes. We
also pass over in silence here the phenomena special to characteristic $p>0$
(some of which are studied, in the affine setting, in (0, 21)). For some further
material on the differential formalism in preschemes, the reader may consult
Exposés II and VII of [42], as well as the later chapters of this Treatise.

\subsection{Normal invariants of an immersion}
\EGAReaderTarget{ega4.c04.subsec.16-1}{IV.16.1}

\begin{env}[16.1.1]
\label{IV.16.1.1}
Let $(X,\mathcal{O}_X)$, $(Y,\mathcal{O}_Y)$ be two ringed spaces and
$f=(\psi,\theta):Y\to X$ a morphism of ringed spaces (0\textsubscript{I}, 4.1.1)
such that the homomorphism
\[
  \theta^\#:\psi^*(\mathcal{O}_X)\longrightarrow\mathcal{O}_Y
\]
is surjective, so that $\mathcal{O}_Y$ is identified with a quotient sheaf of rings
$\psi^*(\mathcal{O}_X)/\mathscr{I}_f$. We can then endow
$\psi^*(\mathcal{O}_X)$ with the $\mathscr{I}_f$-preadic filtration.
\end{env}

\begin{definition}[16.1.2]\label{IV.16.1.2}
\EGAReaderTarget{ega4.c04.def.16-1-2}{IV.16.1.2}
The $\mathcal{O}_Y$-augmented sheaf of rings
$\psi^*(\mathcal{O}_X)/\mathscr{I}_f^{n+1}$ is called the $n$-th normal invariant
of $f$; the ringed space
$(Y,\psi^*(\mathcal{O}_X)/\mathscr{I}_f^{n+1})$ is called the $n$-th infinitesimal
neighbourhood of $Y$ for the morphism $f$, and is denoted by $Y_f^{(n)}$ or simply
$Y^{(n)}$. The sheaf of graded rings associated to the sheaf of filtered rings
$\psi^*(\mathcal{O}_X)$
\[
  \mathscr{G}r_\bullet(f)=\bigoplus_{n\geq 0}
  (\mathscr{I}_f^n/\mathscr{I}_f^{n+1})
  \tag{16.1.2.1}
  \label{IV.16.1.2.1}
\]
is called the sheaf of graded rings associated to $f$. The sheaf
$\mathscr{G}r_1(f)=\mathscr{I}_f/\mathscr{I}_f^2$ is called the conormal sheaf
of $f$ (which is also denoted by $\mathscr{N}_{Y/X}$ if no confusion results).
\end{definition}

It is clear that the
$\mathcal{O}_{Y^{(n)}}=\psi^*(\mathcal{O}_X)/\mathscr{I}_f^{n+1}$
(which are also denoted by $\mathcal{O}_{Y_f^{(n)}}$) form a
\oldpage[IV-4]{6}
projective system of sheaves of rings on $Y$, the transition homomorphism
$\varphi_{nm}:\mathcal{O}_{Y^{(m)}}\to\mathcal{O}_{Y^{(n)}}$ for $n\leq m$
identifying $\mathcal{O}_{Y^{(n)}}$ with the quotient of $\mathcal{O}_{Y^{(m)}}$ by the
power $(\mathscr{I}_f/\mathscr{I}_f^{n+1})^m$ of the augmentation ideal of
$\mathcal{O}_{Y^{(n)}}$, the kernel of
$\varphi_{0n}:\mathcal{O}_{Y^{(n)}}\to\mathcal{O}_Y$. The $Y^{(n)}$ consequently
form an inductive system of ringed spaces, all having as underlying space
the space $Y$, and we have canonical morphisms of ringed spaces
$h_n:Y^{(n)}\to X$ equal to $(\psi,\theta_n)$, where $\theta_n^\#$ is the
canonical morphism
$\psi^*(\mathcal{O}_X)\to\psi^*(\mathcal{O}_X)/\mathscr{I}_f^{n+1}$.
It is clear that the sheaf $\mathscr{G}r_\bullet(f)$ is a sheaf of graded
algebras over the sheaf of rings
$\mathcal{O}_Y=\mathscr{G}r_0(f)$, and the $\mathscr{G}r_k(f)$ are
$\mathcal{O}_Y$-modules.

As for every filtered sheaf of rings, we have a canonical surjective
homomorphism of graded $\mathcal{O}_Y$-algebras
\[
  \mathscr{S}_{\mathcal{O}_Y}^{\bullet}(\mathscr{G}r_1(f))
  \longrightarrow\mathscr{G}r_\bullet(f)
  \tag{16.1.2.2}
  \label{IV.16.1.2.2}
\]
which coincides in degrees 0 and 1 with the identity homomorphisms.

\begin{examples}[16.1.3]
\label{IV.16.1.3}
\begin{enumerate}
\item[(i)] Suppose that $X$ is a ringed space in local rings, that $Y$ is
reduced to a single point $y$ (endowed with a ring $\mathcal{O}_y$), and that, if
$x=\psi(y)$, $\theta^\#:\mathcal{O}_x\to\mathcal{O}_y$ is a surjective
ring homomorphism having as kernel the maximal ideal $\mathfrak{m}_x$ of
$\mathcal{O}_x$. Then the $\mathcal{O}_{Y^{(n)}}$ are identified with the rings
$\mathcal{O}_x/\mathfrak{m}_x^{n+1}$, and $\mathscr{G}r_\bullet(f)$ with the graded
ring associated to the local ring $\mathcal{O}_x$ endowed with its
$\mathfrak{m}_x$-preadic filtration.
\item[(ii)] Suppose that $Y$ is a closed subset of an open subspace $U$
of $X$ and that $\mathcal{O}_Y$ is induced on $Y$ by a quotient sheaf
$\mathcal{O}_U/\mathscr{I}$, where $\mathscr{I}$ is an ideal of $\mathcal{O}_U$
such that $\mathscr{I}_x=\mathcal{O}_x$ for every $x\notin Y$; if $X$ is a
ringed space in local rings, we will suppose moreover that
$\mathscr{I}_x\neq\mathcal{O}_x$ for $x\in Y$, so that
$(Y,\mathcal{O}_Y)$ is again a ringed space in local rings.

Let $\psi_0:Y\to U$ be the canonical injection, and denote by
$\theta_0:\mathcal{O}_U\to(\psi_0)_*(\mathcal{O}_Y)$ the homomorphism such that
$\theta_0^\#$ is the canonical homomorphism
$\psi_0^*(\mathcal{O}_U)=\mathcal{O}_U|Y\to
(\mathcal{O}_U/\mathscr{I})|Y$, so that
$j_0=(\psi_0,\theta_0):Y\to U$ is a morphism of ringed spaces (and of ringed
spaces in local rings if $X$ is a ringed space in local rings); if
$i:U\to X$ is the canonical injection (a morphism of ringed spaces),
$j=i\circ j_0$ is the morphism $(\psi,\theta)$ from $Y$ to $X$, where
$\psi:Y\to X$ is the canonical injection and
$\theta:\mathcal{O}_X\to\psi_*(\mathcal{O}_Y)$ is the homomorphism such that
$\theta^\#=\theta_0^\#$. Since $\theta^\#$ is surjective, we can apply the
preceding definitions; $\mathcal{O}_{Y^{(n)}}$ is equal to
$\psi_0^*(\mathcal{O}_U/\mathscr{I}^{n+1})$, and we have
$(\psi_0)_*(\mathcal{O}_{Y^{(n)}})=\mathcal{O}_U/\mathscr{I}^{n+1}$ and
$\mathscr{G}r_n(j)=\mathscr{G}r_n(j_0)=
\psi_0^*(\mathscr{I}^n/\mathscr{I}^{n+1})=
j_0^*(\mathscr{I}^n/\mathscr{I}^{n+1})$.
\end{enumerate}
\end{examples}

\begin{env}[16.1.4]
\label{IV.16.1.4}
Example (16.1.3, (ii)) shows that in general the
$\mathcal{O}_{Y^{(n)}}$ are not canonically endowed with a structure of
$\mathcal{O}_Y$-module, nor \emph{a fortiori} with a structure of
$\mathcal{O}_Y$-algebra. Giving such a structure is equivalent to giving a
homomorphism of sheaves of rings
$\lambda_n:\mathcal{O}_Y\to\mathcal{O}_{Y^{(n)}}$, right inverse to
the augmentation homomorphism $\varphi_{0n}$; this also amounts to giving
a morphism of ringed spaces
$(1_Y,\lambda_n):Y^{(n)}\to Y$, left inverse to the canonical morphism
$(1_Y,\varphi_{0n}):Y\to Y^{(n)}$.
\end{env}

\begin{proposition}[16.1.5]\label{IV.16.1.5}
\EGAReaderTarget{ega4.c04.prop.16-1-5}{IV.16.1.5}
Let $f=(\psi,\theta):Y\to X$ be an immersion of preschemes. Then:
\begin{enumerate}
\item[(i)] $\mathscr{G}r_\bullet(f)$ is a quasi-coherent graded
$\mathcal{O}_Y$-algebra.
\oldpage[IV-4]{7}
\item[(ii)] The $Y^{(n)}$ are preschemes, canonically isomorphic to
sub-preschemes of $X$.
\item[(iii)] Every homomorphism of sheaves of rings
$\lambda_n:\mathcal{O}_Y\to\mathcal{O}_{Y^{(n)}}$ right inverse to
the augmentation homomorphism $\varphi_{0n}$ makes
$\mathcal{O}_{Y^{(n)}}$ and the $\mathcal{O}_{Y^{(k)}}$ for $k\leq n$ into
quasi-coherent $\mathcal{O}_Y$-algebras; the
$\mathcal{O}_Y$-module structures deduced from the preceding structures on the
$\mathscr{G}r_k(f)$ for $k\leq n$ coincide with those defined in
(16.1.2).
\end{enumerate}
\end{proposition}

\noindent(i) The question being local on $X$ and on $Y$, we may restrict to the case
where $Y$ is a closed sub-prescheme of $X$ defined by a quasi-coherent ideal
$\mathscr{I}$ of $\mathcal{O}_X$; since $\mathcal{O}_Y$ is the restriction to
$Y$ of $\mathcal{O}_X/\mathscr{I}$, assertion (i) is evident, and $Y^{(n)}$
is the closed sub-prescheme of $X$ defined by the quasi-coherent ideal
$\mathscr{I}^{n+1}$ of $\mathcal{O}_X$. Finally, to prove (iii), note that
giving $\lambda_n$ makes the ideal
$\mathscr{I}/\mathscr{I}^{n+1}$ of the augmentation $\varphi_{0n}$ and its
quotients $\mathscr{I}/\mathscr{I}^{k+1}$ ($1\leq k\leq n$) into
$\mathcal{O}_Y$-modules, and it suffices to prove by induction on $k$ that the
$\mathscr{I}/\mathscr{I}^{k+1}$ are quasi-coherent $\mathcal{O}_Y$-modules
and that the quotient $\mathcal{O}_Y$-module structure that we
deduce from this on $\mathscr{I}^k/\mathscr{I}^{k+1}$ is the same as the one defined in
(16.1.2). The second assertion is immediate,
$\mathscr{I}^k/\mathscr{I}^{k+1}$ being annihilated by
$\mathscr{I}/\mathscr{I}^{n+1}$; the first follows, by induction on $k$,
from the fact that it is trivial for $k=1$ and that
$\mathscr{I}/\mathscr{I}^{k+1}$ is an extension of
$\mathscr{I}/\mathscr{I}^k$ by
$\mathscr{I}^k/\mathscr{I}^{k+1}$ (III, 1.4.17).

\begin{corollary}[16.1.6]
\label{IV.16.1.6}
Under the general hypotheses of (16.1.5), if the immersion $f$ is locally of
finite presentation, the $\mathscr{G}r_n(f)$ are quasi-coherent
$\mathcal{O}_Y$-modules of finite type.
\end{corollary}

Indeed, with the notation of the proof of (16.1.5), $\mathscr{I}$ is an
ideal of finite type of $\mathcal{O}_X$ (1.4.7), so the
$\mathscr{I}^n/\mathscr{I}^{n+1}$ are $\mathcal{O}_Y$-modules of finite type,
whence the conclusion.

\begin{corollary}[16.1.7]
\label{IV.16.1.7}
Under the general hypotheses of (16.1.5), let $g:X\to Y$ be a morphism of
preschemes left inverse to $f$. Then, for every $n$, the composite morphism
$(1,\lambda_n):Y^{(n)}\xrightarrow{h_n}X\xrightarrow{g}Y$ defines a
homomorphism of sheaves of rings
$\lambda_n:\mathcal{O}_Y\to\mathcal{O}_{Y^{(n)}}$ right inverse to
the augmentation $\varphi_{0n}$, making $\mathcal{O}_{Y^{(n)}}$ a
quasi-coherent $\mathcal{O}_Y$-algebra; for these homomorphisms, the
transition homomorphisms
$\varphi_{nm}:\mathcal{O}_{Y^{(m)}}\to\mathcal{O}_{Y^{(n)}}$ ($n\leq m$) are
homomorphisms of $\mathcal{O}_Y$-algebras. Moreover, if $g$ is locally
of finite type, the $\mathcal{O}_{Y^{(n)}}$ are quasi-coherent
$\mathcal{O}_Y$-modules of finite type.
\end{corollary}

The first assertion follows immediately from the definitions and from (16.1.5). On the
other hand, if $g$ is locally of finite type, then $f$ is locally of
finite presentation (1.4.3, (v)); the $\mathscr{G}r_n(f)$ then being
quasi-coherent $\mathcal{O}_Y$-modules of finite type by (16.1.6), the same holds
for the $\mathcal{O}_Y$-modules
$\mathscr{I}/\mathscr{I}^{n+1}$, which are extensions of finitely many
$\mathscr{G}r_k(f)$ (III, 1.4.17).

\begin{proposition}[16.1.8]
\label{IV.16.1.8}
Let $X$ be a locally Noetherian prescheme and $j:Y\to X$ an immersion. Then
the $Y^{(n)}$ are locally Noetherian preschemes, the
$\mathscr{G}r_n(j)$ are coherent $\mathcal{O}_Y$-modules, and
$\mathscr{G}r_\bullet(j)$ is a coherent sheaf of rings on the space $Y$.
\end{proposition}

Everything being local on $X$ and $Y$, we are reduced to the case where $X$ is affine
and $j$ is a closed immersion, and then all the assertions are evident except the
last, which follows from the fact that if $A$ is a Noetherian ring and
$\mathfrak{I}$ an ideal of $A$, $\operatorname{gr}_{\mathfrak{I}}^\bullet(A)$
is a Noetherian ring, taking into account the exactness of the functor $\psi^*$ and
(0\textsubscript{I}, 5.3.7).

\begin{proposition}[16.1.9]
\label{IV.16.1.9}
\oldpage[IV-4]{8}
Let $X$ be a prescheme, $j:Y\to X$ an immersion locally of finite
presentation, and $y$ a point of $Y$. The following conditions are equivalent:
\begin{enumerate}
\item[(a)] There exists an open neighbourhood $U$ of $y$ in $Y$ such that $j|U$ is
a homeomorphism of $U$ onto an open subset of $X$.
\item[(b)] There exists an integer $n>0$ such that the canonical homomorphism
\[
  (\varphi_{n,n-1})_y:
  \mathcal{O}_{Y^{(n)},y}\longrightarrow\mathcal{O}_{Y^{(n-1)},y}
\]
is bijective.
\item[(c)] There exists an integer $n>0$ such that
$(\mathscr{G}r_n(j))_y=0$.
\end{enumerate}
Moreover, if the integer $n$ satisfies b) or c), there exists a neighbourhood $V$ of $y$
in $Y$ such that $\mathscr{G}r_m(j)|V=0$ for $m\geq n$ and
$\varphi_{nm}|V:\mathcal{O}_{Y^{(m)}}|V\to\mathcal{O}_{Y^{(n)}}|V$ is
bijective for $m\geq n$.
\end{proposition}

The question being local on $Y$, we may restrict to the case where $j$ is a
closed immersion, $Y$ being defined by a quasi-coherent ideal of finite type
$\mathscr{I}$ of $\mathcal{O}_X$. The equivalence of b) and c), for a given $n$,
is then immediate; moreover, since
$\mathscr{I}^n/\mathscr{I}^{n+1}$ is an $\mathcal{O}_X$-module of finite type,
there exists an open neighbourhood $U$ of $y$ in $X$ such that
$\mathscr{I}^n|U=\mathscr{I}^{n+1}|U$ (0\textsubscript{I}, 5.2.2), hence also
$\mathscr{I}^n|U=\mathscr{I}^m|U$ for $m\geq n$, which proves the last
assertions. To prove that a) implies b), we may restrict to the case where
the underlying space of $Y$ is equal to the underlying space of $X$, and where
$\mathscr{I}$ is generated by finitely many of its sections over $X$:
since $\mathscr{I}$ is then contained in the nilradical $\mathscr{N}$ of
$\mathcal{O}_X$ (I, 5.1.2), it is nilpotent, which proves b). Finally, to
prove that b) implies a), we may also restrict to the case where
$\mathscr{I}^n=\mathscr{I}^{n+1}$; then, for every $y\in Y$, since
$\mathscr{I}_y\subset\mathfrak{m}_y$, the maximal ideal of
$\mathcal{O}_{X,y}$, we necessarily have $\mathscr{I}_y^n=0$ by virtue of
Nakayama's lemma, since $\mathscr{I}_y$ is an ideal of finite type.
The set of $x\in X$ such that $\mathscr{I}_x^n=0$ is therefore an open subset $U$ of
$X$ containing $Y$ (0\textsubscript{I}, 5.2.2); since on the other hand
$\mathscr{I}_x\neq 0$ for $x\notin Y$, we necessarily have $U=Y$.

\begin{corollary}[16.1.10]
\label{IV.16.1.10}
For the restriction of the immersion $j$ to a neighbourhood of $y$ in $Y$ to be
an open immersion (in other words, for $j$ to be a local isomorphism at the
point $y$), it is necessary and sufficient that
$(\mathscr{G}r_1(j))_y=(\mathscr{N}_{Y/X})_y=0$.
\end{corollary}

The condition is evidently necessary, and the preceding argument, applied
for $n=1$, proves that it is sufficient.

\begin{remarks}[16.1.11]\label{IV.16.1.11}
\EGAReaderTarget{ega4.c04.rem.16-1-11}{IV.16.1.11}
\begin{enumerate}
\item[(i)] Under the conditions of the definition (16.1.1), the projective limit
of the projective system
$(\mathcal{O}_{Y^{(n)}},\varphi_{nm})$ of sheaves of rings on $Y$ is
called the normal invariant of infinite order of $f$, and is sometimes denoted by
$\mathcal{O}_{Y^{(\infty)}}$. When $X$ is a locally
Noetherian prescheme and $j:Y\to X$ a closed immersion, $Y$ thus being a closed
sub-prescheme of $X$ defined by a coherent ideal $\mathscr{I}$,
$\mathcal{O}_{Y^{(\infty)}}$ is none other than the formal completion of
$\mathcal{O}_X$ along $Y$ (I, 10.8.4), and
$Y^{(\infty)}=(Y,\mathcal{O}_{Y^{(\infty)}})$ the formal prescheme completion of
$X$ along $Y$ (I, 10.8.5). In all cases, we may say that
$Y^{(\infty)}$ is the formal neighbourhood of $Y$ in $X$ (for the morphism
$f$). In the particular case just considered, it is therefore the formal
prescheme that is the inductive limit of the infinitesimal neighbourhoods of order $n$.
\item[(ii)] Note that for a morphism of preschemes
$f=(\psi,\theta):Y\to X$, it can happen that the homomorphism
$\theta^\#:\psi^*(\mathcal{O}_X)\to\mathcal{O}_Y$ is surjective without $f$
being a local
\oldpage[IV-4]{9}
immersion, and without $f$ being injective. An example is obtained by taking for $Y$
a sum of preschemes $Y_\lambda$ all isomorphic to
$\operatorname{Spec}(\mathcal{O}_x)$, where $x\in X$, and for $f$ the morphism
equal to the canonical morphism on each of the $Y_\lambda$.
\end{enumerate}
\end{remarks}

\subsection{Functorial properties of the normal invariants of an immersion}
\EGAReaderTarget{ega4.c04.subsec.16-2}{IV.16.2}

\begin{env}[16.2.1]
\label{IV.16.2.1}
Let $f=(\psi,\theta):Y\to X$,
$f'=(\psi',\theta'):Y'\to X'$ be two morphisms of ringed spaces such that the
homomorphisms $\theta^\#$ and $\theta'^\#$ are surjective; consider a
commutative diagram of morphisms of ringed spaces
\[
  \xymatrix{
    Y \ar[r]^f & X \\
    Y' \ar[r]_{f'} \ar[u]^u & X' \ar[u]_v
  }
  \tag{16.2.1.1}
  \label{IV.16.2.1.1}
\]
Set $u=(\rho,\lambda)$, $v=(\sigma,\mu)$. We have
$\rho^*(\psi^*(\mathcal{O}_X))=
\psi'^*(\sigma^*(\mathcal{O}_X))$, and consequently we have a commutative diagram
of homomorphisms of sheaves of rings on $Y'$
\[
  \xymatrix{
    \rho^*(\psi^*(\mathcal{O}_X))
      =\psi'^*(\sigma^*(\mathcal{O}_X))
      \ar[r]^-{\psi'^*(\mu^\#)}
      \ar[d]_{\rho^*(\theta^\#)}
      & \psi'^*(\mathcal{O}_{X'}) \ar[d]^{\theta'^\#} \\
    \rho^*(\mathcal{O}_Y) \ar[r]_-{\lambda^\#}
      & \mathcal{O}_{Y'}
  }
\]
from which we conclude, if $\mathscr{I}$ and $\mathscr{I}'$ are the kernels of
$\theta^\#$ and $\theta'^\#$, that
$\psi'^*(\mu^\#)(\rho^*(\mathscr{I}))\subset\mathscr{I}'$, taking into account
the exactness of the functor $\rho^*$. We deduce immediately that for every integer
$n$,
$\psi'^*(\mu^\#)(\rho^*(\mathscr{I}^n))\subset\mathscr{I}'^n$, which shows
that $\psi'^*(\mu^\#)$ defines, by passage to quotients, a homomorphism of
sheaves of rings
\[
  \nu_n:
  \rho^*(\psi^*(\mathcal{O}_X)/\mathscr{I}^{n+1})
  \longrightarrow
  \psi'^*(\mathcal{O}_{X'})/\mathscr{I}'^{n+1}
  \tag{16.2.1.2}
  \label{IV.16.2.1.2}
\]
and consequently a morphism of ringed spaces
$w_n=(\rho,\nu_n):Y'^{(n)}\to Y^{(n)}$ (which, for $n=0$, is none other than
$u$). It follows immediately from this definition that the diagrams
\[
  \xymatrix@R=1pc{
    Y^{(n)} \ar[r]^-{h_{mn}} & Y^{(m)} \ar[r]^-{h_m} & X \\
    & & & (n\leq m) \\
    Y'^{(n)} \ar[r]_-{h'_{mn}} \ar[uu]^-{w_n}
      & Y'^{(m)} \ar[r]_-{h'_m} \ar[uu]^-{w_m}
      & X' \ar[uu]_-v
  }
\]
(where the horizontal arrows are the canonical morphisms (16.1.2)) are
commutative.

By passage to quotients from the homomorphisms (16.2.1.2), and taking
\oldpage[IV-4]{10}
into account the exactness of the functor $\rho^*$, we obtain a di-homomorphism
of graded algebras (relative to the homomorphism
$\lambda^\#:\rho^*(\mathcal{O}_Y)\to\mathcal{O}_{Y'}$)
\[
  \operatorname{gr}(u):
  \rho^*(\mathscr{G}r_\bullet(f))
  \longrightarrow\mathscr{G}r_\bullet(f')
  \tag{16.2.1.3}
  \label{IV.16.2.1.3}
\]
(or, if one prefers, a $\rho$-morphism (0\textsubscript{I}, 3.5.1)
$\mathscr{G}r_\bullet(f)\to\mathscr{G}r_\bullet(f')$), and in particular a
di-homomorphism of the conormal sheaves
\[
  \operatorname{gr}_1(u):
  \rho^*(\mathscr{G}r_1(f))\longrightarrow\mathscr{G}r_1(f').
\]
It is moreover immediate that these homomorphisms give rise to a commutative
diagram
\[
  \xymatrix{
    \rho^*(\mathscr{S}_{\mathcal{O}_Y}^{\bullet}
      (\mathscr{G}r_1(f)))
      \ar[r] \ar[d]_-{\mathscr{S}(\operatorname{gr}_1(u))}
      & \rho^*(\mathscr{G}r_\bullet(f))
        \ar[d]^-{\operatorname{gr}(u)} \\
    \mathscr{S}_{\mathcal{O}_{Y'}}^{\bullet}
      (\mathscr{G}r_1(f'))
      \ar[r]
      & \mathscr{G}r_\bullet(f')
  }
  \tag{16.2.1.4}
  \label{IV.16.2.1.4}
\]
where the horizontal arrows are the canonical homomorphisms (16.1.2.2).

Finally, if we have a commutative diagram of morphisms of ringed spaces
\[
  \xymatrix{
    Y \ar[r]^f & X \\
    Y' \ar[r]_{f'} \ar[u]^u & X' \ar[u]_v \\
    Y'' \ar[r]_{f''} \ar[u]^{u'} & X'' \ar[u]_{v'}
  }
\]
where $f''=(\psi'',\theta'')$ is such that $\theta''^\#$ is surjective, and if
$w'_n$ and $w''_n$ are defined from $u'$, $v'$ on the one hand, and from
$u''=u\circ u'$, $v''=v\circ v'$ on the other, then we have
$w''_n=w_n\circ w'_n$, as follows immediately from the definitions and from
(0\textsubscript{I}, 3.5.5); likewise we have
$\operatorname{gr}(u'')=
\operatorname{gr}(u')\circ\rho'^*(\operatorname{gr}(u))$ if
$u'=(\rho',\lambda')$. We can therefore say that the $Y^{(n)}$ and
$\mathscr{G}r_\bullet(f)$ depend functorially on $f$.
\end{env}

\begin{proposition}[16.2.2]
\label{IV.16.2.2}
With the notation and hypotheses of (16.2.1), suppose moreover that $f$, $f'$,
$u$ and $v$ are morphisms of preschemes. Then:
\begin{enumerate}
\item[(i)] The morphisms $w_n:Y'^{(n)}\to Y^{(n)}$ are morphisms of
preschemes.
\item[(ii)] If $Y'=Y\times_X X'$, $u$ and $f'$ being the canonical
projections, and if $f$ is an immersion or $v$ is flat, we have
$Y'^{(n)}=Y^{(n)}\times_X X'$.
\item[(iii)] If $Y'=Y\times_X X'$ and if $v$ is flat (resp. if $f$ is an
immersion), the homomorphism
\[
  \operatorname{Gr}(u)=\operatorname{gr}(u)\otimes 1:
  \mathscr{G}r_\bullet(f)\otimes_{\mathcal{O}_Y}\mathcal{O}_{Y'}
  \longrightarrow\mathscr{G}r_\bullet(f')
\]
is bijective (resp. surjective).
\end{enumerate}
\end{proposition}

\noindent(i) The hypotheses immediately imply that for every $y'\in Y'$,
$\rho_{y'}^*(\theta_{\psi'(y')}^\#)$ is a local homomorphism (I, 1.6.2),
so $w_n$ is a morphism of preschemes (I, 2.2.1).

\oldpage[IV-4]{11}
\noindent(ii) and (iii) If $f$ is an immersion, we may restrict to the case where
$f$ is a closed immersion, $Y$ thus being defined by the quasi-coherent
ideal $\mathscr{I}$ of $\mathcal{O}_X$, and $Y^{(n)}$ by the ideal
$\mathscr{I}^{n+1}$; the assertions then follow from (I, 4.4.5).

Suppose secondly that $v$ is flat; we may restrict to the case where
$X=\operatorname{Spec}(A)$, $Y=\operatorname{Spec}(B)$,
$X'=\operatorname{Spec}(A')$ are affine, $A'$ being a flat $A$-module;
then $Y'=\operatorname{Spec}(B')$ with $B'=B\otimes_A A'$; moreover, if
$\mathfrak{I}$ is the kernel of the homomorphism $A\to B$, the kernel
$\mathfrak{I}'$ of $A'\to B'$ is identified with $\mathfrak{I}\otimes_A A'$ by
flatness, and
\[
  \mathfrak{I}'^n/\mathfrak{I}'^{n+1}
  =(\mathfrak{I}^n/\mathfrak{I}^{n+1})\otimes_A A'.
\]
We deduce immediately, taking into account (0\textsubscript{I}, 4.3.3), that the
$\mathcal{O}_{Y'}$-module
$\mathscr{I}'^n/\mathscr{I}'^{n+1}$ is equal to
\begin{align*}
&\psi'^*\!\left(\sigma^*
  \left((\mathfrak{I}^n/\mathfrak{I}^{n+1})^{\sim}\right)\right)
  \otimes_{\sigma^*(\mathcal{O}_X)}\mathcal{O}_{X'} \\
={}&
\psi'^*\!\left(\sigma^*
  \left((\mathfrak{I}^n/\mathfrak{I}^{n+1})^{\sim}\right)\right)
  \otimes_{\psi'^*(\sigma^*(\mathcal{O}_X))}
  \psi'^*(\mathcal{O}_{X'}) \\
={}&
\rho^*(\mathscr{I}^n/\mathscr{I}^{n+1})
  \otimes_{\rho^*(\psi^*(\mathcal{O}_X))}
  \psi'^*(\mathcal{O}_{X'}),
\end{align*}
and since in particular for $n=0$ we have
\[
  \mathcal{O}_{Y'}=
  \rho^*(\mathcal{O}_Y)
  \otimes_{\rho^*(\psi^*(\mathcal{O}_X))}
  \psi'^*(\mathcal{O}_{X'}),
\]
we obtain a canonical isomorphism of
$\mathscr{I}'^n/\mathscr{I}'^{n+1}$ onto
\[
  \rho^*(\mathscr{I}^n/\mathscr{I}^{n+1})
  \otimes_{\rho^*(\mathcal{O}_Y)}\mathcal{O}_{Y'}
  =
  (\mathscr{I}^n/\mathscr{I}^{n+1})
  \otimes_{\mathcal{O}_Y}\mathcal{O}_{Y'},
\]
which proves (iii). Now set
$C_n=\Gamma(Y,\mathcal{O}_{Y^{(n)}})$,
$C'_n=\Gamma(Y',\mathcal{O}_{Y'^{(n)}})$. Since $Y^{(n)}$ and
$Y'^{(n)}$ are affine schemes (16.1.5), the kernel
$\mathfrak{K}_n$ (resp. $\mathfrak{K}'_n$) of the homomorphism
$C_n\to C_{n-1}$ (resp. $C'_n\to C'_{n-1}$) is
$\Gamma(Y,\mathscr{I}^n/\mathscr{I}^{n+1})$ (resp.
$\Gamma(Y',\mathscr{I}'^n/\mathscr{I}'^{n+1})$); we therefore deduce from the
above that $\mathfrak{K}'_n=\mathfrak{K}_n\otimes_A A'$. Now, we have a
commutative diagram
\[
  \xymatrix{
    0 \ar[r] &
    \mathfrak{K}_n\otimes_A A' \ar[r]\ar[d]^r &
    C_n\otimes_A A' \ar[r]\ar[d]^{s_n} &
    C_{n-1}\otimes_A A' \ar[r]\ar[d]^{s_{n-1}} & 0 \\
    0 \ar[r] &
    \mathfrak{K}'_n \ar[r] &
    C'_n \ar[r] &
    C'_{n-1} \ar[r] & 0
  }
\]
where the left vertical arrow is bijective and the two rows are exact
($A'$ being a flat $A$-module). We deduce by induction that $s_n$ is
bijective for every $n$, since this is true by hypothesis for $n=0$, and follows
for arbitrary $n$ by application of the five lemma. This proves the second
assertion of (ii).

\begin{corollary}[16.2.3]
\label{IV.16.2.3}
Let $g:X\to Y$, $u:Y'\to Y$ be two morphisms of preschemes,
$X'=X\times_Y Y'$, and $g':X'\to Y'$ and $v:X'\to X$ the canonical projections.
Let $f:Y\to X$ be a $Y$-section of $X$ (hence an immersion), and
$f'=f_{(Y')}:Y'\to X'$ the $Y'$-section of $X'$ deduced from $f$ by the
base change $u$. Then:
\begin{enumerate}
\item[(i)] The morphism
$w_n:Y_{f'}'^{(n)}\to Y_f^{(n)}$ corresponding to $f$, $f'$, $u$, $v$
(16.2.1) and the canonical morphism
$h'_n:Y_{f'}'^{(n)}\to X'$ identify $Y_{f'}'^{(n)}$ with the product
$Y_f^{(n)}\times_X X'$.
\item[(ii)] If we endow $\mathcal{O}_{Y_f^{(n)}}$ (resp.
$\mathcal{O}_{Y_{f'}'^{(n)}}$) with the $\mathcal{O}_Y$-algebra structure
defined by $g$ (resp. with the $\mathcal{O}_{Y'}$-algebra structure defined
by $g'$) (16.1.6), the homomorphism of $\mathcal{O}_{Y'}$-algebras
\[
  \rho^*(\mathcal{O}_{Y_f^{(n)}})
  \otimes_{\mathcal{O}_Y}\mathcal{O}_{Y'}
  \longrightarrow
  \mathcal{O}_{Y_{f'}'^{(n)}}
  \tag{16.2.3.1}
  \label{IV.16.2.3.1}
\]
\oldpage[IV-4]{12}
deduced from the homomorphism $\nu_n$ (16.2.1.2) is bijective. Moreover,
the homomorphism of $\mathcal{O}_{Y'}$-modules
\[
  \operatorname{Gr}_1(u):
  \mathscr{G}r_1(f)\otimes_{\mathcal{O}_Y}\mathcal{O}_{Y'}
  \longrightarrow\mathscr{G}r_1(f')
  \tag{16.2.3.2}
  \label{IV.16.2.3.2}
\]
is bijective.
\end{enumerate}
\end{corollary}

\noindent(i) Note first that the morphisms $f':Y'\to X'$ and $u:Y'\to Y$
identify $Y'$ with the product $Y\times_X X'$ (for the structural morphisms
$f:Y\to X$ and $v:X'\to X$) (14.5.12.1). The conclusion of (i) then follows from
(16.2.2, (ii)), the morphism $g$ being an immersion.

\noindent(ii) The commutative diagram
\[
  \xymatrix{
    Y_f^{(n)} \ar[d]^{h_n}
      & Y_{f'}'^{(n)} \ar[d]^-{h'_n} \ar[l]^{w_n} \\
    X \ar[d]^g
      & X' \ar[d]^{g'} \ar[l]^v \\
    Y
      & Y' \ar[l]^u
  }
\]
identifies $Y_{f'}'^{(n)}$ with the product $Y_f^{(n)}\times_X X'$ and $X'$ with the product
$X\times_Y Y'$, so (I, 3.3.9) it identifies (for the morphisms
$g'\circ h'_n$ and $w_n$) $Y_{f'}'^{(n)}$ with the product
$Y_f^{(n)}\times_Y Y'$. Since $Y_f^{(n)}$ (resp. $Y_{f'}'^{(n)}$) is the
prescheme affine over $Y$ (resp. $Y'$) associated to the
$\mathcal{O}_Y$-algebra $\mathcal{O}_{Y_f^{(n)}}$ (resp. to the
$\mathcal{O}_{Y'}$-algebra $\mathcal{O}_{Y_{f'}'^{(n)}}$), the fact that
the canonical homomorphism (16.2.3.1) is bijective follows from (II, 1.5.2).
Finally, the canonical homomorphism (16.2.3.1) is compatible with the
augmentations
$\mathcal{O}_{Y_f^{(n)}}\to\mathcal{O}_Y$ and
$\mathcal{O}_{Y_{f'}'^{(n)}}\to\mathcal{O}_{Y'}$; since
$\mathcal{O}_{Y_f^{(n)}}$ is the direct sum (as an
$\mathcal{O}_Y$-module) of $\mathcal{O}_Y$ and the augmentation ideal
$\mathscr{I}/\mathscr{I}^{n+1}$, we therefore see that the canonical homomorphism
(16.2.3.1), restricted to
$(\mathscr{I}/\mathscr{I}^{n+1})\otimes_{\mathcal{O}_Y}\mathcal{O}_{Y'}$,
is a bijection of the latter onto
$\mathscr{I}'/\mathscr{I}'^{n+1}$. For $n=1$, this shows that
$\operatorname{Gr}_1(u)$ is bijective.

Note that, under the hypotheses of (16.2.3), the homomorphisms
$\operatorname{Gr}_n(u)$ are surjective by virtue of the above, but are
not bijective in general for $n\geq 2$. However:

\begin{corollary}[16.2.4]
\label{IV.16.2.4}
Under the hypotheses of (16.2.3), suppose that $u:Y'\to Y$ is a flat
morphism (resp. that the $\mathscr{G}r_n(f)$ are flat $\mathcal{O}_Y$-modules
for $n\leq m$). Then the homomorphism
\[
  \operatorname{Gr}_n(u):
  \mathscr{G}r_n(f)\otimes_{\mathcal{O}_Y}\mathcal{O}_{Y'}
  \longrightarrow\mathscr{G}r_n(f')
\]
is bijective for every $n$ (resp. for $n\leq m$).
\end{corollary}

If $u$ is flat, so is $v:X'\to X$, which is deduced from it by
base change, and we already know in this case that $\operatorname{Gr}(u)$ is
bijective (16.2.2, (iii)). If the $\mathscr{G}r_n(f)$ are flat for
$n\leq m$, we first see by induction on $n$ that the same holds for the
$\mathscr{I}/\mathscr{I}^{n+1}$ for $n\leq m$, by virtue of the exact sequences
\[
  0\longrightarrow\mathscr{I}^n/\mathscr{I}^{n+1}
  \longrightarrow\mathscr{I}/\mathscr{I}^{n+1}
  \longrightarrow\mathscr{I}/\mathscr{I}^n
  \longrightarrow 0
\]
\oldpage[IV-4]{13}
(0\textsubscript{I}, 6.1.2); moreover, we then have commutative diagrams
\[
  \xymatrix{
    0 \ar[r] &
    (\mathscr{I}^n/\mathscr{I}^{n+1})
      \otimes_{\mathcal{O}_Y}\mathcal{O}_{Y'}
      \ar[r]\ar[d] &
    (\mathscr{I}/\mathscr{I}^{n+1})
      \otimes_{\mathcal{O}_Y}\mathcal{O}_{Y'}
      \ar[r]\ar[d] &
    (\mathscr{I}/\mathscr{I}^n)
      \otimes_{\mathcal{O}_Y}\mathcal{O}_{Y'}
      \ar[r]\ar[d] & 0 \\
    0 \ar[r] &
    \mathscr{I}'^n/\mathscr{I}'^{n+1}
      \ar[r] &
    \mathscr{I}'/\mathscr{I}'^{n+1}
      \ar[r] &
    \mathscr{I}'/\mathscr{I}'^n
      \ar[r] & 0
  }
\]
in which the rows are exact (the first by flatness
(0\textsubscript{I}, 6.1.2)) and the last two vertical arrows are
bijective by virtue of (16.2.3, (ii)); whence the conclusion.

\begin{remarks}[16.2.5]
\label{IV.16.2.5}
\begin{enumerate}
\item[(i)] The argument of (16.2.2, (i)) still applies when in
(16.2.1.1) we are dealing with morphisms of ringed spaces in local rings
($\mathbf{Err}_{\mathrm{II}}$, (1.8.2)).
\item[(ii)] In (16.2.2, (ii)), the conclusion is no longer necessarily valid
when we only suppose that $v$ and $f$ are morphisms of preschemes
($f$ satisfying the condition of (16.1.1)). For example (with the notation of the
proof of (16.2.2, (ii)), it can happen that $\mathfrak{I}=0$ but that
the kernel $\mathfrak{I}'$ of $A'\to B'=B\otimes_A A'$ is not zero and
$B'\neq 0$, in which case we have $Y^{(n)}=Y$ for every $n$, but
$Y'^{(n)}\neq Y'$. An example of this is obtained by taking
$A=\mathbb{Z}$, $B=\mathbb{Q}$,
$A'=\prod_{h=1}^{\infty}(\mathbb{Z}/m^h\mathbb{Z})$ where $m>1$.
\end{enumerate}
\end{remarks}

\begin{env}[16.2.6]
\label{IV.16.2.6}
Consider the particular case of the diagram (16.2.1.1) where $X'=X$, $v$ being
the identity, $X$ is a prescheme, $Y$ a sub-prescheme of $X$, $Y'$ a
sub-prescheme of $Y$, and $f$, $u$ and $f'=f\circ u$ the canonical injections; the
di-homomorphism (16.2.1.3) gives, by tensoring with
$\mathcal{O}_{Y'}$ over $\rho^*(\mathcal{O}_Y)$, a homomorphism of
graded $\mathcal{O}_{Y'}$-algebras
\[
  u^*(\mathscr{G}r_\bullet(f))
  \longrightarrow\mathscr{G}r_\bullet(f')
  \tag{16.2.6.1}
  \label{IV.16.2.6.1}
\]
On the other hand, we identify $\mathcal{O}_Y$ with
$\psi^*(\mathcal{O}_X)/\mathscr{I}_f$ and $\mathcal{O}_{Y'}$ with
$\rho^*(\mathcal{O}_Y)/\mathscr{I}_u$; since $\rho^*$ is an exact functor,
we have
$\rho^*(\mathcal{O}_Y)=
\rho^*(\psi^*(\mathcal{O}_X))/\rho^*(\mathscr{I}_f)=
\psi'^*(\mathcal{O}_X)/\rho^*(\mathscr{I}_f)$, and since
$\mathcal{O}_{Y'}$ is moreover identified with
$\psi'^*(\mathcal{O}_X)/\mathscr{I}_{f'}$, we see that we have
$\mathscr{I}_u=\mathscr{I}_{f'}/\rho^*(\mathscr{I}_f)$. We deduce from this, for
every integer $n$, a canonical homomorphism
$\mathscr{I}_{f'}^n/\mathscr{I}_{f'}^{n+1}\to
\mathscr{I}_u^n/\mathscr{I}_u^{n+1}$, whence a canonical homomorphism of
graded $\mathcal{O}_{Y'}$-algebras
\[
  \mathscr{G}r_\bullet(f')
  \longrightarrow\mathscr{G}r_\bullet(u).
  \tag{16.2.6.2}
  \label{IV.16.2.6.2}
\]
\end{env}

\begin{proposition}[16.2.7]
\label{IV.16.2.7}
Let $X$ be a prescheme, $Y$ a sub-prescheme of $X$, $Y'$ a
sub-prescheme of $Y$, and $j:Y'\to Y$ the canonical injection. We then have an exact
sequence of conormal sheaves ($\mathcal{O}_{Y'}$-modules)
\[
  \xymatrix{
    j^*(\mathscr{N}_{Y/X}) \ar[r] &
    \mathscr{N}_{Y'/X} \ar[r] &
    \mathscr{N}_{Y'/Y} \ar[r] & 0
  }
  \tag{16.2.7.1}
  \label{IV.16.2.7.1}
\]
where the arrows are the degree 1 components of the canonical homomorphisms
(16.2.6.1) and (16.2.6.2).
\end{proposition}

The question being local, we may restrict to the case where
$X=\operatorname{Spec}(A)$, $Y=\operatorname{Spec}(A/\mathfrak{I})$,
$Y'=\operatorname{Spec}(A/\mathfrak{K})$, $\mathfrak{I}$ and $\mathfrak{K}$
being two ideals of $A$ such that $\mathfrak{I}\subset\mathfrak{K}$; everything
then comes down to seeing
\oldpage[IV-4]{14}
that the sequence of canonical homomorphisms
$\mathfrak{I}/\mathfrak{K}\mathfrak{I}\to
\mathfrak{K}/\mathfrak{K}^2\to
(\mathfrak{K}/\mathfrak{I})/(\mathfrak{K}/\mathfrak{I})^2\to 0$
is exact, which is immediate since the image of
$\mathfrak{I}/\mathfrak{K}\mathfrak{I}$ in
$\mathfrak{K}/\mathfrak{K}^2$ is
$(\mathfrak{I}+\mathfrak{K}^2)/\mathfrak{K}^2$ and
$(\mathfrak{K}/\mathfrak{I})/(\mathfrak{K}/\mathfrak{I})^2$ is identified with
$\mathfrak{K}/(\mathfrak{I}+\mathfrak{K}^2)$.

It is easy to give examples where the sequence (16.2.7.1), extended on the left
by a 0, is no longer exact; with the preceding notation, it suffices to
take $A=k[T]$, $\mathfrak{I}=AT^2$, $\mathfrak{K}=AT$, since we then have
$(\mathfrak{I}+\mathfrak{K}^2)/\mathfrak{K}^2=0$ and
$\mathfrak{I}/\mathfrak{K}\mathfrak{I}\neq 0$. See however (16.9.13) and
(19.1.5) for useful cases where the extended sequence still remains exact.

\subsection{Fundamental differential invariants of a morphism of preschemes}
\EGAReaderTarget{ega4.c04.subsec.16-3}{IV.16.3}

\begin{definition}[16.3.1]\label{IV.16.3.1}
\EGAReaderTarget{ega4.c04.def.16-3-1}{IV.16.3.1}
Let $f:X\to S$ be a morphism of preschemes and $\Delta_f:X\to X\times_S X$
the corresponding diagonal morphism, which is an immersion (I, 5.3.9 and
$\mathbf{Err}_{\mathrm{III}}$, 10). We denote by $\mathscr{P}_f^n$ or
$\mathscr{P}_{X/S}^n$, and call the sheaf of principal parts of order $n$
of the $S$-prescheme $X$, the sheaf of $\mathcal{O}_X$-augmented rings that is the
$n$-th normal invariant of $\Delta_f$ (16.1.2). We set
$\mathscr{P}_f^\infty=\mathscr{P}_{X/S}^\infty=
\varprojlim_n\mathscr{P}_{X/S}^n$,
$\mathscr{G}r_n(\mathscr{P}_f)=\mathscr{G}r_n(\mathscr{P}_{X/S})=
\mathscr{G}r_n(\Delta_f)$ (16.1.2); the $\mathcal{O}_X$-module
$\mathscr{G}r_1(\Delta_f)$, the augmentation ideal of $\mathscr{P}_{X/S}^1$, is
also denoted by $\Omega_f^1$ or $\Omega_{X/S}^1$, and called the $\mathcal{O}_X$-module
of first differentials of $f$, or of $X$ relative to $S$, or of the $S$-prescheme $X$.
\end{definition}

It follows from this definition that $\mathscr{P}_{X/S}^0$ is
canonically identified with $\mathcal{O}_X$ (16.1.2).

We have (16.1.2.2) a canonical surjective homomorphism of
graded $\mathcal{O}_X$-algebras
\[
  \mathscr{S}_{\mathcal{O}_X}^{\bullet}(\Omega_{X/S}^1)
  \longrightarrow\mathscr{G}r_\bullet(\mathscr{P}_{X/S}).
  \tag{16.3.1.1}
  \label{IV.16.3.1.1}
\]
It also follows from the definition (16.3.1) that for every open subset $U$ of $X$ we have
$\mathscr{P}_{f|U}^n=\mathscr{P}_f^n|U$,
$\mathscr{P}_{f|U}^\infty=\mathscr{P}_f^\infty|U$,
$\mathscr{G}r_n(\mathscr{P}_{f|U})=\mathscr{G}r_n(\mathscr{P}_f)|U$,
$\Omega_{f|U}^1=\Omega_f^1|U$ (in other words, the notions introduced are
local on $X$).

\begin{env}[16.3.2]
\label{IV.16.3.2}
Denote by $p_1$, $p_2$ the two canonical projections of the product $X\times_S X$;
since $\Delta_f$ is an $X$-section of $X\times_S X$ for both of the
morphisms $p_1$ and $p_2$, each of these morphisms defines, for every $n$, a
homomorphism of sheaves of rings $\mathcal{O}_X\to\mathscr{P}_{X/S}^n$,
right inverse to the augmentation $\mathscr{P}_{X/S}^n\to\mathcal{O}_X$
(16.1.7); we can also say that we thus define on
$\mathscr{P}_{X/S}^n$ two structures of quasi-coherent augmented
$\mathcal{O}_X$-algebra; the corresponding $\mathcal{O}_X$-module structures on
$\mathscr{G}r_n(\mathscr{P}_{X/S})$ are the same. Likewise, by passage to
the limit, we have two $\mathcal{O}_X$-algebra structures on
$\mathscr{P}_{X/S}^\infty$.
\end{env}

\begin{env}[16.3.3]\label{IV.16.3.3}
\EGAReaderTarget{ega4.c04.statement.16-3-3}{IV.16.3.3}
The morphism $s=(p_2,p_1)_S:X\times_S X\to X\times_S X$ is an involutive
automorphism of $X\times_S X$, called the canonical symmetry, such that
\[
  p_1\circ s=p_2,\qquad p_2\circ s=p_1,\qquad s\circ\Delta_f=\Delta_f.
  \tag{16.3.3.1}
  \label{IV.16.3.3.1}
\]
If we set $s=(\rho,\lambda)$, $p_i=(\pi_i,\mu_i)$ ($i=1,2$),
$\Delta_f=(\delta,\nu)$, then $\lambda^\#$ is an isomorphism of
$\rho^*(\pi_1^*(\mathcal{O}_X))$ onto $\pi_2^*(\mathcal{O}_X)$, and
$\delta^*(\lambda^\#)$ leaves invariant $\delta^*(\mathcal{O}_{X\times_S X})$
and the kernel $\mathscr{I}$ of the homomorphism
$\nu^\#:\delta^*(\mathcal{O}_{X\times_S X})\to\mathcal{O}_X$. Consequently:
\end{env}

\begin{proposition}[16.3.4]\label{IV.16.3.4}
\EGAReaderTarget{ega4.c04.prop.16-3-4}{IV.16.3.4}
The homomorphism $\sigma=\delta^*(\lambda^\#)$ deduced from $s$ (and also called
the canonical symmetry) is an involutive automorphism of the projective system
$(\mathscr{P}_{X/S}^n)$ of sheaves
\oldpage[IV-4]{15}
of $\mathcal{O}_X$-augmented rings, and consequently also of their projective limit
$\mathscr{P}_{X/S}^\infty$. This automorphism interchanges the two
$\mathcal{O}_X$-algebra structures on the $\mathscr{P}_{X/S}^n$ and on
$\mathscr{P}_{X/S}^\infty$.
\end{proposition}

\begin{env}[16.3.5]
\label{IV.16.3.5}
In what follows, the two $\mathcal{O}_X$-algebra structures defined on the
$\mathscr{P}_{X/S}^n$ and on $\mathscr{P}_{X/S}^\infty$ will play very
different roles: we henceforth adopt the convention, unless expressly stated
otherwise, that when $\mathscr{P}_{X/S}^n$ or $\mathscr{P}_{X/S}^\infty$ is
considered as an $\mathcal{O}_X$-algebra, it is the algebra structure defined by
$p_1$ that is meant.
\end{env}

For every open subset $U$ of $X$ and every section $t\in\Gamma(U,\mathcal{O}_X)$, we
will denote simply by $t.1$ or even $t$ the image of $t$ under the structural
homomorphism $\Gamma(U,\mathcal{O}_X)\to\Gamma(U,\mathscr{P}_{X/S}^n)$ (resp.
$\Gamma(U,\mathcal{O}_X)\to\Gamma(U,\mathscr{P}_{X/S}^\infty)$)
(that is, the homomorphism corresponding to $p_1$).

\begin{definition}[16.3.6]\label{IV.16.3.6}
\EGAReaderTarget{ega4.c04.def.16-3-6}{IV.16.3.6}
We denote by $d_f^n$, or $d_{X/S}^n$ (resp. $d_f^\infty$, or
$d_{X/S}^\infty$), or simply $d^n$ (resp. $d^\infty$), the homomorphism of
sheaves of rings $\mathcal{O}_X\to\mathscr{P}_f^n=\mathscr{P}_{X/S}^n$
(resp. $\mathcal{O}_X\to\mathscr{P}_f^\infty=\mathscr{P}_{X/S}^\infty$)
deduced from $p_2$. For every open subset $U$ of $X$ and every
$t\in\Gamma(U,\mathcal{O}_X)$, $d^nt$ (resp. $d^\infty t$) is called the principal
part of order $n$ (resp. the principal part of infinite order) of $t$. We set
$dt=d^1t-t$, and we say that $dt$ is the differential of $t$ (an element of
$\Gamma(U,\Omega_{X/S}^1)$, also denoted by $d_{X/S}(t)$).
\end{definition}

It follows immediately from this definition that we have
\[
  d(t_1t_2)=t_1dt_2+t_2dt_1
  \tag{16.3.6.1}
  \label{IV.16.3.6.1}
\]
for all $t_1,t_2$ in $\Gamma(U,\mathcal{O}_X)$; in other words, $d$
is a derivation of the ring $\Gamma(U,\mathcal{O}_X)$ into the
$\Gamma(U,\mathcal{O}_X)$-module $\Gamma(U,\Omega_{X/S}^1)$.

In all the notation introduced in (16.3.1) and (16.3.6), we will
sometimes replace $S$ by $A$ when $S=\operatorname{Spec}(A)$.

\begin{env}[16.3.7]
\label{IV.16.3.7}
Suppose in particular that $S=\operatorname{Spec}(A)$ and
$X=\operatorname{Spec}(B)$ are affine schemes, $B$ thus being an
$A$-algebra. Then $\Delta_f$ corresponds to the canonical surjective homomorphism
$\pi:B\otimes_A B\to B$ such that $\pi(b\otimes b')=bb'$, with kernel
$\mathfrak{I}=\mathfrak{I}_{B/A}$ (0, 20.4.1); $\mathscr{P}_f^n$ is the
structural sheaf of the prescheme $\operatorname{Spec}(P_{B/A}^n)$, where
\[
  P_{B/A}^n=(B\otimes_A B)/\mathfrak{I}^{n+1};
\]
$\mathscr{G}r_\bullet(\mathscr{P}_f)$ is the quasi-coherent
$\mathcal{O}_X$-module corresponding to the graded $B$-module
\[
  \operatorname{gr}_{\mathfrak{I}}^\bullet(B\otimes_A B)
  =\bigoplus_{n\geq 0}(\mathfrak{I}^n/\mathfrak{I}^{n+1});
\]
in particular $\Omega_f^1=\Omega_{X/S}^1$ is the quasi-coherent
$\mathcal{O}_X$-module corresponding to the $B$-module of first differentials of $B$
relative to $A$, $\Omega_{B/A}^1$ (0, 20.4.3). The projection morphisms
$p_1:X\times_S X\to X$, $p_2:X\times_S X\to X$ correspond to the two
ring homomorphisms $j_1:B\to B\otimes_A B$, $j_2:B\to B\otimes_A B$
such that $j_1(b)=b\otimes 1$, $j_2(b)=1\otimes b$, so that (by the
convention of (16.3.5)), $P_{B/A}^n$ is always considered as a
$B$-algebra via the composite homomorphism
$B\xrightarrow{j_1}B\otimes_A B\to P_{B/A}^n$; the ring homomorphism
$B\xrightarrow{j_2}B\otimes_A B\to P_{B/A}^n$ is denoted by $d_{B/A}^n$ and corresponds
to $d_{X/S}^n$ acting on $\Gamma(X,\mathcal{O}_X)$; for every $t\in B$, $dt$
is equal to $d_{B/A}t$, defined in (0, 20.4.6) (cf. $\mathbf{Err}_{\mathrm{IV}}$,
11).

If $\pi_n:B\otimes_A B\to P_{B/A}^n$ is the canonical homomorphism, we therefore have,
by virtue of the preceding definitions,
\[
  \pi_n(b\otimes b')=b.\pi_n(1\otimes b')=b.d_{B/A}^n(b')
  \quad\text{for }b\in B,\ b'\in B.
  \tag{16.3.7.1}
  \label{IV.16.3.7.1}
\]
\end{env}

\oldpage[IV-4]{16}
\begin{proposition}[16.3.8]
\label{IV.16.3.8}
The image of the homomorphism $d_{X/S}^n:\mathcal{O}_X\to\mathscr{P}_{X/S}^n$
generates the $\mathcal{O}_X$-module $\mathscr{P}_{X/S}^n$.
\end{proposition}

We reduce immediately to the case where $X=\operatorname{Spec}(B)$ and
$S=\operatorname{Spec}(A)$ are affine, and the proposition follows from (16.3.7.1),
since $\pi_n$ is surjective. Note that in general $d_{X/S}^n$ is not
surjective (already for $n=1$).

\begin{proposition}[16.3.9]
\label{IV.16.3.9}
Suppose that $f:X\to S$ is a morphism locally of finite type. Then the
$\mathscr{P}_f^n$ and the $\mathscr{G}r_n(\mathscr{P}_f)$ are quasi-coherent
$\mathcal{O}_X$-modules of finite type.
\end{proposition}

This follows from (16.1.6) and from the fact that $\Delta_f$ is locally of finite
presentation (I, 4.3.1).


\subsection{Functorial properties of the differential invariants}
\EGAReaderTarget{ega4.c04.subsec.16-4}{IV.16.4}

\begin{env}[16.4.1]
\label{IV.16.4.1}
Consider a commutative diagram of morphisms of preschemes
\[
  \xymatrix{
    X \ar[d]_-{f} & X' \ar[l]_-u \ar[d]^-{f'} \\
    S & S' \ar[l]^-w
  }
  \tag{16.4.1.1}
  \label{IV.16.4.1.1}
\]
From it we deduce a commutative diagram
\[
  \xymatrix{
    X \ar[d]_-{\Delta_f} & X' \ar[l]_-u \ar[d]^-{\Delta_{f'}} \\
    X\times_S X & X'\times_{S'}X' \ar[l]^-v
  }
\]
where $v$ is the composite homomorphism (I, 5.3.5 and 5.3.15)
\[
  X'\times_{S'}X'
  \xrightarrow{(p'_1,p'_2)_S}X'\times_S X'
  \xrightarrow{u\times_S u}X\times_S X.
  \tag{16.4.1.2}
  \label{IV.16.4.1.2}
\]
From $u$ and $v$ we therefore deduce, as was explained in (16.2.1),
homomorphisms of sheaves of augmented rings
\[
  \nu_n:\rho^*(\mathscr{P}_{X/S}^n)\longrightarrow
  \mathscr{P}_{X'/S'}^n
  \tag{16.4.1.3}
  \label{IV.16.4.1.3}
\]
(where we have set $u=(\rho,\lambda)$); these homomorphisms form a projective
system, and therefore give in the limit a homomorphism of sheaves of augmented
rings
\[
  \nu_\infty:\rho^*(\mathscr{P}_{X/S}^\infty)\longrightarrow
  \mathscr{P}_{X'/S'}^\infty;
  \tag{16.4.1.4}
  \label{IV.16.4.1.4}
\]
on the other hand, by passing to the quotients, the homomorphisms $\nu_n$ give a
di-homomorphism of graded algebras (relative to $\lambda^\#$):
\[
  \operatorname{gr}(u):\rho^*(\mathscr{G}r_\bullet(\mathscr{P}_{X/S}))
  \longrightarrow\mathscr{G}r_\bullet(\mathscr{P}_{X'/S'}).
  \tag{16.4.1.5}
  \label{IV.16.4.1.5}
\]
\end{env}

\begin{env}[16.4.2]
\label{IV.16.4.2}
\oldpage[IV-4]{17}
If we have a commutative diagram
\[
  \xymatrix{
    X \ar[d]_-{f} & X' \ar[l]_-u \ar[d]^-{f'} &
      X'' \ar[l]_-{u'} \ar[d]^-{f''} \\
    S & S' \ar[l]^-w & S'' \ar[l]^-{w'}
  }
\]
we deduce from it a commutative diagram
\[
  \xymatrix{
    X \ar[d]_-{\Delta_f} & X' \ar[l]_-u \ar[d]^-{\Delta_{f'}} &
      X'' \ar[l]_-{u'} \ar[d]^-{\Delta_{f''}} \\
    X\times_S X & X'\times_{S'}X' \ar[l]^-v &
      X''\times_{S''}X'' \ar[l]^-{v'}
  }
\]
where $v'$ is defined from $u'$, $w'$, $f'$, $f''$ as $v$ is from
$u$, $w$, $f$, $f'$. We verify immediately that if $u''=u\circ u'$,
$w''=w\circ w'$, then the composite morphism $v\circ v'$ is equal to the morphism
$v''$ defined from $u''$, $w''$, $f$, $f''$ as $v$ is from $u$,
$w$, $f$, $f'$. If we set $u'=(\rho',\lambda')$,
$u''=(\rho'',\lambda'')$, it then follows from (16.2.1) that the homomorphism
$\nu_n'':\rho''^*(\mathscr{P}_{X/S}^n)\to\mathscr{P}_{X''/S''}^n$ is equal to the
composite
\[
  \rho'^*(\rho^*(\mathscr{P}_{X/S}^n))
  \xrightarrow{\rho'^*(\nu_n^\#)}
  \rho'^*(\mathscr{P}_{X'/S'}^n)
  \xrightarrow{\nu_n'}\mathscr{P}_{X''/S''}^n
\]
and we have analogous transitivity properties for the homomorphisms
(16.4.1.4) and (16.4.1.5), which allows us to say that the
$\mathscr{P}_{X/S}^n$, $\mathscr{P}_{X/S}^\infty$ and
$\mathscr{G}r_\bullet(\mathscr{P}_{X/S})$ depend functorially on $f$.
\end{env}

\begin{env}[16.4.3]
\label{IV.16.4.3}
We verify immediately (for example by reducing to the affine case by means of
(16.3.7)) that, with the notation of (16.4.1), the diagram
\[
  \xymatrix{
    \rho^*(\mathcal{O}_X) \ar[r]^{\lambda^\#} \ar[d] &
      \mathcal{O}_{X'} \ar[d] \\
    \rho^*(\mathscr{P}_{X/S}^n) \ar[r]_{\nu_n} &
      \mathscr{P}_{X'/S'}^n
  }
  \tag{16.4.3.1}
  \label{IV.16.4.3.1}
\]
where the vertical arrows are those defining the algebra structures
chosen in (16.3.5) (that is, those coming from the first projections),
is commutative; the same holds for the diagram
\[
  \xymatrix{
    \rho^*(\mathcal{O}_X) \ar[r]^{\lambda^\#}
      \ar[d]_{\rho^*(d_{X/S}^n)} &
      \mathcal{O}_{X'} \ar[d]^{d_{X'/S'}^n} \\
    \rho^*(\mathscr{P}_{X/S}^n) \ar[r]_{\nu_n} &
      \mathscr{P}_{X'/S'}^n
  }
  \tag{16.4.3.2}
  \label{IV.16.4.3.2}
\]
\oldpage[IV-4]{18}
the vertical arrows here defining the algebra structures coming from the
second projections; moreover, if $\sigma$ and $\sigma'$ are the canonical
symmetries corresponding to $f$ and $f'$ (16.3.4), we have
\[
  \nu_n\circ\rho^*(\sigma)=\sigma'\circ\nu_n
\]
which allows one to pass from one of the preceding diagrams to the other. We
therefore deduce from (16.4.3.1) a canonical homomorphism of augmented
$\mathcal{O}_{X'}$-algebras
\[
  P^n(u):u^*(\mathscr{P}_{X/S}^n)
  =\mathscr{P}_{X/S}^n\otimes_{\mathcal{O}_X}\mathcal{O}_{X'}
  \longrightarrow\mathscr{P}_{X'/S'}^n
  \tag{16.4.3.3}
  \label{IV.16.4.3.3}
\]
and it follows from (16.4.3.2) that the diagram
\[
  \xymatrix{
    \mathcal{O}_{X'} \ar[r]^{\operatorname{id}}
      \ar[d]_{u^*(d_{X/S}^n)} &
      \mathcal{O}_{X'} \ar[d]^{d_{X'/S'}^n} \\
    u^*(\mathscr{P}_{X/S}^n) \ar[r]_{P^n(u)} &
      \mathscr{P}_{X'/S'}^n
  }
  \tag{16.4.3.4}
  \label{IV.16.4.3.4}
\]
is commutative. From this we deduce a homomorphism of graded
$\mathcal{O}_{X'}$-algebras
\[
  \operatorname{Gr}_\bullet(u):u^*(\mathscr{G}r_\bullet(\mathscr{P}_{X/S}))
  \longrightarrow\mathscr{G}r_\bullet(\mathscr{P}_{X'/S'})
  \tag{16.4.3.5}
  \label{IV.16.4.3.5}
\]
and in particular a homomorphism of $\mathcal{O}_{X'}$-modules
\[
  \operatorname{Gr}_1(u):
  \Omega_{X/S}^1\otimes_{\mathcal{O}_X}\mathcal{O}_{X'}
  \longrightarrow\Omega_{X'/S'}^1
  \tag{16.4.3.6}
  \label{IV.16.4.3.6}
\]
giving rise to a commutative diagram
\[
  \xymatrix{
    \mathcal{O}_{X'} \ar[r]^{\operatorname{id}}
      \ar[d]_{d_{X/S}\otimes 1} &
      \mathcal{O}_{X'} \ar[d]^{d_{X'/S'}} \\
    \Omega_{X/S}^1\otimes_{\mathcal{O}_X}\mathcal{O}_{X'} \ar[r] &
      \Omega_{X'/S'}^1
  }
  \tag{16.4.3.7}
  \label{IV.16.4.3.7}
\]
\end{env}

\begin{env}[16.4.4]
\label{IV.16.4.4}
When $S=\operatorname{Spec}(A)$, $S'=\operatorname{Spec}(A')$,
$X=\operatorname{Spec}(B)$, $X'=\operatorname{Spec}(B')$ are affine, so
that we have a commutative diagram of ring homomorphisms
\[
  \xymatrix{
    B \ar[r] & B' \\
    A \ar[u] \ar[r] & A' \ar[u]
  }
\]
the image of $\mathfrak{I}_{B/A}$ in $B'\otimes_{A'}B'$ is contained in
$\mathfrak{I}_{B'/A'}$, and the homomorphism $\nu_n$ corresponds to the ring
homomorphism $P_{B/A}^n\to P_{B'/A'}^n$ deduced from the homomorphism
$B\otimes_A B\to B'\otimes_{A'}B'$ by passing to the quotients. The homomorphism
(16.4.3.6) corresponds to the homomorphism defined in (0, 20.5.4.1), and the
commutative diagram (16.4.3.7) to the diagram (0, 20.5.4.2).
\end{env}

\oldpage[IV-4]{19}
\begin{proposition}[16.4.5]
\label{IV.16.4.5}
Suppose that $X'=X\times_S S'$, $f'$ and $u$ being the canonical projections.
Then the canonical homomorphisms $P^n(u)$ (16.4.3.3) and
$\operatorname{Gr}_1(u)$ (16.4.3.6) are bijective.
\end{proposition}

Indeed, we then have $X'\times_{S'}X'=(X\times_S X)\times_S S'$, and we can
therefore apply (16.2.3, (ii)), replacing $g$ by the first projection
$p_1:X\times_S X\to X$ and $f$ by the diagonal $\Delta_f$.

Note also that under the hypotheses of (16.4.5) the homomorphism
$\operatorname{Gr}_\bullet(u)$ (16.4.3.5) is surjective, but not bijective in
general. However (16.2.4):

\begin{corollary}[16.4.6]
\label{IV.16.4.6}
With the hypotheses of (16.4.5), suppose moreover that $w:S'\to S$ is
flat (resp. that the $\mathscr{G}r_n(\mathscr{P}_{X/S})$ are flat
$\mathcal{O}_X$-modules for $n\leq m$); then the homomorphism
\[
  \operatorname{Gr}_n(u):u^*(\mathscr{G}r_n(\mathscr{P}_{X/S}))
  \longrightarrow\mathscr{G}r_n(\mathscr{P}_{X'/S'})
\]
is bijective for every $n$ (resp. for $n\leq m$).
\end{corollary}

Indeed, if $w$ is flat, so is
$v:X'\times_{S'}X'\to X\times_S X$, hence the conclusion follows from (16.2.4).

\begin{env}[16.4.7]
\label{IV.16.4.7}
Let $S$ be a prescheme, $\mathscr{E}$ a quasi-coherent $\mathcal{O}_S$-module,
and set $X=\mathbf{V}(\mathscr{E})$ (II, 1.7.8), the vector bundle
associated to $\mathscr{E}$, equal to
$\operatorname{Spec}(\mathscr{S}_{\mathcal{O}_S}(\mathscr{E}))$. Let
$f:X\to S$ be the structural morphism. For every open subset $U$ of $S$ and every
section $t\in\Gamma(U,\mathscr{E})$, $t$ is identified with a section of
$\mathscr{S}_{\mathcal{O}_S}(\mathscr{E})$ over $U$; let $t'$ be its
image in
$\Gamma(f^{-1}(U),\mathcal{O}_X)=\Gamma(U,f_*(\mathcal{O}_X))=
\Gamma(U,\mathscr{S}_{\mathcal{O}_S}(\mathscr{E}))$, and set
\[
  \delta(t)=d_{X/S}^n(t')-t'\in
  \Gamma(f^{-1}(U),\mathscr{P}_{X/S}^n);
  \tag{16.4.7.1}
  \label{IV.16.4.7.1}
\]
it is clear that $\delta$ is a di-homomorphism of modules (corresponding to
the ring homomorphism
$\Gamma(U,\mathcal{O}_S)\to\Gamma(f^{-1}(U),\mathcal{O}_X)$) from
$\Gamma(U,\mathscr{E})$ to $\Gamma(f^{-1}(U),\mathscr{P}_{X/S}^n)$, whose
image moreover lies in the augmentation ideal of
$\Gamma(f^{-1}(U),\mathscr{P}_{X/S}^n)$. From this we deduce (by varying $U$)
a canonical homomorphism of $\mathcal{O}_X$-algebras
\[
  f^*(\mathscr{S}_{\mathcal{O}_S}(\mathscr{E}))
  \longrightarrow\mathscr{P}_{X/S}^n
  \tag{16.4.7.2}
  \label{IV.16.4.7.2}
\]
and by the preceding remark, if $\mathscr{K}$ is the ideal that is the kernel of
the augmentation $\mathscr{S}_{\mathcal{O}_S}(\mathscr{E})\to\mathcal{O}_S$,
the image of $\mathscr{K}^{n+1}$ under (16.4.7.2) is zero, so that by
factoring out $\mathscr{K}^{n+1}$ we finally obtain a canonical homomorphism
\[
  \delta_n:f^*(\mathscr{S}_{\mathcal{O}_S}(\mathscr{E})/
  \mathscr{K}^{n+1})\longrightarrow\mathscr{P}_{X/S}^n.
  \tag{16.4.7.3}
  \label{IV.16.4.7.3}
\]
\end{env}

\begin{proposition}[16.4.8]
\label{IV.16.4.8}
Under the conditions of (16.4.7), the homomorphisms $\delta_n$ are bijective and
form a projective system of isomorphisms; from this we deduce an isomorphism of
graded $\mathcal{O}_X$-algebras
\[
  f^*(\mathscr{S}_{\mathcal{O}_S}^{\bullet}(\mathscr{E}))
  \longrightarrow\mathscr{G}r_\bullet(\mathscr{P}_{X/S}).
  \tag{16.4.8.1}
  \label{IV.16.4.8.1}
\]
\end{proposition}

The fact that the homomorphisms (16.4.7.3) form a projective system follows
immediately from their definition. To prove that they are isomorphisms, it
will suffice to
\oldpage[IV-4]{20}
prove that (16.4.8.1) is an isomorphism, the filtrations of the two sides of
(16.4.7.3) being finite (Bourbaki, \emph{Alg. comm.}, chap. III, \S~2,
n\textsuperscript{o} 8, cor. 3 of th. 1). For this, consider the split exact
sequence of $\mathcal{O}_S$-modules
\[
  \xymatrix{
    0 \ar[r] & \mathscr{E} \ar[r]^{u} &
    \mathscr{E}\oplus\mathscr{E} \ar[r]^{v} &
    \mathscr{E} \ar[r] & 0
  }
  \tag{16.4.8.2}
  \label{IV.16.4.8.2}
\]
where, for every pair of sections $s$, $t$ of $\mathscr{E}$ over an open subset
$U$ of $S$, we take $u(s)=(-s,s)$ and $v(s,t)=s+t$. We have
\[
  X\times_S X=\operatorname{Spec}
  (\mathscr{S}_{\mathcal{O}_S}(\mathscr{E})\otimes_{\mathcal{O}_S}
  \mathscr{S}_{\mathcal{O}_S}(\mathscr{E}))
  =\operatorname{Spec}(\mathscr{S}_{\mathcal{O}_S}
  (\mathscr{E}\oplus\mathscr{E}))
\]
(II, 1.4.6 and 1.7.11), and the diagonal morphism $X\to X\times_S X$ corresponds
(II, 1.2.7) to the homomorphism of $\mathcal{O}_S$-algebras
$\mathscr{S}(v):\mathscr{S}_{\mathcal{O}_S}(\mathscr{E}\oplus\mathscr{E})
\to\mathscr{S}_{\mathcal{O}_S}(\mathscr{E})$ (II, 1.7.4), so that if
$\mathscr{I}$ is the kernel of this homomorphism, we have
\[
  \mathscr{P}_{X/S}^n=f^*(\mathscr{S}_{\mathcal{O}_S}
  (\mathscr{E}\oplus\mathscr{E})/\mathscr{I}^{n+1}).
\]

The proposition will be a consequence of the following lemma:

\begin{lemma}[16.4.8.3]
\label{IV.16.4.8.3}
Let $Y$ be a ringed space and
$0\to\mathscr{F}'\xrightarrow{u}\mathscr{F}\xrightarrow{v}\mathscr{F}''\to0$
an exact sequence of $\mathcal{O}_Y$-modules such that every point $y\in Y$
has an open neighbourhood $V$ for which the sequence
$0\to\mathscr{F}'|V\to\mathscr{F}|V\to\mathscr{F}''|V\to0$ is split.
Let $\mathscr{I}$ be the ideal that is the kernel of
$\mathscr{S}(v):\mathscr{S}_{\mathcal{O}_Y}(\mathscr{F})\to
\mathscr{S}_{\mathcal{O}_Y}(\mathscr{F}'')$, and let
$\operatorname{gr}_{\mathscr{I}}^\bullet
(\mathscr{S}_{\mathcal{O}_Y}(\mathscr{F}))$ be the graded $\mathcal{O}_Y$-algebra
associated to the $\mathcal{O}_Y$-algebra
$\mathscr{S}_{\mathcal{O}_Y}(\mathscr{F})$ endowed with the
$\mathscr{I}$-preadic filtration. Then the homomorphism of graded
$\mathcal{O}_Y$-algebras
\[
  \mathscr{S}_{\mathcal{O}_Y}^{\bullet}(\mathscr{F}')
  \otimes_{\mathcal{O}_Y}
  \mathscr{S}_{\mathcal{O}_Y}^{\bullet}(\mathscr{F}'')
  \longrightarrow
  \operatorname{gr}_{\mathscr{I}}^\bullet
  (\mathscr{S}_{\mathcal{O}_Y}(\mathscr{F}))
  \tag{16.4.8.4}
  \label{IV.16.4.8.4}
\]
(where the left-hand side is the graded tensor product of the symmetric
$\mathcal{O}_Y$-algebras endowed with their canonical gradings
(II, 1.7.4 and 2.1.2)), arising from the canonical injection
\[
  \mathscr{F}'\longrightarrow\mathscr{I}
  =\operatorname{gr}_{\mathscr{I}}^1
  (\mathscr{S}_{\mathcal{O}_Y}(\mathscr{F})),
\]
is bijective.
\end{lemma}

Indeed, the injection $\mathscr{F}'\to\mathscr{I}$ canonically gives a
homomorphism of graded $\mathcal{O}_Y$-algebras\linebreak
$\mathscr{S}_{\mathcal{O}_Y}^{\bullet}(\mathscr{F}')\to\allowbreak
\operatorname{gr}_{\mathscr{I}}^\bullet
(\mathscr{S}_{\mathcal{O}_Y}(\mathscr{F}))$, and since the right-hand side is by
definition a graded $\mathscr{S}_{\mathcal{O}_Y}^{\bullet}(\mathscr{F}'')$-algebra,
we deduce the canonical homomorphism (16.4.8.4) by tensoring the preceding one
with $\mathscr{S}_{\mathcal{O}_Y}^{\bullet}(\mathscr{F}'')$. To
prove the lemma, we may, the question being local, restrict to the case where
$\mathscr{F}=\mathscr{F}'\oplus\mathscr{F}''$, $u$ and $v$ being the
canonical homomorphisms. Then the graded algebra
$\mathscr{S}_{\mathcal{O}_Y}^{\bullet}(\mathscr{F})$ is canonically identified
with the graded tensor product
$\mathscr{S}_{\mathcal{O}_Y}^{\bullet}(\mathscr{F}')
\otimes_{\mathcal{O}_Y}
\mathscr{S}_{\mathcal{O}_Y}^{\bullet}(\mathscr{F}'')$ (II, 1.7.4), and it is
immediate that $\mathscr{I}$ is then the ideal
$\mathscr{I}'\otimes_{\mathcal{O}_Y}
\mathscr{S}_{\mathcal{O}_Y}^{\bullet}(\mathscr{F}'')$, where $\mathscr{I}'$ is
the augmentation ideal of
$\mathscr{S}_{\mathcal{O}_Y}^{\bullet}(\mathscr{F}')$, that is, the
(direct) sum of the $\mathscr{S}_{\mathcal{O}_Y}^m(\mathscr{F}')$ for $m\geq1$.
We conclude that $\mathscr{I}^n=\mathscr{I}'^n\otimes_{\mathcal{O}_Y}
\mathscr{S}_{\mathcal{O}_Y}^{\bullet}(\mathscr{F}'')$, where this time
$\mathscr{I}'^n$ is the (direct) sum of the
$\mathscr{S}_{\mathcal{O}_Y}^m(\mathscr{F}')$ for $m\geq n$; consequently
$\mathscr{I}^n/\mathscr{I}^{n+1}=
\mathscr{S}_{\mathcal{O}_Y}^n(\mathscr{F}')\otimes_{\mathcal{O}_Y}
\mathscr{S}_{\mathcal{O}_Y}^{\bullet}(\mathscr{F}'')$, which proves that
(16.4.8.4) is bijective.

\oldpage[IV-4]{21}
This lemma having been proved, it remains to see that the homomorphism
(16.4.8.1) is indeed the image under $f^*$ of the homomorphism (16.4.8.4)
corresponding to the exact sequence (16.4.8.2); one easily checks that this
follows from the definition of $u$ (16.4.8.2) and that of $\delta$ (16.4.7.1),
taking into account the definition of the $\mathcal{O}_X$-algebra structure on
$\mathscr{P}_{X/S}^n$ and that of $d_{X/S}^n$ (16.3.5 and 16.3.6).

In particular:

\begin{corollary}[16.4.9]
\label{IV.16.4.9}
Under the conditions of (16.4.7), we have a canonical isomorphism
\[
  \operatorname{gr}_1(\delta):f^*(\mathscr{E})
  \xrightarrow{\sim}\Omega_{X/S}^1.
  \tag{16.4.9.1}
  \label{IV.16.4.9.1}
\]
\end{corollary}

\begin{corollary}[16.4.10]
\label{IV.16.4.10}
If $S=\operatorname{Spec}(A)$ and $\mathscr{E}=\mathcal{O}_S^m$, so that
\[
  X=\operatorname{Spec}(A[T_1,\ldots,T_m]),
\]
then $\mathscr{P}_{X/S}^n$ is canonically identified with the
$\mathcal{O}_X$-algebra corresponding to the quotient $A[T_1,\ldots,T_m]$-algebra
\[
  A[T_1,\ldots,T_m,U_1,\ldots,U_m]/\mathfrak{K}^{n+1},
\]
where the $U_i$ ($1\leq i\leq m$) are $m$ new indeterminates and
$\mathfrak{K}$ is the ideal generated by $U_1,\ldots,U_m$.
\end{corollary}

In particular, we thus recover the structure of $\Omega_{X/S}^1$ in this case
(0, 20.5.13).

Note moreover that $d_{X/S}^n$ then sends a polynomial
$F(T_1,\ldots,T_m)$ to the class mod.\ $\mathfrak{K}^{n+1}$ of
$F(T_1+U_1,\ldots,T_m+U_m)$, as follows from the definition (16.4.7.1).

\begin{proposition}[16.4.11]
\label{IV.16.4.11}
Let $f:X\to S$ be a morphism, $g:S\to X$ an $S$-section of $X$, and
$S_g^{(n)}$ the $n$-th infinitesimal neighbourhood of $S$ for the immersion $g$
(16.1.2). Then there exists one and only one isomorphism of
$\mathcal{O}_S$-algebras
\[
  \varpi_n:g^*(\mathscr{P}_{X/S}^n)\longrightarrow
  \mathcal{O}_{S_g^{(n)}}
  \tag{16.4.11.1}
  \label{IV.16.4.11.1}
\]
(for the $\mathcal{O}_S$-algebra structure on
$\mathcal{O}_{S_g^{(n)}}$ defined by $f$ (16.1.7)) such that the
diagram
\[
  \xymatrix{
    \mathcal{O}_S=g^*(\mathcal{O}_X)
      \ar[rr]^-{\lambda_n}\ar[dr]_-{g^*(d_{X/S}^n)}
      &&\mathcal{O}_{S_g^{(n)}}\\
    &g^*(\mathscr{P}_{X/S}^n)\ar[ur]_-{\varpi_n}
  }
  \tag{16.4.11.2}
  \label{IV.16.4.11.2}
\]
(where $\lambda_n$ is the structural homomorphism) is commutative.
\end{proposition}

By virtue of (I, 5.3.7), where we replace $X$, $Y$, $S$ by $S$, $X$, $S$
respectively and $f$ by $g$, the diagrams
\[
  \xymatrix@=3pc{
    S\ar[r]^-g\ar[d]_-g&X\ar[d]^-{\Delta_f}
      &S\ar[r]^-g\ar[d]_-g&X\ar[d]^-{\Delta_f}\\
    X\ar[r]_{(g\circ f,1_X)_S}&X\times_S X
      &X\ar[r]_{(1_X,g\circ f)_S}&X\times_S X
  }
  \tag{16.4.11.3}
  \label{IV.16.4.11.3}
\]
\oldpage[IV-4]{22}
identify $S$ with the product of the $(X\times_S X)$-preschemes $X$ and $X$ for the
morphisms $\Delta_f$ and $(g\circ f,1_X)_S$ (resp.\ $(1_X,g\circ f)_S$).
On the other hand, the diagrams
\[
  \xymatrix@=3pc{
    X\ar[r]^{(g\circ f,1_X)_S}\ar[d]_-f&X\times_S X\ar[d]^-{p_1}
      &X\ar[r]^{(1_X,g\circ f)_S}\ar[d]_-f&X\times_S X\ar[d]^-{p_2}\\
    S\ar[r]_{g}&X&S\ar[r]_{g}&X
  }
  \tag{16.4.11.4}
  \label{IV.16.4.11.4}
\]
identify $X$ with the product of the $X$-preschemes $S$ and $X\times_S X$ for the
morphisms $g$ and $p_1$ (resp.\ $p_2$) (a particular case of the associativity
formula (I, 3.3.9.1)). We can say that $\Delta_f$, regarded as an
$X$-section of $X\times_S X$ (relative to $p_1$ or $p_2$), plays the role
of a \emph{universal section} for the $S$-sections of $X$: indeed, each of these
sections $g$ is deduced from it by the \emph{base change}
$(g\circ f,1_X)_S:X\to X\times_S X$. The definition of the homomorphism
$\varpi_n$ and the fact that it is bijective therefore follow from these
remarks and from (16.2.3, (ii)) applied to the first diagram (16.4.11.4). The
commutativity of the diagram (16.4.11.2) follows in the same way from
(16.2.3, (ii)), applied this time to the second diagram (16.4.11.4). To make
$\varpi_n$ explicit, we may restrict to the case where $g$ is a \emph{closed}
immersion: indeed, for every $s\in S$ there is an open neighbourhood $W$ of $s$
in $S$ such that $g(W)$ is closed in an open subset $U$ of $X$, and it is clear
that $g|W$ is a $W$-section of the morphism $U\cap f^{-1}(W)\to W$, the
restriction of $f$, and $g(W)$ is \emph{a fortiori} closed in
$U\cap f^{-1}(W)$. We may therefore suppose that $S$ is a closed
sub-prescheme of $X$ defined by a quasi-coherent ideal $\mathscr{K}$.
Then the preceding definitions show that if $W$ is an open subset of $S$ and
$t$ a section of $\mathcal{O}_X$ over $f^{-1}(W)$,
$\varpi_n(d^n t|W)$ is equal to the canonical image of $t$ in
$\Gamma(W,(\mathcal{O}_X/\mathscr{K}^{n+1})|W)$. The uniqueness of $\varpi_n$
therefore follows from the fact that the image of $\mathcal{O}_X$ under
$d_{X/S}^n$ generates the $\mathcal{O}_X$-module $\mathscr{P}_{X/S}^n$ (16.3.8).

\begin{corollary}[16.4.12]
\label{IV.16.4.12}
Let $k$ be a field, $X$ a $k$-prescheme, and $x$ a point of $X$ rational over
$k$. Then $(\mathscr{P}_{X/k}^n)_x\otimes_{\mathcal{O}_x}k(x)$ is
canonically isomorphic (as an augmented $k(x)$-algebra) to
$\mathcal{O}_x/\mathfrak{m}_x^{n+1}$.
\end{corollary}

It suffices to consider the unique $k$-section $g$ of $X$ such that
$g(\operatorname{Spec}(k))=\{x\}$.

\begin{corollary}[16.4.13]
\label{IV.16.4.13}
Let $f:X\to S$ be a morphism, $s$ a point of $S$, and
$X_s=X\times_S\operatorname{Spec}(k(s))$ the fibre of $f$ at $s$. If
$x\in X_s$ is rational over $k(s)$, then
$(\mathscr{P}_{X/S}^n)_x\otimes_{\mathcal{O}_s}k(s)$ is canonically
isomorphic to $\mathcal{O}_{X_s,x}/\mathfrak{m}'_x{}^{\,n+1}$, where
$\mathfrak{m}'_x$ is the maximal ideal of $\mathcal{O}_{X_s,x}$; more
precisely, this isomorphism sends $(d^n t)_x\otimes1$ (where $t$ is a
section of $\mathcal{O}_X$ over an open neighbourhood of $x$ in
$X$) to the class of $t_x\otimes1$ modulo $\mathfrak{m}'_x{}^{\,n+1}$.
\end{corollary}

This follows from (16.4.5) and (16.4.12).

The preceding corollaries justify the terminology ``sheaf of principal parts
of order $n$''.

\begin{proposition}[16.4.14]
\label{IV.16.4.14}
Let $\rho:A\to B$ be a ring homomorphism and $S$ a multiplicatively stable
subset of $B$. Then the canonical homomorphisms
\[
  S^{-1}P_{B/A}^n\longrightarrow P_{S^{-1}B/A}^n
  \tag{16.4.14.1}
  \label{IV.16.4.14.1}
\]
\oldpage[IV-4]{23}
deduced from the canonical homomorphisms
$P_{B/A}^n\to P_{S^{-1}B/A}^n$ (16.4.4) form a projective system and are
bijective.
\end{proposition}

It suffices to remark that
\[
  S^{-1}((B\otimes_A B)/\mathfrak{I}^{n+1})
  =S^{-1}(B\otimes_A B)/(S^{-1}\mathfrak{I})^{n+1}
\]
by flatness, and that
$S^{-1}(B\otimes_A B)=(S^{-1}B)\otimes_A(S^{-1}B)$ (I, 1.3.4).

\begin{corollary}[16.4.15]
\label{IV.16.4.15}
With the notation of (16.4.14), let $R$ be a multiplicative subset of
$A$ such that $\rho(R)\subset S$. Then we have canonical isomorphisms
\[
  S^{-1}P_{B/A}^n\xrightarrow{\sim}P_{S^{-1}B/R^{-1}A}^n
  \tag{16.4.15.1}
  \label{IV.16.4.15.1}
\]
forming a projective system.
\end{corollary}

It evidently suffices to define canonical isomorphisms
\[
  P_{S^{-1}B/A}^n\xrightarrow{\sim}P_{S^{-1}B/R^{-1}A}^n,
  \tag{16.4.15.2}
  \label{IV.16.4.15.2}
\]
that is, we are reduced to the case where $\rho(R)$ consists of
\emph{invertible} elements of $B$. But then the isomorphism (16.4.15.2) is simply
deduced from the canonical isomorphism
$B\otimes_A B\to B\otimes_{R^{-1}A}B$ by passing to the quotients
(0$_{\rm I}$, 1.5.3).

\begin{corollary}[16.4.16]
\label{IV.16.4.16}
Let $f:X\to S$ be a morphism of preschemes, $x$ a point of $X$, and
$s=f(x)$. Then we have canonical isomorphisms
\[
  (\mathscr{P}_{X/S}^n)_x\xrightarrow{\sim}
  P_{\mathcal{O}_x/\mathcal{O}_s}^n
  \tag{16.4.16.1}
  \label{IV.16.4.16.1}
\]
forming a projective system.
\end{corollary}

From this we deduce isomorphisms for the associated graded objects, and in
particular a canonical isomorphism
\[
  (\Omega_{X/S}^1)_x\xrightarrow{\sim}
  \Omega_{\mathcal{O}_x/\mathcal{O}_s}^1.
  \tag{16.4.16.2}
  \label{IV.16.4.16.2}
\]

\begin{corollary}[16.4.17]
\label{IV.16.4.17}
Let $k$ be a field and $K$ the field of rational functions
$k(T_1,\ldots,T_r)$. Then, for every integer $n$, the homomorphism from
$K[U_1,\ldots,U_r]$ ($U_i$ indeterminates) to $P_{K/k}^n$ which sends
each $U_i$ to $d^nT_i-T_i\cdot1$ is surjective and defines an isomorphism
of the quotient $K[U_1,\ldots,U_r]/\mathfrak{m}^{n+1}$ (where $\mathfrak{m}$ is
the ideal generated by the $U_i$) onto $P_{K/k}^n$.
\end{corollary}

This follows from (16.4.8), (16.4.10) and (16.4.14), where we take $A=k$,
$B=k[T_1,\ldots,T_r]$ and $S=B-\{0\}$.

We thus recover the fact that the $dT_i$ form a basis of the $K$-vector
space $\Omega_{K/k}^1$ (0, 20.5.10).

\begin{proposition}[16.4.18]
\label{IV.16.4.18}
Let $f:X\to Y$, $g:Y\to Z$ be two morphisms of preschemes, and consider
the canonical homomorphisms of augmented $\mathcal{O}_X$-algebras
(16.4.3.3)
\[
  g_{X/Y/Z}:\mathscr{P}_{X/Z}^n\longrightarrow\mathscr{P}_{X/Y}^n
  \tag{16.4.18.1}
  \label{IV.16.4.18.1}
\]
\[
  f_{X/Y/Z}:f^*(\mathscr{P}_{Y/Z}^n)
  \longrightarrow\mathscr{P}_{X/Z}^n.
  \tag{16.4.18.2}
  \label{IV.16.4.18.2}
\]
Then $g_{X/Y/Z}$ is surjective, and its kernel is the ideal generated by the
image under $f_{X/Y/Z}$ of the augmentation ideal of
$f^*(\mathscr{P}_{Y/Z}^n)$.
\end{proposition}

\oldpage[IV-4]{24}
Note first that $g_{X/Y/Z}$ corresponds to the case where in (16.4.3.3) we take
$X'=X$, $S'=Y$, $S=Z$ and $u=1_X$, and $f_{X/Y/Z}$ to the case where we replace
$X'$, $X$, $S$, $S'$ by $X$, $Y$, $Z$, $Z$ respectively and $u$, $f$ by
$f$, $g$ respectively.

We have a commutative diagram (I, 5.3.5)
\[
  \xymatrix{
    X\ar[r]^-{\Delta_f}\ar[dr]_-f
      &X\times_Y X\ar[d]^-p\ar[r]^-j
      &X\times_Z X\ar[d]^-{f\times_Z f}\\
    &Y\ar[r]_-{\Delta_g}&Y\times_Z Y
  }
  \tag{16.4.18.3}
  \label{IV.16.4.18.3}
\]
where $j=(1_X,1_X)_Z$ is an immersion, $j\circ\Delta_f=\Delta_{g\circ f}$, and
$p$ is the structural morphism. Since we may restrict to the case where $X$, $Y$,
$Z$ are affine, we may suppose the immersions $\Delta_f$, $\Delta_g$ and
$j$ to be closed, so that $\mathcal{O}_X$ and
$\mathcal{O}_{X\times_Y X}$ are identified respectively with
$\mathcal{O}_{X\times_Z X}/\mathscr{J}$ and
$\mathcal{O}_{X\times_Z X}/\mathscr{L}$, where
$\mathscr{J}\supset\mathscr{L}$ are the two quasi-coherent ideals
corresponding respectively to the immersions $\Delta_{g\circ f}$ and $j$. The
$\mathcal{O}_X$-algebra $\mathscr{P}_{X/Z}^n$ is therefore identified with
$\mathcal{O}_{X\times_Z X}/\mathscr{J}^{n+1}$, and
$\mathscr{P}_{X/Y}^n$ is identified with
$\mathcal{O}_{X\times_Y X}/(\mathscr{J}/\mathscr{L})^{n+1}$, that is, with
$\mathcal{O}_{X\times_Z X}/(\mathscr{J}^{n+1}+\mathscr{L})$, and consequently
with the quotient of $\mathscr{P}_{X/Z}^n$ by
$(\mathscr{J}^{n+1}+\mathscr{L})/\mathscr{J}^{n+1}$. But we know
(\emph{loc. cit.}) that $p$ and $j$ make $X\times_Y X$ the product of the
$(Y\times_Z Y)$-preschemes $Y$ and $X\times_Z X$; hence, if $\mathcal{O}_Y$ is
identified with $\mathcal{O}_{Y\times_Z Y}/\mathscr{K}$, where $\mathscr{K}$ is
the ideal corresponding to $\Delta_g$, then $\mathscr{L}$ is equal to
$(f\times_Z f)^*(\mathscr{K})\cdot\mathcal{O}_{X\times_Z X}$ (I, 4.4.5).
Since $(\mathscr{J}^{n+1}+\mathscr{L})/\mathscr{J}^{n+1}$ is the ideal of
$\mathscr{P}_{X/Z}^n$ generated by the image of $\mathscr{L}$, the
proposition follows.

\begin{corollary}[16.4.19]
\label{IV.16.4.19}
With the notation of (16.4.18), we have an exact sequence of quasi-coherent
$\mathcal{O}_X$-modules
\[
  \xymatrix{
    f^*(\Omega_{Y/Z}^1)\ar[r]^-{f_{X/Y/Z}}&
    \Omega_{X/Z}^1\ar[r]^-{g_{X/Y/Z}}&
    \Omega_{X/Y}^1\ar[r]&0.
  }
  \tag{16.4.19.1}
  \label{IV.16.4.19.1}
\]
\end{corollary}

When $X$, $Y$, $Z$ are affine, we thus recover the exact sequence
(0, 20.5.7.1).

\begin{proposition}[16.4.20]
\label{IV.16.4.20}
Let $f:Y\to Z$ be a morphism, $j:X\to Y$ a closed immersion, and
$\mathscr{K}$ the quasi-coherent ideal of $\mathcal{O}_Y$ corresponding to $j$.
Then we have
$\mathscr{P}_{X/Y}^n=\mathcal{O}_X=\mathcal{O}_Y/\mathscr{K}$,
the canonical homomorphism
$j_{X/Y/Z}:j^*(\mathscr{P}_{Y/Z}^n)\to\mathscr{P}_{X/Z}^n$ is surjective, and
its kernel is the ideal of $j^*(\mathscr{P}_{Y/Z}^n)$ generated by
$j^*(\mathcal{O}_Y\cdot d_{Y/Z}^n(\mathscr{K}))$ (it should be noted that
$d_{Y/Z}^n(\mathscr{K})$ is a sub-sheaf of commutative groups of
$\mathscr{P}_{Y/Z}^n$, but not an $\mathcal{O}_Y$-module in general).
\end{proposition}

We know (I, 5.3.8) that the diagonal
$\Delta_j:X\to X\times_Y X$ is an isomorphism, whence the first assertion.
If $\varpi_1$ and $\varpi_2$ are the two algebra homomorphisms
$\mathcal{O}_Y\to\mathscr{P}_{Y/Z}^n$ corresponding respectively to the two
canonical projections $p_1$, $p_2$ of $Y\times_Z Y$ onto $Y$, recall that
by definition (16.3.5 and 16.3.6) $\varpi_1$ is the structural homomorphism of
the $\mathcal{O}_Y$-algebra $\mathscr{P}_{Y/Z}^n$ and
$\varpi_2=d_{Y/Z}^n$. The $\mathcal{O}_X$-algebra
$j^*(\mathscr{P}_{Y/Z}^n)$ is therefore identified with
$\mathscr{P}_{Y/Z}^n/\varpi_1(\mathscr{K})\mathscr{P}_{Y/Z}^n$, and its quotient
by the ideal generated by $j^*(d_{Y/Z}^n(\mathscr{K}))$ with
\[
  \mathscr{P}_{Y/Z}^n/
  (\varpi_1(\mathscr{K})+\varpi_2(\mathscr{K}))\mathscr{P}_{Y/Z}^n.
\]
Note now that we have a commutative diagram
\oldpage[IV-4]{25}
\[
  \xymatrix{
    Y\ar[d]_-{\Delta_f}&X\ar[l]_-j\ar[d]^-{\Delta_{f\circ j}}\\
    Y\times_Z Y&X\times_Z X\ar[l]^-{j\times_Z j}
  }
\]
identifying $X$ with the product of the $(Y\times_Z Y)$-preschemes $Y$ and
$X\times_Z X$ (I, 5.3.7). Since $j\times_Z j$ is an immersion, we therefore
deduce from this remark and from (16.2.2) that, if $\Delta_{Y/Z}^n$ and
$\Delta_{X/Z}^n$ denote the infinitesimal neighbourhoods of order $n$ of $Y$
and $X$ for the canonical immersions $\Delta_f$ and $\Delta_{f\circ j}$
respectively, we have a diagram
\[
  \xymatrix{
    \Delta_{Y/Z}^n\ar[d]&\Delta_{X/Z}^n\ar[l]\ar[d]\\
    Y\times_Z Y&X\times_Z X\ar[l]^-{j\times_Z j}
  }
\]
making $\Delta_{X/Z}^n$ the product of the $(Y\times_Z Y)$-preschemes
$\Delta_{Y/Z}^n$ and $X\times_Z X$. We can also say that
$\mathscr{P}_{X/Z}^n$ is identified with the sheaf of rings
$\mathscr{P}_{Y/Z}^n\otimes_{\mathcal{O}_{Y\times_Z Y}}
\mathcal{O}_{X\times_Z X}$. But one sees immediately (for example by reducing
to the affine case) that
\[
  \mathcal{O}_{X\times_Z X}=
  \mathcal{O}_{Y\times_Z Y}/
  (p_1^*(\mathscr{K})+p_2^*(\mathscr{K}))
  \mathcal{O}_{Y\times_Z Y}.
\]
Hence $\mathscr{P}_{X/Z}^n$ is identified with the quotient of
$\mathscr{P}_{Y/Z}^n$ by the ideal generated by the image in
$\mathscr{P}_{Y/Z}^n$ of $p_1^*(\mathscr{K})+p_2^*(\mathscr{K})$. But by
definition this ideal is also generated by
$\varpi_1(\mathscr{K})+\varpi_2(\mathscr{K})$. Q.E.D.

\begin{corollary}[16.4.21]
\label{IV.16.4.21}
Let $f:Y\to Z$ be a morphism and $j:X\to Y$ an immersion. We have an exact
sequence of quasi-coherent $\mathcal{O}_X$-modules
\[
  \xymatrix{
    \mathscr{N}_{X/Y}\ar[r]&j^*(\Omega_{Y/Z}^1)\ar[r]&
    \Omega_{X/Z}^1\ar[r]&0.
  }
  \tag{16.4.21.1}
  \label{IV.16.4.21.1}
\]
\end{corollary}

When $X$, $Y$ and $Z$ are affine, we thus recover the exact sequence
(0, 20.5.12.1).

\begin{corollary}[16.4.22]
\label{IV.16.4.22}
If $f:X\to S$ is a morphism locally of finite presentation, then
$\mathscr{P}_{X/S}^n$ and $\Omega_{X/S}^1$ are quasi-coherent
$\mathcal{O}_X$-modules of finite presentation.
\end{corollary}

We are immediately reduced to the case where $S=\operatorname{Spec}(A)$ is affine
and $X=\operatorname{Spec}(B)$, where
$B=A[T_1,\ldots,T_r]/\mathfrak{K}$, $\mathfrak{K}$ being an ideal of finite
type of $C=A[T_1,\ldots,T_r]$. We then apply (16.4.20) with $Z=S$,
$Y=\operatorname{Spec}(C)$ and $\mathscr{K}=\widetilde{\mathfrak{K}}$. Then
$j^*(\mathscr{P}_{Y/Z}^n)$ is a free $\mathcal{O}_X$-module of finite rank
(16.4.10), and the hypothesis on $\mathfrak{K}$ implies that
$j^*(\mathcal{O}_Y\cdot d_{Y/Z}^n(\mathscr{K}))$ generates a quasi-coherent
$\mathcal{O}_X$-module of finite type; whence the conclusion.

\begin{proposition}[16.4.23]
\label{IV.16.4.23}
Let $X$, $Y$ be two $S$-preschemes, $Z=X\times_S Y$ their product, and
$p:X\times_S Y\to X$ and $q:X\times_S Y\to Y$ the canonical projections.
Then the canonical homomorphism
\[
  p_{Z/X/S}\oplus q_{Z/Y/S}:
  p^*(\Omega_{X/S}^1)\oplus q^*(\Omega_{Y/S}^1)
  \longrightarrow\Omega_{(X\times_S Y)/S}^1
  \tag{16.4.23.1}
  \label{IV.16.4.23.1}
\]
is bijective.
\end{proposition}

\oldpage[IV-4]{26}
The commutative diagram
\[
  \xymatrix{
    Y\ar[d]_-g&X\times_S Y\ar[l]_-q\ar[d]_-h&
      X\times_S Y\ar[l]_-{\operatorname{id}}\ar[d]^-p\\
    S&S\ar[l]^-{\operatorname{id}}&X\ar[l]^-f
  }
\]
gives a factorization of the canonical \emph{isomorphism} $P^n(p)$ (16.4.5)
\[
  p^*(\mathscr{P}_{X/S}^n)\longrightarrow
  \mathscr{P}_{Z/S}^n\longrightarrow\mathscr{P}_{Z/Y}^n
\]
and similarly, interchanging the roles of $X$ and $Y$, we have a factorization
of the \emph{isomorphism} $P^n(q)$
\[
  q^*(\mathscr{P}_{Y/S}^n)\longrightarrow
  \mathscr{P}_{Z/S}^n\longrightarrow\mathscr{P}_{Z/X}^n.
\]
This proves that the canonical homomorphism (16.4.18.1)
\[
  p_{Z/X/S}:p^*(\mathscr{P}_{X/S}^n)\longrightarrow\mathscr{P}_{Z/S}^n
  \quad\text{(resp.\
  $q_{Z/Y/S}:q^*(\mathscr{P}_{Y/S}^n)\to\mathscr{P}_{Z/S}^n$)}
\]
is \emph{injective}, and that the kernel of the canonical surjective homomorphism
(16.4.18.2)
\[
  \mathscr{P}_{Z/S}^n\longrightarrow\mathscr{P}_{Z/Y}^n
  \quad\text{(resp.\
  $\mathscr{P}_{Z/S}^n\to\mathscr{P}_{Z/X}^n$)}
\]
is complementary to the image of $p_{Z/X/S}$ (resp.\ $q_{Z/Y/S}$). But
on the other hand, by virtue of (16.4.18), this kernel is generated by the image
under $q_{Z/Y/S}$ (resp.\ $p_{Z/X/S}$) of the augmentation ideal of
$q^*(\mathscr{P}_{Y/S}^n)$ (resp.\ $p^*(\mathscr{P}_{X/S}^n)$). We obtain
the proposition by considering the case $n=1$.

We immediately generalize (16.4.23) to the case of a product of an arbitrary
finite number of $S$-preschemes.

\begin{remarks}[16.4.24]
\label{IV.16.4.24}
\begin{enumerate}
\item[(i)] We will see (17.2.3) that when the morphism $f:X\to Y$ in
(16.4.18) is \emph{smooth}, the homomorphism $f_{X/Y/Z}$ in (16.4.19.1) is
locally \emph{left invertible}, and in particular injective. Similarly,
when the morphism $f\circ j:X\to Z$ of (16.4.20) is \emph{smooth},
the left-hand homomorphism in (16.4.21.2) is locally
\emph{left invertible} and \emph{a fortiori} injective (17.2.5). In Chapter V,
we will also give a variant, in the case of modules on preschemes, of the
``modules of imperfection'' studied in (0, 20.6), and of the exact sequences in
which they appear.

\item[(ii)] Let $X$ be a topological space, $\mathscr{A}$ a sheaf
of rings on $X$, and $\mathscr{B}$ an $\mathscr{A}$-algebra on $X$. Then it
is clear that
\[
  U\longmapsto
  P_{\Gamma(U,\mathscr{B})/\Gamma(U,\mathscr{A})}^n
  \qquad (U\text{ open in }X)
\]
is a presheaf of augmented $\Gamma(U,\mathscr{B})$-algebras, hence the
associated sheaf $\mathscr{P}_{\mathscr{B}/\mathscr{A}}^n$ is an augmented
$\mathscr{B}$-algebra. In the particular case where $X$ is a
prescheme and $f=(\psi,\theta):X\to S$ a morphism of preschemes, it follows
easily from (16.4.16) and from the exactness of the functor $\varinjlim$ that
$\mathscr{P}_{X/S}^n$ is canonically isomorphic to
$\mathscr{P}_{\mathcal{O}_X/\psi^*(\mathcal{O}_S)}^n$. It follows that the
formalism developed in the present section could be regarded as
a
\oldpage[IV-4]{27}
particular case of a differential formalism for ringed spaces endowed with
a sheaf of algebras over the structural sheaf. However, we did not wish
to start from this point of view, which is less intuitive and less convenient
for applications. It seems moreover that, for the various kinds of
``varieties'', the ``global'' construction of the $\mathscr{P}^n$ analogous to
the one we use here is likewise better suited to applications.
\end{enumerate}
\end{remarks}


\subsection{Relative tangent sheaves and tangent bundles; derivations.}
\EGAReaderTarget{ega4.c04.subsec.16-5}{IV.16.5}

\begin{env}[16.5.1]\label{IV.16.5.1}
\EGAReaderTarget{ega4.c04.statement.16-5-1}{IV.16.5.1}
Let $f=(\psi,\theta):X\to S$ be a morphism of ringed spaces. For every
$\mathcal{O}_X$-module $\mathscr{F}$, we call \emph{$S$-derivation}
(or \emph{$(X/S)$-derivation}, or \emph{$f$-derivation}) \emph{of
$\mathcal{O}_X$ into $\mathscr{F}$} any homomorphism of \emph{sheaves of
additive groups} $D:\mathcal{O}_X\to\mathscr{F}$ satisfying the following
conditions:
\begin{enumerate}
\item[(a)] for every open subset $V$ of $X$ and every pair of sections
$(t_1,t_2)$ of $\mathcal{O}_X$ over $V$, we have
\[
  D(t_1t_2)=t_1D(t_2)+D(t_1)t_2;
  \tag{16.5.1.1}
  \label{IV.16.5.1.1}
\]
\item[(b)] for every open subset $V$ of $X$, every section $t$ of $\mathcal{O}_X$
over $V$ and every section $s$ of $\mathcal{O}_S$ over an open subset
$U$ of $S$ such that $V\subset f^{-1}(U)$, we have
\[
  D((s|V)t)=(s|V)D(t).
  \tag{16.5.1.2}
  \label{IV.16.5.1.2}
\]
\end{enumerate}
It is clear that this amounts to saying that for every $x\in X$,
the homomorphism of additive groups
$D_x:\mathcal{O}_x\to\mathscr{F}_x$ is an
$\mathcal{O}_{f(x)}$-derivation.

Another interpretation consists in considering the $\mathcal{O}_X$-algebra
$\mathscr{D}_{\mathcal{O}_X}(\mathscr{F})$ equal to
$\mathcal{O}_X\oplus\mathscr{F}$, the algebra structure being defined by the
condition that, for every open subset $V$ of $X$, the product of two sections of
$\mathcal{O}_X$ (resp. of a section of $\mathcal{O}_X$ and a section of
$\mathscr{F}$) over $V$ is defined by the ring structure of
$\Gamma(V,\mathcal{O}_X)$ (resp. the
$\Gamma(V,\mathcal{O}_X)$-module structure on $\Gamma(V,\mathscr{F})$), and the
product of two sections of $\mathscr{F}$ over $V$ is taken to be 0; then
$\mathscr{F}$ is an ideal of
$\mathscr{D}_{\mathcal{O}_X}(\mathscr{F})$, the kernel of the canonical augmentation
$\mathscr{D}_{\mathcal{O}_X}(\mathscr{F})\to\mathcal{O}_X$, and to say that $D$
is an $S$-derivation of $\mathcal{O}_X$ into $\mathscr{F}$ means that
$1_{\mathcal{O}_X}+D$ is an $\mathcal{O}_S$-homomorphism of algebras from
$\mathcal{O}_X$ to $\mathscr{D}_{\mathcal{O}_X}(\mathscr{F})$ which, composed
with the augmentation, gives $1_{\mathcal{O}_X}$.

The $S$-derivations of $\mathcal{O}_X$ into $\mathscr{F}$ evidently form
a $\Gamma(X,\mathcal{O}_X)$-module
$\operatorname{D\acute{e}r}_S(\mathcal{O}_X,\mathscr{F})$.

When $\mathscr{F}=\mathcal{O}_X$, an $S$-derivation of $\mathcal{O}_X$ into
itself is simply called an \emph{$S$-derivation of $\mathcal{O}_X$}.
\end{env}

\begin{proposition}[16.5.2]
\label{IV.16.5.2}
Let $A$ be a ring, $B$ an $A$-algebra, and $L$ a $B$-module; set
$S=\operatorname{Spec}(A)$, $X=\operatorname{Spec}(B)$,
$\mathscr{F}=\widetilde{L}$. Then the map $D\mapsto\Gamma(D)$, which sends
each $S$-derivation $D$ of $\mathcal{O}_X$ into $\mathscr{F}$ to the map
$\Gamma(D):t\mapsto D(t)$ from $B$ to $L$, is an
isomorphism of $B$-modules from
$\operatorname{D\acute{e}r}_S(\mathcal{O}_X,\mathscr{F})$ onto
$\operatorname{D\acute{e}r}_A(B,L)$ (cf. (0, 20.1.2)).
\end{proposition}

This follows immediately from the interpretation of $S$-derivations given above
in terms
\oldpage[IV-4]{28}
of algebra homomorphisms, from the analogous interpretation given in
(0, 20.1.6), and from the canonical correspondence between homomorphisms of
$\mathcal{O}_X$-algebras and homomorphisms of $B$-algebras (I, 1.3.13 and
1.3.8).

\begin{proposition}[16.5.3]
\label{IV.16.5.3}
Let $f=(\psi,\theta):X\to S$ be a morphism of preschemes.
\begin{enumerate}
\item[(i)] The differential
$d_{X/S}:\mathcal{O}_X\to\Omega_{X/S}^1$ (16.3.6) is an $S$-derivation.
\item[(ii)] For every $\mathcal{O}_X$-module $\mathscr{F}$, the map
$u\mapsto u\circ d_{X/S}$ is an isomorphism of
$\Gamma(X,\mathcal{O}_X)$-modules
\[
  \operatorname{Hom}_{\mathcal{O}_X}(\Omega_{X/S}^1,\mathscr{F})
  \xrightarrow{\sim}
  \operatorname{D\acute{e}r}_S(\mathcal{O}_X,\mathscr{F}).
  \tag{16.5.3.1}
  \label{IV.16.5.3.1}
\]
\end{enumerate}
\end{proposition}

Assertion (i) has already been noted (16.3.6). On the other hand, it is immediate
(by virtue of (0, 20.4.8)) that $u\mapsto u\circ d_{X/S}$ is injective, by
considering the restrictions of both sides to a stalk $\mathcal{O}_x$ and
using (16.4.16.2). To see that the homomorphism (16.5.3.1) is surjective,
consider an $S$-derivation $D:\mathcal{O}_X\to\mathscr{F}$; for every
affine open subset $V=\operatorname{Spec}(B)$ of $X$ such that $f(V)$ is contained
in an affine open subset $U=\operatorname{Spec}(A)$ of $S$,
$D_V:B\to\Gamma(V,\mathscr{F})$ is an $A$-derivation, hence there exists a
\emph{unique} $B$-homomorphism
$u_V:\Omega_{B/A}^1\to\Gamma(V,\mathscr{F})$ such that
$D_V=u_V\circ d_{B/A}$ (0, 20.4.8); moreover, the uniqueness of $u_V$ shows
immediately that for every affine open subset $W\subset V$ we have $u_W=u_V|W$, hence
the $u_V$ define a homomorphism of $\mathcal{O}_X$-modules
$u:\mathcal{O}_X\to\mathscr{F}$ that answers the question.

\begin{env}[16.5.4]\label{IV.16.5.4}
\EGAReaderTarget{ega4.c04.statement.16-5-4}{IV.16.5.4}
With the notation of (16.5.1), for every open subset $U$ of $X$,
$\operatorname{D\acute{e}r}_S(\mathcal{O}_U,\mathscr{F}|U)$ is a
$\Gamma(U,\mathcal{O}_X)$-module, and it is clear that the map
$U\mapsto\operatorname{D\acute{e}r}_S(\mathcal{O}_U,\mathscr{F}|U)$ is a
presheaf; in fact, it is even a \emph{sheaf} (hence an
$\mathcal{O}_X$-module), by virtue of the pointwise characterization of
$S$-derivations seen in (16.5.1). This $\mathcal{O}_X$-module is denoted by
$\mathscr{D}\!\mathrm{\acute{e}r}_S(\mathcal{O}_X,\mathscr{F})$ and is called the
\emph{sheaf of $S$-derivations of $\mathcal{O}_X$ into $\mathscr{F}$}, and
what we have just seen can also be expressed by the following corollary:
\end{env}

\begin{corollary}[16.5.5]
\label{IV.16.5.5}
For every $\mathcal{O}_X$-module $\mathscr{F}$, the homomorphism of
$\mathcal{O}_X$-modules deduced from $u\mapsto u\circ d_{X/S}$
\[
  \mathscr{H}\!om_{\mathcal{O}_X}(\Omega_{X/S}^1,\mathscr{F})
  \longrightarrow
  \mathscr{D}\!\mathrm{\acute{e}r}_S(\mathcal{O}_X,\mathscr{F})
  \tag{16.5.5.1}
  \label{IV.16.5.5.1}
\]
is bijective.
\end{corollary}

\begin{corollary}[16.5.6]
\label{IV.16.5.6}
\begin{enumerate}
\item[(i)] If the morphism $f:X\to S$ is locally of finite presentation and
$\mathscr{F}$ is a quasi-coherent $\mathcal{O}_X$-module, then
$\mathscr{D}\!\mathrm{\acute{e}r}_S(\mathcal{O}_X,\mathscr{F})$ is a
quasi-coherent $\mathcal{O}_X$-module.
\item[(ii)] If moreover $S$ is locally Noetherian and $\mathscr{F}$ is
coherent, then $\mathscr{D}\!\mathrm{\acute{e}r}_S(\mathcal{O}_X,\mathscr{F})$ is a
coherent $\mathcal{O}_X$-module.
\end{enumerate}
\end{corollary}

Assertion (i) follows from the isomorphism (16.5.5.1), from (16.4.22) and
(I, 1.3.12); assertion (ii) follows from (0$_{\rm I}$, 5.3.5).

\begin{env}[16.5.7]\label{IV.16.5.7}
\EGAReaderTarget{ega4.c04.statement.16-5-7}{IV.16.5.7}
We set
\[
  \mathfrak{G}_{X/S}=
  \mathscr{H}\!om_{\mathcal{O}_X}(\Omega_{X/S}^1,\mathcal{O}_X)
  =\mathscr{D}\!\mathrm{\acute{e}r}_S(\mathcal{O}_X,\mathcal{O}_X)
  \tag{16.5.7.1}
  \label{IV.16.5.7.1}
\]
and we say that it is the \emph{sheaf of $S$-derivations of
$\mathcal{O}_X$}, or also the \emph{tangent sheaf of $X$ relative to
$S$}: it is thus the \emph{dual} of the $\mathcal{O}_X$-module
$\Omega_{X/S}^1$. If $f$ is locally of finite
\oldpage[IV-4]{29}
presentation, $\mathfrak{G}_{X/S}$ is a quasi-coherent $\mathcal{O}_X$-module; if
moreover $S$ is locally Noetherian, $\mathfrak{G}_{X/S}$ is coherent
(16.5.6).
\end{env}

\begin{env}[16.5.8]
\label{IV.16.5.8}
Suppose more particularly that $\Omega_{X/S}^1$ is a \emph{locally free}
$\mathcal{O}_X$-module (of finite rank) (which will be the
case when $f$ is \emph{smooth} (17.2.3)); then $\mathfrak{G}_{X/S}$ is a locally
free $\mathcal{O}_X$-module and has the same rank as $\Omega_{X/S}^1$
at every point. More precisely, suppose that $\Omega_{X/S}^1$ has
rank $n$ at a point $x$; there are then $n$ sections $s_i$ ($1\leq i\leq n$)
of $\mathcal{O}_X$ over an affine neighbourhood $U$ of $x$ such that the
canonical images of the $ds_i$ in
$\Omega_{X/S}^1\otimes_{\mathcal{O}_X}k(x)$ form a basis of this
$k(x)$-vector space; by Nakayama's lemma, the germs
$(ds_i)_x=d(s_i)_x$ of the $ds_i$ at the point $x$ form a basis of the
$\mathcal{O}_x$-module $(\Omega_{X/S}^1)_x$; hence, by shrinking $U$, we
may suppose that the $ds_i$ form a \emph{basis} of the
$\Gamma(U,\mathcal{O}_X)$-module $\Gamma(U,\Omega_{X/S}^1)$. Then the
$\Gamma(U,\mathcal{O}_X)$-module $\Gamma(U,\mathfrak{G}_{X/S})$ is the dual of
the preceding one; we denote by $(D_i)_{1\leq i\leq n}$ or
$\left(\frac{\partial}{\partial s_i}\right)_{1\leq i\leq n}$ the dual basis of
$(ds_i)_{1\leq i\leq n}$, so that, by (16.5.3), we have
\[
  D_is_j=\langle D_i,ds_j\rangle
  =\left\langle\frac{\partial}{\partial s_i},ds_j\right\rangle
  =\delta_{ij}\qquad\text{(Kronecker delta).}
  \tag{16.5.8.1}
  \label{IV.16.5.8.1}
\]
Every $\Gamma(S,\mathcal{O}_S)$-derivation of the
$\Gamma(S,\mathcal{O}_S)$-algebra $\Gamma(U,\mathcal{O}_X)$ can therefore be
written in one and only one way as
\[
  D=\sum_{i=1}^n a_iD_i=\sum_{i=1}^n a_i\frac{\partial}{\partial s_i}
\]
where the $a_i$ ($1\leq i\leq n$) are sections of $\mathcal{O}_X$ over
$U$. For every section $g\in\Gamma(U,\mathcal{O}_X)$, if we set
$dg=\sum_{i=1}^n c_i ds_i$, we have
$c_i=\langle D_i,dg\rangle=D_ig$ by virtue of (16.5.8.1); in other words
\[
  dg=\sum_{i=1}^n(D_ig)ds_i
  =\sum_{i=1}^n\frac{\partial g}{\partial s_i}ds_i.
  \tag{16.5.8.2}
  \label{IV.16.5.8.2}
\]
\end{env}

\begin{env}[16.5.9]
\label{IV.16.5.9}
Let $D_1$, $D_2$ be two $S$-derivations of $\mathcal{O}_X$. For every open
subset $U$ of $X$, if $D_1^U$, $D_2^U$ are the corresponding derivations of
the ring $\Gamma(U,\mathcal{O}_X)$, the \emph{bracket}
\[
  [D_1^U,D_2^U]=D_1^U\circ D_2^U-D_2^U\circ D_1^U
\]
is also a derivation of this ring; hence the
$\psi^*(\mathcal{O}_S)$-endomorphism of $\mathcal{O}_X$
\[
  [D_1,D_2]=D_1\circ D_2-D_2\circ D_1
  \tag{16.5.9.1}
  \label{IV.16.5.9.1}
\]
is again an $S$-derivation; since one verifies immediately that this bracket
satisfies the Jacobi identity, we see that we have thus defined on
$\operatorname{D\acute{e}r}_S(\mathcal{O}_X,\mathcal{O}_X)$ a structure of
Lie $\Gamma(S,\mathcal{O}_S)$-algebra. Since the definition of this
structure commutes with restriction to an open subset of $X$, we see that
$\mathfrak{G}_{X/S}$ is canonically endowed with a structure of
Lie $\psi^*(\mathcal{O}_S)$-algebra. Note that the map
$(D_1,D_2)\mapsto[D_1,D_2]$ is not
$\Gamma(X,\mathcal{O}_X)$-bilinear.
\end{env}

\begin{env}[16.5.10]
\label{IV.16.5.10}
For every base change $g:S'\to S$, if we set $X'=X\times_S S'$, we
have seen (16.4.5) that we have a canonical isomorphism
\[
  \Omega_{X/S}^1\otimes_S S'\xrightarrow{\sim}\Omega_{X'/S'}^1
  \tag{16.5.10.1}
  \label{IV.16.5.10.1}
\]
\oldpage[IV-4]{30}
from which we deduce, by virtue of (16.5.10.1), a canonical homomorphism
(Bourbaki, \emph{Alg.}, chap. II, 3\textsuperscript{e} éd., \S~5,
n\textsuperscript{o} 3)
\[
  \mathfrak{G}_{X/S}\otimes_{\mathcal{O}_S}\mathcal{O}_{S'}
  \longrightarrow\mathfrak{G}_{X'/S'}
  \tag{16.5.10.2}
  \label{IV.16.5.10.2}
\]
which in general is neither injective nor surjective. However:
\end{env}

\begin{proposition}[16.5.11]
\label{IV.16.5.11}
\begin{enumerate}
\item[(i)] If $g:S'\to S$ is a flat morphism and $f$ is locally of
finite type (resp.\ locally of finite presentation), the homomorphism
(16.5.10.2) is injective (resp.\ bijective).
\item[(ii)] If $\Omega_{X/S}^1$ is a locally free $\mathcal{O}_X$-module
of finite type, the homomorphism (16.5.10.2) is bijective.
\end{enumerate}
\end{proposition}

Indeed, assertion (ii) follows from Bourbaki, \emph{Alg.}, chap. II,
3\textsuperscript{e} éd., \S~5, n\textsuperscript{o} 3, prop. 7. Assertion
(i) follows likewise from Bourbaki, \emph{Alg. comm.}, chap. I, \S~2,
n\textsuperscript{o} 10, prop. 11, and from the fact that if $f$ is locally of
finite type (resp.\ locally of finite presentation), $\Omega_{X/S}^1$ is an
$\mathcal{O}_X$-module of finite type (resp.\ of finite presentation) ((16.3.9)
and (16.4.22)).

\begin{env}[16.5.12]\label{IV.16.5.12}
\EGAReaderTarget{ega4.c04.statement.16-5-12}{IV.16.5.12}
Since $\Omega_{X/S}^1$ is a quasi-coherent $\mathcal{O}_X$-module, we can
consider the vector bundle over $X$ defined by $\Omega_{X/S}^1$
(II, 1.7.8)
\[
  T_{X/S}=\mathbf{V}(\Omega_{X/S}^1)
  \tag{16.5.12.1}
  \label{IV.16.5.12.1}
\]
which is called the \emph{tangent bundle of $X$ relative to $S$}. We thus have a
canonical bijection (II, 1.7.9)
\[
  \Gamma(T_{X/S}/S)\xrightarrow{\sim}
  \operatorname{Hom}_{\mathcal{O}_X}(\Omega_{X/S}^1,\mathcal{O}_X)
  =\Gamma(X,\mathfrak{G}_{X/S})
\]
by definition of $\mathfrak{G}_{X/S}$, and in this isomorphism we can
replace $X$ by an arbitrary open subset $U$ of $X$; we can therefore say that the
\emph{tangent sheaf} of $X$ relative to $S$ is isomorphic to the
\emph{sheaf of germs of $S$-sections} of the tangent bundle of $X$
relative to $S$. If $f:X\to Y$ is an $S$-morphism, we have seen (16.4.19)
that we have a canonical homomorphism
$f_{X/Y/S}:f^*(\Omega_{Y/S}^1)\to\Omega_{X/S}^1$, which gives, taking into
account the fact that
\[
  \mathbf{V}(f^*(\Omega_{Y/S}^1))
  =\mathbf{V}(\Omega_{Y/S}^1)\times_Y X\qquad\text{(II, 1.7.11),}
\]
an $X$-morphism
$T_{X/S}(f):T_{X/S}\to T_{Y/S}\times_Y X$. If $g:Y\to Z$ is a second
$S$-morphism, we have
$T_{X/S}(g\circ f)=(T_{Y/S}(g)\times1_X)\circ T_{X/S}(f)$ (0, 20.5.4.1).

It follows from (16.5.10.1) and (II, 1.7.11) that for every base change
$g:S'\to S$ we have a canonical isomorphism
\[
  T_{X'/S'}\xrightarrow{\sim}T_{X/S}\times_S S'
  =T_{X/S}\times_X X'.
  \tag{16.5.12.2}
  \label{IV.16.5.12.2}
\]
\end{env}

\begin{env}[16.5.13]\label{IV.16.5.13}
\EGAReaderTarget{ega4.c04.statement.16-5-13}{IV.16.5.13}
For every point $x\in X$, we call \emph{tangent space to $X$ at the point $x$}
(relative to $S$) the set of points of the fibre
$T_{X/S}\times_X\operatorname{Spec}(k(x))$ that are \emph{rational over $k(x)$},
hence the set
\[
  T_{X/S}(x)=
  \operatorname{Hom}_{k(x)}
  (\Omega_{X/S}^1\otimes_{\mathcal{O}_x}k(x),k(x))
  \tag{16.5.13.1}
  \label{IV.16.5.13.1}
\]
which is the \emph{dual} of the $k(x)$-vector space
$\Omega_{\mathcal{O}_x/\mathcal{O}_s}^1/
\mathfrak{m}_x\cdot\Omega_{\mathcal{O}_x/\mathcal{O}_s}^1$. When
$\Omega_{X/S}^1$ is an $\mathcal{O}_X$-module of \emph{finite type},
$T_{X/S}(x)$ is therefore a vector space of finite rank over $k(x)$, and for
every
\oldpage[IV-4]{31}
base change $g:S'\to S$ and every point
$x'\in X'=X\times_S S'$ above $x$, we have a canonical isomorphism
\[
  T_{X'/S'}(x')\xrightarrow{\sim}
  T_{X/S}(x)\otimes_{k(x)}k(x').
  \tag{16.5.13.2}
  \label{IV.16.5.13.2}
\]

If $x$ is \emph{rational over $k(s)$}, where $s=f(x)$ (so that
$k(s)\to k(x)$ is an isomorphism), it follows from (16.4.13) that we have a
canonical isomorphism
\[
  T_{X/S}(x)=T_{X_s/k(s)}(x)
  =\operatorname{Hom}_{k(s)}
  (\mathfrak{m}'_x/\mathfrak{m}'_x{}^2,k(s))
  \tag{16.5.13.3}
  \label{IV.16.5.13.3}
\]
where $\mathfrak{m}'_x$ is the maximal ideal of
$\mathcal{O}_{X_s,x}=\mathcal{O}_{X,x}/
\mathfrak{m}_s\mathcal{O}_{X,x}$. In the case where $S$ is the spectrum of a
field $k$, we thus recover Zariski's definition of the tangent space at a
point $x\in X$ \emph{rational over $k$}, as the dual of
$\mathfrak{m}_x/\mathfrak{m}_x^2$.

Let $Y$ be a second $S$-prescheme and let $g:Y\to X$ be an $S$-morphism; we
then have a canonical homomorphism of $\mathcal{O}_Y$-modules (16.4.19)
\[
  g_{Y/X/S}:g^*(\Omega_{X/S}^1)\longrightarrow\Omega_{Y/S}^1.
  \tag{16.5.13.4}
  \label{IV.16.5.13.4}
\]
Note now that if $y\in Y$ and $x=g(y)$, we have
\[
  g^*(\Omega_{X/S}^1)\otimes_{\mathcal{O}_Y}k(y)
  =(\Omega_{X/S}^1\otimes_{\mathcal{O}_X}k(x))
  \otimes_{k(x)}k(y)
\]
and consequently, if $\Omega_{X/S}^1$ is an $\mathcal{O}_X$-module
\emph{of finite type}, we can identify
\[
  \operatorname{Hom}_{k(y)}
  (g^*(\Omega_{X/S}^1)\otimes_{\mathcal{O}_Y}k(y),k(y))
\]
with $T_{X/S}(x)\otimes_{k(x)}k(y)$. We therefore deduce from the homomorphism
(16.5.13.4) a homomorphism of \emph{$k(y)$-vector spaces}
\[
  T_y(g):T_{Y/S}(y)\longrightarrow
  T_{X/S}(x)\otimes_{k(x)}k(y)
  \tag{16.5.13.5}
  \label{IV.16.5.13.5}
\]
called the \emph{tangent linear map to $g$ at the point $y$}. When $y$ is
\emph{rational over $k(s)$}, we can identify $k(s)$, $k(y)$ and $k(x)$, and
$T_y(g)$ is then a homomorphism of $k(s)$-vector spaces
$T_{Y/S}(y)\to T_{X/S}(x)$; note moreover that in this case
$g^*(\Omega_{X/S}^1)\otimes_{\mathcal{O}_Y}k(y)$ is identified with
$\Omega_{X/S}^1\otimes_{\mathcal{O}_X}k(x)$, and the preceding homomorphism is
therefore defined without any finiteness condition on $\Omega_{X/S}^1$ and is
none other than the homomorphism $T_{Y/S}(g)$ (16.5.12) restricted to the fibre
of $T_{Y/S}$ at the point $y$.
\end{env}

\begin{env}[16.5.14]
\label{IV.16.5.14}
The interpretation of the derivations of an $A$-algebra $B$ into a $B$-module
$L$, given in (0, 20.1.1), translates into the language of preschemes as
follows.

Consider two morphisms of preschemes $f:X\to S$, $g:Y\to S$, and a
closed sub-prescheme $Y_0$ of $Y$ defined by an ideal $\mathscr{J}$ of
$\mathcal{O}_Y$ of square zero (so that $Y$ and $Y_0$ have the
\emph{same underlying topological space}). Suppose given an $S$-morphism
$u_0:Y_0\to X$, so that we have a commutative diagram
\[
  \xymatrix{
    X\ar[d]_-f & Y_0\ar[l]_-{u_0}\ar[d]^j\\
    S & Y\ar@{-->}[ul]_-u\ar[l]^-g
  }
  \tag{16.5.14.1}
  \label{IV.16.5.14.1}
\]
\oldpage[IV-4]{32}
and let us investigate whether there exist \emph{$S$-morphisms}
$u:Y\to X$ such that $u_0=u\circ j$ (in other words, whether it is possible to
complete the preceding diagram by the dotted arrow $u$ so as to
keep it \emph{commutative}).

To this end, consider an affine open subset $U=\operatorname{Spec}(C)$ of $Y$;
its inverse image $j^{-1}(U)$ is the affine open subset
$U_0=\operatorname{Spec}(C/\mathfrak{L})$, where
$\mathfrak{L}=\Gamma(U,\mathscr{J})$ is an ideal \emph{of square zero} in $C$;
we will suppose $U$ sufficiently small for $u_0(U_0)$ to be contained in an
affine open subset $V=\operatorname{Spec}(B)$ of $X$, and
$g(U)=f(u_0(U_0))$ to be contained in an affine open subset
$W=\operatorname{Spec}(A)$ of $S$, so that $B$ and $C$ are
$A$-algebras and $u_0|U_0$ corresponds to an $A$-homomorphism $\psi$ from
$B$ to $C/\mathfrak{L}$; let $P(U_0)$ be the set of the restrictions $u|U$
of the homomorphisms sought, which correspond canonically to the
$A$-homomorphisms of algebras $\varphi:B\to C$ such that the \emph{composite}
\[
  B\xrightarrow{\varphi}C\longrightarrow C/\mathfrak{L}
\]
\emph{is equal to $\psi$}. We therefore know (0, 20.1.1) that the set of these
homomorphisms is either empty or of the form
$\varphi_1+\operatorname{D\acute{e}r}_A(B,\mathfrak{L})$; when $P(U_0)$ is
non-empty, the additive group $\operatorname{D\acute{e}r}_A(B,\mathfrak{L})$ acts by
addition on $P(U_0)$, which is therefore an \emph{affine space} for the
additive group $\operatorname{D\acute{e}r}_A(B,\mathfrak{L})$ (or also a
\emph{principal homogeneous space} (or \emph{torsor}) \emph{under}
$\operatorname{D\acute{e}r}_A(B,\mathfrak{L})$).

Now observe that, since $\mathfrak{L}$ is endowed with a
$B$-module structure by means of $\psi$, we have an \emph{isomorphism}
$v\mapsto v\circ d_{B/A}$ of
$\operatorname{Hom}_B(\Omega_{B/A}^1,\mathfrak{L})$ onto
$\operatorname{D\acute{e}r}_A(B,\mathfrak{L})$ (0, 20.4.8). Moreover, since
$\mathfrak{L}$ is of square zero, hence a $(C/\mathfrak{L})$-module, every
$B$-homomorphism $v:\Omega_{B/A}^1\to\mathfrak{L}$ can be regarded
as a $(C/\mathfrak{L})$-homomorphism
$\Omega_{B/A}^1\otimes_B(C/\mathfrak{L})\to\mathfrak{L}$. Since
$\mathscr{J}$ is of square zero, it can be regarded as a quasi-coherent
$\mathcal{O}_{Y_0}$-module; let us introduce the
$\mathcal{O}_{Y_0}$-module
\[
  \mathscr{G}=\mathscr{H}om(u_0^*(\Omega_{X/S}^1),\mathscr{J});
  \tag{16.5.14.2}
  \label{IV.16.5.14.2}
\]
it then follows from the fact that
$\Omega_{B/A}^1=\Gamma(V,\Omega_{X/S}^1)$ (16.3.7) that we can write
$\operatorname{D\acute{e}r}_A(B,\mathfrak{L})=\Gamma(U_0,\mathscr{G})$.

Since $P(U_0)$ is defined as a set of $S$-morphisms $U\to X$, it is
clear that $U_0\mapsto P(U_0)$ is a \emph{sheaf of sets}
$\mathscr{P}$ on $Y_0$. Let us use this fact to prove that the map
\[
  h:\Gamma(U_0,\mathscr{G})\times P(U_0)\longrightarrow P(U_0)
\]
defining the torsor structure on $P(U_0)$ is also independent of the
choice of $V$ and $W$, and moreover that, if $U'\subset U$ is a second affine
open subset of $Y$ and $U'_0$ its inverse image in $Y_0$, the diagram
\[
  \xymatrix{
    \Gamma(U_0,\mathscr{G})\times P(U_0)\ar[d]\ar[r]^h
      & P(U_0)\ar[d]\\
    \Gamma(U'_0,\mathscr{G})\times P(U'_0)\ar[r]_{h'}
      & P(U'_0)
  }
  \tag{16.5.14.3}
  \label{IV.16.5.14.3}
\]
is commutative (the vertical arrows being the restriction operators).
By virtue of the preceding remark, we are reduced to proving the commutativity
of the preceding diagram when $h$ is defined as above from the
affine open subsets $V$, $W$ and $h'$ from affine open
\oldpage[IV-4]{33}
subsets $V'\subset V$ and $W'\subset W$. But by virtue of the preceding
description of $h$, this follows from the commutativity of the diagram
(0, 20.5.4.2).

The maps
$\Gamma(U_0,\mathscr{G})\times P(U_0)\to P(U_0)$ therefore define a
homomorphism of \emph{sheaves of sets}
$m:\mathscr{G}\times\mathscr{P}\to\mathscr{P}$ such that, for every open subset
$U_0$ for which $\Gamma(U_0,\mathscr{P})\neq\varnothing$,
\[
  m_{U_0}:\Gamma(U_0,\mathscr{G})\times\Gamma(U_0,\mathscr{P})
  \longrightarrow\Gamma(U_0,\mathscr{P})
\]
is an external composition law defining on $\Gamma(U_0,\mathscr{P})$ a
structure of torsor under the group $\Gamma(U_0,\mathscr{G})$.
\end{env}

\begin{env}[16.5.15]\label{IV.16.5.15}
\EGAReaderTarget{ega4.c04.statement.16-5-15}{IV.16.5.15}
In general, when we are given on a topological space $Z$ a sheaf of
sets $\mathscr{P}$, a sheaf of groups $\mathscr{G}$ (not necessarily
commutative) and a homomorphism of sheaves of sets
$m:\mathscr{G}\times\mathscr{P}\to\mathscr{P}$ such that, for every open subset
$U\subset Z$ such that $\Gamma(U,\mathscr{P})\neq\varnothing$,
\[
  m_U:\Gamma(U,\mathscr{G})\times\Gamma(U,\mathscr{P})
  \longrightarrow\Gamma(U,\mathscr{P})
\]
makes $\Gamma(U,\mathscr{P})$ a \emph{torsor} under the group
$\Gamma(U,\mathscr{G})$, we say that $\mathscr{P}$ is a
\emph{pseudo-torsor} (or \emph{formally principal homogeneous sheaf})
under the sheaf of groups $\mathscr{G}$. We say that $\mathscr{P}$ is a
\emph{torsor} (or \emph{principal homogeneous sheaf}) under $\mathscr{G}$ if
moreover $\Gamma(U,\mathscr{P})\neq\varnothing$ for every open subset
$U\neq\varnothing$ of a suitable basis of the topology of $Z$.

For the general theory of torsors we refer to [42]; here we will
merely recall the canonical correspondence between the isomorphism classes
of these torsors (for a given $\mathscr{G}$) and the elements of
the cohomology set $H^1(Z,\mathscr{G})$. Indeed, consider a
torsor $\mathscr{P}$ under $\mathscr{G}$ and an open covering
$(U_\lambda)$ of $Z$ such that
$\Gamma(U_\lambda,\mathscr{P})\neq\varnothing$ for every $\lambda$;
let $p_\lambda$ denote an element of
$\Gamma(U_\lambda,\mathscr{P})$. For every pair of indices $\lambda$, $\mu$
such that $U_\lambda\cap U_\mu\neq\varnothing$, there then exists one and
only one element $\gamma_{\lambda\mu}$ of
$\Gamma(U_\lambda\cap U_\mu,\mathscr{G})$ such that
\[
  \gamma_{\lambda\mu}\cdot
  (p_\mu|U_\lambda\cap U_\mu)=p_\lambda|U_\lambda\cap U_\mu;
\]
moreover, if $\lambda$, $\mu$, $\nu$ are three indices such that
$U_\lambda\cap U_\mu\cap U_\nu\neq\varnothing$, the restrictions
$\gamma'_{\lambda\mu}$, $\gamma'_{\mu\nu}$, $\gamma'_{\lambda\nu}$ of
$\gamma_{\lambda\mu}$, $\gamma_{\mu\nu}$, $\gamma_{\lambda\nu}$ to
$U_\lambda\cap U_\mu\cap U_\nu$ satisfy the condition
$\gamma'_{\lambda\nu}=\gamma'_{\lambda\mu}\gamma'_{\mu\nu}$; in other words,
$(\lambda,\mu)\mapsto\gamma_{\lambda\mu}$ is a \emph{1-cocycle} of the
covering $(U_\lambda)$ with values in $\mathscr{G}$. If, for every
$\lambda$, $p'_\lambda$ is a second element of
$\Gamma(U_\lambda,\mathscr{P})$, there exists one and only one
$\beta_\lambda\in\Gamma(U_\lambda,\mathscr{G})$ such that
$p'_\lambda=\beta_\lambda\cdot p_\lambda$, and the 1-cocycle
$(\gamma'_{\lambda\mu})$ corresponding to the family $(p'_\lambda)$ is given
by
\[
  \gamma'_{\lambda\mu}=\beta_\lambda\gamma_{\lambda\mu}\beta_\mu^{-1},
\]
in other words, is \emph{cohomologous} to $(\gamma_{\lambda\mu})$.
Conversely, the datum of a 1-cocycle $(\gamma_{\lambda\mu})$ defines, for
every pair $(\lambda,\mu)$, an automorphism $\theta_{\lambda\mu}$ of the
sheaf \emph{of sets}
$\mathscr{G}|U_\lambda\cap U_\mu$, namely the right translation by
$\gamma_{\lambda\mu}$, and the fact that this is a cocycle shows that we
can \emph{glue} the sheaves of sets $\mathscr{G}|U_\lambda$ by
means of the automorphisms $\theta_{\lambda\mu}$ (0, 3.3.1); we thus obtain
in an evident way a torsor under $\mathscr{G}$, say $\mathscr{P}$, and if
we take for $p_\lambda$ the unit section over $U_\lambda$, the
corresponding 1-cocycle is none other than the given 1-cocycle
$(\gamma_{\lambda\mu})$; moreover, if we replace
$(\gamma_{\lambda\mu})$ by a 1-cocycle
$\gamma'_{\lambda\mu}=\beta_\lambda\gamma_{\lambda\mu}\beta_\mu^{-1}$
cohomologous to it, one verifies immediately that the torsor obtained is
isomorphic to $\mathscr{P}$.

In particular, if $(\gamma_{\lambda\mu})$ is a \emph{1-coboundary}, in other
words of the form
$\gamma_{\lambda\mu}=\beta_\lambda\beta_\mu^{-1}$, the torsor
$\mathscr{P}$ obtained is \emph{isomorphic to $\mathscr{G}$} (regarded as a torsor
under itself for the left translations); in this case we say that
$\mathscr{P}$ is \emph{trivial}, and the converse is evident.

More particularly, it follows from (III, 1.3.1) that we have:
\end{env}

\begin{proposition}[16.5.16]
\label{IV.16.5.16}
Let $Z$ be an affine scheme and $\mathscr{G}$ a quasi-coherent
$\mathcal{O}_Z$-module; then every torsor under $\mathscr{G}$ is trivial.
\end{proposition}

\oldpage[IV-4]{34}
Returning to the problem considered in (16.5.13), we therefore obtain:

\begin{proposition}[16.5.17]
\label{IV.16.5.17}
Let $X$, $Y$ be two $S$-preschemes, $Y_0$ a closed sub-prescheme of $Y$
defined by a quasi-coherent ideal $\mathscr{J}$ of $\mathcal{O}_Y$ such that
$\mathscr{J}^2=0$, and $j:Y_0\to Y$ the canonical injection. Let
$u_0:Y_0\to X$ be an $S$-morphism, and $\mathscr{P}$ the sheaf of sets on
$Y$ such that, for every open subset $U$ of $Y$, $\Gamma(U,\mathscr{P})$ is
the set of $S$-morphisms $u:U\to X$ such that
$u_0|U_0=u\circ(j|U_0)$, where $U_0=j^{-1}(U)$. Then there exists on
$\mathscr{P}$ a structure of pseudo-torsor under the
$\mathcal{O}_{Y_0}$-module
\[
  \mathscr{G}=\mathscr{H}om_{\mathcal{O}_{Y_0}}
  (u_0^*(\Omega_{X/S}^1),\mathscr{J}).
\]
\end{proposition}

In particular:

\begin{corollary}[16.5.18]
\label{IV.16.5.18}
With the notation of (16.5.16), suppose that $Y$ is affine and $\Omega_{X/S}^1$
is of finite presentation; if there exist an open covering $(U_\alpha)$ of $Y$
and, for each index $\alpha$, an $S$-morphism $v_\alpha:U_\alpha\to X$
such that, setting $U_\alpha^0=j^{-1}(U_\alpha)$, we have
$v_\alpha\circ(j|U_\alpha^0)=u_0|U_\alpha^0$, then there exists an
$S$-morphism $u:Y\to X$ such that $u\circ j=u_0$.
\end{corollary}

Indeed, $\mathscr{G}$ is then a quasi-coherent
$\mathcal{O}_{Y_0}$-module (I, 1.3.12); by virtue of (16.5.16) and the fact that
$Y_0$ is then affine, the sheaf $\mathscr{P}$, which by hypothesis is a torsor
under $\mathscr{G}$ and not merely a pseudo-torsor, is \emph{trivial}; but if
$w$ is an isomorphism of $\mathscr{G}$ onto $\mathscr{P}$ (as
torsors under $\mathscr{G}$), the image under $w$ of the section 0 of
$\mathscr{G}$ is the required $S$-morphism $u$.

\subsection{\texorpdfstring{Sheaves of $p$-differentials and the exterior differential.}{Sheaves of p-differentials and the exterior differential.}}
\EGAReaderTarget{ega4.c04.subsec.16-6}{IV.16.6}

\begin{env}[16.6.1]\label{IV.16.6.1}
\EGAReaderTarget{ega4.c04.statement.16-6-1}{IV.16.6.1}
Let $f:X\to S$ be a morphism of preschemes. We call \emph{sheaf of
$p$-differentials of $X$ relative to $S$} ($p$ an integer) the
\emph{$p$-th exterior power} (0, 4.1.5) of the $\mathcal{O}_X$-module
$\Omega_{X/S}^1$, denoted
\[
  \Omega_{X/S}^p=\bigwedge^p(\Omega_{X/S}^1).
  \tag{16.6.1.1}
  \label{IV.16.6.1.1}
\]
We thus have $\Omega_{X/S}^0=\mathcal{O}_X$, and $\Omega_{X/S}^p=0$ for $p<0$;
the $\Omega_{X/S}^p$ are the homogeneous components of the exterior algebra
over $\Omega_{X/S}^1$
\[
  \Omega_{X/S}^{\bullet}=\bigwedge(\Omega_{X/S}^1)
  =\bigoplus_{p\in\mathbf{Z}}\bigwedge^p(\Omega_{X/S}^1),
  \tag{16.6.1.2}
  \label{IV.16.6.1.2}
\]
which is therefore a quasi-coherent graded $\mathcal{O}_X$-algebra,
anticommutative and whose elements of degree 1 are of square zero. For every
affine open subset $U$ of $X$, we have
\[
  \Gamma(U,\Omega_{X/S}^{\bullet})
  =\bigwedge(\Gamma(U,\Omega_{X/S}^1)),
\]
where $\Gamma(U,\Omega_{X/S}^1)$ is regarded as a
$\Gamma(U,\mathcal{O}_X)$-module.

When $S=\operatorname{Spec}(A)$ and $X=\operatorname{Spec}(B)$ are
affine, $B$ thus being an $A$-algebra, we have (0, 4.1.5)
$\Omega_{X/S}^p=(\Omega_{B/A}^p)^{\sim}$, setting
$\Omega_{B/A}^p=\bigwedge^p\Omega_{B/A}^1$.
\end{env}

\begin{theorem}[16.6.2]
\label{IV.16.6.2}
There exists one and only one endomorphism $d$ of the sheaf of additive groups
$\Omega_{X/S}^{\bullet}$ having the following properties:
\begin{enumerate}
\item[(i)] $d\circ d=0$.
\item[(ii)] For every open subset $U$ of $X$ and every section
$f\in\Gamma(U,\mathcal{O}_X)$, we have $df=d_{X/S}f$.
\oldpage[IV-4]{35}
\item[(iii)] For every open subset $U$ of $X$, every pair of integers $p$, $q$
and every pair of sections
$\omega'_p\in\Gamma(U,\Omega_{X/S}^p)$,
$\omega''_q\in\Gamma(U,\Omega_{X/S}^q)$, we have
\[
  d(\omega'_p\wedge\omega''_q)
  =(d\omega'_p)\wedge\omega''_q
  +(-1)^p\omega'_p\wedge d\omega''_q.
  \tag{16.6.2.1}
  \label{IV.16.6.2.1}
\]
\end{enumerate}
Moreover, $d$ is an endomorphism of graded $\psi^*(\mathcal{O}_X)$-modules,
of degree $+1$.
\end{theorem}

Suppose the existence of the endomorphism $d$ proved. For every affine open
subset $U$ of $X$, every section of $\Omega_{X/S}^p$ over $U$ is (by virtue
of (i)) a linear combination of finitely many elements of the form
$g(df_1\wedge df_2\wedge\cdots\wedge df_p)$, where $g$ and the $f_i$ are
sections of $\mathcal{O}_X$ over $U$ (0, 20.4.7). Conditions (i)
and (iii) then show, by induction on $p$, that we necessarily have
\[
  d(g(df_1\wedge df_2\wedge\cdots\wedge df_p))
  =dg\wedge df_1\wedge df_2\wedge\cdots\wedge df_p.
  \tag{16.6.2.2}
  \label{IV.16.6.2.2}
\]
This therefore proves the \emph{uniqueness} of $d$ and the last assertion of the
theorem. By virtue of this uniqueness property, to prove the existence of $d$, we
may restrict to the case where $S=\operatorname{Spec}(A)$ and
$X=\operatorname{Spec}(B)$ are affine. Now (Bourbaki, \emph{Alg.}, chap.
III, 3\textsuperscript{e} éd., \S~10), to define an $A$-antiderivation
$D$ of degree $+1$ of an exterior algebra $\bigwedge(M)$ (where $M$ is a
$B$-module and $B$ an $A$-algebra), this antiderivation taking its values
in an \emph{anticommutative graded} $A$-algebra
$C=\bigoplus_{n=0}^{\infty}C_n$ whose elements of degree 1 are of square
zero, it suffices to give arbitrarily an $A$-derivation $D_0$ of $B$
into $C_1$ and an $A$-homomorphism $D_1$ from $M$ to $C_2$; there then exists
one and only one $A$-antiderivation $D$ of $\bigwedge(M)$ into $C$
coinciding with $D_0$ on $B$ and with $D_1$ on $M$.

In the present case, $D_0$ is necessarily equal to $d_{B/A}$ by virtue of
(ii); taking into account (16.6.2.2), everything comes down to seeing that there
is an $A$-homomorphism $u$ from $\Omega_{B/A}^1$ to $\Omega_{B/A}^2$ such that
\[
  u(g.df)=dg\wedge df
  \tag{16.6.2.3}
  \label{IV.16.6.2.3}
\]
for all $f$, $g$ in $A$; for this it will suffice to show that there
exists an $A$-homomorphism
$v:B\otimes_A\Omega_{B/A}^1\to\Omega_{B/A}^2$ such that
\[
  v(g.\omega)=dg\wedge\omega
  \tag{16.6.2.4}
  \label{IV.16.6.2.4}
\]
for $g\in B$ and $\omega\in\Omega_{B/A}^1$. Finally, since
$\Omega_{B/A}^1=\mathfrak{I}/\mathfrak{I}^2$ (where
$\mathfrak{I}=\mathfrak{I}_{B/A}$ is the kernel of the canonical homomorphism
$B\otimes_A B\to B$) and $\Omega_{B/A}^1$ is generated by the elements
of the form $g.df$, it suffices to define an $A$-homomorphism
$w:B\otimes_A(B\otimes_A B)\to\Omega_{B/A}^2$ such that
\[
  w(g'\otimes g\otimes f)=dg'\wedge(g.df)
  \tag{16.6.2.5}
  \label{IV.16.6.2.5}
\]
and such that $w$ vanishes on the image of $B\otimes_A\mathfrak{I}^2$. Now,
since the right-hand side of (16.6.2.5) is $A$-trilinear in $g'$, $g$, $f$,
the existence of $w$ satisfying (16.6.2.5) is immediate. Since, on the other hand,
$\mathfrak{I}$ is generated by the elements $1\otimes x-x\otimes1$
($x\in B$), we are reduced to verifying that when
$z=(1\otimes x-x\otimes1)(1\otimes y-y\otimes1)$, we have
$w(g'\otimes z)=0$. Now, since
$z=1\otimes(xy)+(xy)\otimes1-x\otimes y-y\otimes x$, formula (16.6.2.4)
shows that it
\oldpage[IV-4]{36}
suffices to see that $d(xy)-x.dy-y.dx=0$, which expresses the fact that $d$ is a
derivation.

It remains to prove that $d$ satisfies condition (i). Now, the square of an
antiderivation is a derivation (Bourbaki, \emph{loc. cit.}), and since
$\Omega_{B/A}^{\bullet}$ is generated by $\Omega_{B/A}^1$ as a
$B$-algebra, it suffices to verify that $d(dz)=0$ for $z\in B$ and for
$z\in\Omega_{B/A}^1$; in the first case, this follows from formula
(16.6.2.3) with $g=1$; in the second, we may restrict to the case where $z=g.df$
with $f$, $g$ in $B$, and then, by virtue of (16.6.2.1) and (16.6.2.3), we have
\[
  d(d(g.df))=d(dg\wedge df)
  =(d(dg))\wedge(df)-(dg)\wedge(d(df))=0.
\]
\hfill Q.E.D.

\begin{definition}[16.6.3]\label{IV.16.6.3}
\EGAReaderTarget{ega4.c04.def.16-6-3}{IV.16.6.3}
The antiderivation $d$ defined in (16.6.2) (also denoted $d_{X/S}$) is
called the \emph{exterior differential on $X$} (relative to $S$).
\end{definition}

\begin{proposition}[16.6.4]
\label{IV.16.6.4}
For every base change $g:S'\to S$, if we set
$X'=X\times_S S'$, the canonical homomorphism
\[
  \Omega_{X/S}^{\bullet}\otimes_S S'
  \longrightarrow\Omega_{X'/S'}^{\bullet}
  \tag{16.6.4.1}
  \label{IV.16.6.4.1}
\]
deduced from the isomorphism (16.5.9.1) is bijective. Moreover, if $s$ is a
section of $\Omega_{X/S}^{\bullet}$ over an open subset $U$ of $X$ and
$s\otimes1$ its inverse image, a section of $\Omega_{X'/S'}^{\bullet}$
over the inverse image $U'$ of $U$ in $X'$, we have
$d_{X'/S'}(s\otimes1)=d_{X/S}(s)\otimes1$.
\end{proposition}

The first assertion is immediate, since the formation of the exterior algebra
of a module commutes with every extension of the ring of scalars. To
prove the second, we may, by virtue of (16.6.2.2), restrict to the case where
$s\in\Gamma(U,\mathcal{O}_X)$, and in this case the assertion has already been
proved (16.4.3.7).

\begin{env}[16.6.5]
\label{IV.16.6.5}
Suppose that $\Omega_{X/S}^1$ is an $\mathcal{O}_X$-module that is locally free
of rank $n$ at a point $x$, so that there are $n$ sections
$s_i\in\Gamma(U,\mathcal{O}_X)$ such that the $ds_i$ form a basis of the
$\Gamma(U,\mathcal{O}_X)$-module $\Gamma(U,\Omega_{X/S}^1)$ (16.5.8).
Then, for every integer $p\geq1$, the $p$-differentials
$ds_{i_1}\wedge ds_{i_2}\wedge\cdots\wedge ds_{i_p}$ (for
$i_1<i_2<\cdots<i_p$ elements of $[1,n]$) form a basis of
$\Gamma(U,\Omega_{X/S}^p)$ over $\Gamma(U,\mathcal{O}_X)$ consisting of
$\binom{n}{p}$ elements. Moreover, formula (16.6.2.2) shows that for
every section $g\in\Gamma(U,\mathcal{O}_X)$ we have
\[
  d(g.ds_{i_1}\wedge ds_{i_2}\wedge\cdots\wedge ds_{i_p})
  =\sum_k(-1)^r\frac{\partial g}{\partial s_k}
  ds_{i_1}\wedge\cdots\wedge ds_{i_r}\wedge ds_k\wedge
  ds_{i_{r+1}}\wedge\cdots\wedge ds_{i_p}
  \tag{16.6.5.1}
  \label{IV.16.6.5.1}
\]
where, on the right-hand side, $k$ runs through the set of the $n-p$ indices
distinct from the $i_h$, $i_r$ being the largest index $<k$.

Note that the relation $d(dg)=0$ for every section
$g\in\Gamma(U,\mathcal{O}_X)$ is expressed in the form
\[
  D_i(D_jg)=D_j(D_ig)\qquad\text{for $i\neq j$};
\]
in other words, the derivations $D_i$ defined in (16.5.7) commute
pairwise.
\end{env}


\subsection{\texorpdfstring{The $\mathscr{P}_{X/S}^n(\mathscr{F})$.}{The P\string_\string{X/S\string}\string^n(F).}}
\EGAReaderTarget{ega4.c04.subsec.16-7}{IV.16.7}

\begin{env}[16.7.1]
\label{IV.16.7.1}
Let $f:X\to S$ be a morphism of preschemes, $\mathscr{F}$ an
$\mathcal{O}_X$-module. Denote by $X_{\Delta_f}^{(n)}$ the \emph{$n$-th
infinitesimal neighbourhood of $X$} for the diagonal morphism
\oldpage[IV-4]{37}
$\Delta_f:X\to X\times_S X$, by
$h_n:X_{\Delta_f}^{(n)}\to X\times_S X$ the canonical morphism (16.1.2), and
consider the two composite morphisms
\[
  p_1^{(n)}:X_{\Delta_f}^{(n)}\xrightarrow{h_n}X\times_S X
  \xrightarrow{p_1}X,qquad
  p_2^{(n)}:X_{\Delta_f}^{(n)}\xrightarrow{h_n}X\times_S X
  \xrightarrow{p_2}X
\]
so that, by definition, $p_1^{(n)}$ corresponds to the homomorphism of
sheaves of rings $\mathcal{O}_X\to\mathscr{P}_{X/S}^n$ that we chose in order
to define the $\mathcal{O}_X$-algebra structure on
$\mathscr{P}_{X/S}^n$ (16.3.5), and $p_2^{(n)}$ to the homomorphism of sheaves
of rings $d_{X/S}^n:\mathcal{O}_X\to\mathscr{P}_{X/S}^n$ (16.3.6). Since
$X_{\Delta_f}^{(n)}$ and $X$ have the same underlying space, we can write
\[
  \mathscr{P}_{X/S}^n
  =(p_1^{(n)})_*((p_2^{(n)})^*(\mathcal{O}_X)).
  \tag{16.7.1.1}
  \label{IV.16.7.1.1}
\]
More generally, we will set
\[
  \mathscr{P}_{X/S}^n(\mathscr{F})
  =(p_1^{(n)})_*((p_2^{(n)})^*(\mathscr{F}))
  \tag{16.7.1.2}
  \label{IV.16.7.1.2}
\]
so that
$\mathscr{P}_{X/S}^n=\mathscr{P}_{X/S}^n(\mathcal{O}_X)$; by definition,
$\mathscr{P}_{X/S}^n(\mathscr{F})$ is therefore an $\mathcal{O}_X$-module.
\end{env}

\begin{env}[16.7.2]
\label{IV.16.7.2}
If we go back to the definitions of inverse images of modules on ringed
spaces (0, 4.3.1) and take into account the fact that
$X_{\Delta_f}^{(n)}$ and $X$ have the same underlying space, we see that the
definition (16.7.1.2) can also be written in the form
\[
  \mathscr{P}_{X/S}^n(\mathscr{F})
  =\mathscr{P}_{X/S}^n\otimes_{\mathcal{O}_X}\mathscr{F}
  \tag{16.7.2.1}
  \label{IV.16.7.2.1}
\]
but where one must take care that, in the interpretation of the sign $\otimes$,
$\mathscr{P}_{X/S}^n$ is endowed with its $\mathcal{O}_X$-module structure
defined by the homomorphism of sheaves of rings
$d_{X/S}^n:\mathcal{O}_X\to\mathscr{P}_{X/S}^n$. It follows immediately from
this formula (or directly from (16.7.1.2)) that
$\mathscr{P}_{X/S}^n(\mathscr{F})$ is canonically endowed with a structure of
$\mathscr{P}_{X/S}^n$-module.
\end{env}

\begin{proposition}[16.7.3]
\label{IV.16.7.3}
\begin{enumerate}
\item[(i)] The functor
$\mathscr{F}\mapsto\mathscr{P}_{X/S}^n(\mathscr{F})$ from the category of
$\mathcal{O}_X$-modules to that of $\mathscr{P}_{X/S}^n$-modules is
right exact and commutes with arbitrary inductive limits; it is exact
when $\mathscr{P}_{X/S}^n$ is a flat $\mathcal{O}_X$-module.
\item[(ii)] If $\mathscr{F}$ is a quasi-coherent $\mathcal{O}_X$-module
(resp. of finite type, resp. of finite presentation),
$\mathscr{P}_{X/S}^n(\mathscr{F})$ is a quasi-coherent
$\mathscr{P}_{X/S}^n$-module (resp. of finite type, resp. of finite
presentation).
\end{enumerate}
\end{proposition}

The assertions of (i) follow immediately from formula (16.7.2.1) and from
consideration of the symmetry of $\mathscr{P}_{X/S}^n$ (16.3.4). The
assertions of (ii) follow from the right exactness of the functor
$\mathscr{F}\mapsto\mathscr{P}_{X/S}^n(\mathscr{F})$.

\begin{env}[16.7.4]
\label{IV.16.7.4}
The two $\mathcal{O}_X$-module structures of $\mathscr{P}_{X/S}^n$
define on $\mathscr{P}_{X/S}^n(\mathscr{F})$ two
$\mathcal{O}_X$-module structures, which moreover commute with each other,
hence an $\mathcal{O}_X$-bimodule structure. It is convenient to write
\emph{on the left} the one of these structures that comes from the structural
homomorphism $\mathcal{O}_X\to\mathscr{P}_{X/S}^n$ (chosen in (16.3.5)) and
\emph{on the right} the one that comes from the homomorphism
$d_{X/S}^n:\mathcal{O}_X\to\mathscr{P}_{X/S}^n$. In other words, for every
open subset $U$ of $X$ and every triple of elements
$a\in\Gamma(U,\mathcal{O}_X)$, $b\in\Gamma(U,\mathscr{P}_{X/S}^n)$,
$t\in\Gamma(U,\mathscr{F})$, we have by definition
\[
  a(b\otimes t)=(ab)\otimes t,qquad
  (b\otimes t)a=(b.d^na)\otimes t=b\otimes(at)
  =(d^na).(b\otimes t).
  \tag{16.7.4.1}
  \label{IV.16.7.4.1}
\]
The $\mathcal{O}_X$-module structure coming from the definition (16.7.1.2)
is therefore, with these conventions, the left $\mathcal{O}_X$-module
structure.

\oldpage[IV-4]{38}
If $\mathscr{F}$ is a quasi-coherent $\mathcal{O}_X$-module, so is
$\mathscr{P}_{X/S}^n(\mathscr{F})$ for either of its
$\mathcal{O}_X$-module structures. If moreover $\mathscr{F}$ is of finite
type (resp. of finite presentation) and $f:X\to S$ is locally of finite type
(resp. locally of finite presentation), then
$\mathscr{P}_{X/S}^n(\mathscr{F})$ is (for either of its
$\mathcal{O}_X$-module structures) of finite type (resp. of finite
presentation), as follows from (16.3.9) and (16.4.22).
\end{env}

\begin{env}[16.7.5]
\label{IV.16.7.5}
The definition (16.7.2.1) implies the existence of a homomorphism of sheaves
of commutative groups
\[
  d_{X/S,\mathscr{F}}^n:\mathscr{F}
  \longrightarrow\mathscr{P}_{X/S}^n(\mathscr{F})
  \qquad\text{(also denoted $d_{X/S}^n$)}
  \tag{16.7.5.1}
  \label{IV.16.7.5.1}
\]
such that, with the notation of (16.7.4), we have
\[
  d_{X/S,\mathscr{F}}^n(t)=1\otimes t
  \tag{16.7.5.2}
  \label{IV.16.7.5.2}
\]
and consequently, by virtue of (16.7.4.1),
\[
  d_{X/S,\mathscr{F}}^n(at)
  =(1\otimes t)a=(d_{X/S,\mathscr{F}}^n(t)).a
  \tag{16.7.5.3}
  \label{IV.16.7.5.3}
\]
\[
  d_{X/S,\mathscr{F}}^n(at)
  =(d_{X/S}^n(a)).(1\otimes t)
  =(d_{X/S}^n(a))(d_{X/S,\mathscr{F}}^n(t)).
  \tag{16.7.5.4}
  \label{IV.16.7.5.4}
\]
It is therefore $\mathcal{O}_X$-linear for the right
$\mathcal{O}_X$-module structure on $\mathscr{P}_{X/S}^n(\mathscr{F})$, and
\emph{semi-linear} (relative to the automorphism $\sigma$ (16.3.4)) for
the left $\mathcal{O}_X$-module structure.
\end{env}

\begin{proposition}[16.7.6]
\label{IV.16.7.6}
The left $\mathcal{O}_X$-module $\mathscr{P}_{X/S}^n(\mathscr{F})$ is
generated by the image of $\mathscr{F}$ under the canonical homomorphism
$d_{X/S,\mathscr{F}}^n$.
\end{proposition}

This follows immediately from (16.7.5.3) and from the particular case
corresponding to $\mathscr{F}=\mathcal{O}_X$ (16.3.8).

\begin{env}[16.7.7]
\label{IV.16.7.7}
The canonical homomorphisms of sheaves of rings
\[
  \varphi_{nm}:\mathscr{P}_{X/S}^m\longrightarrow\mathscr{P}_{X/S}^n
\]
for $n\leq m$ (16.1.2) define, by virtue of (16.7.2.1), canonical
homomorphisms
\[
  \mathscr{P}_{X/S}^m(\mathscr{F})
  \longrightarrow\mathscr{P}_{X/S}^n(\mathscr{F})qquad(n\leq m)
\]
which are homomorphisms of $\mathcal{O}_X$-bimodules by virtue of (16.1.6)
and (16.7.4.1); moreover, we have commutative diagrams
\[
  \xymatrix{
    \mathscr{P}_{X/S}^m(\mathscr{F})\ar[rr]
      &&\mathscr{P}_{X/S}^n(\mathscr{F})\\
    &\mathscr{F}\ar[ul]^-{d_{X/S,\mathscr{F}}^m}
      \ar[ur]_-{d_{X/S,\mathscr{F}}^n}
  }
\]
We thus have a projective system of $\mathcal{O}_X$-bimodules
$(\mathscr{P}_{X/S}^n(\mathscr{F}))$, and we set
\[
  \mathscr{P}_{X/S}^{\infty}(\mathscr{F})
  =\varprojlim\mathscr{P}_{X/S}^n(\mathscr{F}).
  \tag{16.7.7.1}
  \label{IV.16.7.7.1}
\]
Moreover, the foregoing shows that the homomorphisms (16.7.5.1) form a
projective system of homomorphisms, and therefore define a canonical
homomorphism
\[
  d_{X/S,\mathscr{F}}^{\infty}:\mathscr{F}
  \longrightarrow\mathscr{P}_{X/S}^{\infty}(\mathscr{F}).
  \tag{16.7.7.2}
  \label{IV.16.7.7.2}
\]
\end{env}

\oldpage[IV-4]{39}

\begin{env}[16.7.8]
\label{IV.16.7.8}
Let $\mathscr{F}$, $\mathscr{G}$ be two $\mathcal{O}_X$-modules; it follows
immediately from the definition (16.7.2.1) that we have a canonical
isomorphism of $\mathscr{P}_{X/S}^n$-modules
\[
  \mathscr{P}_{X/S}^n(\mathscr{F}\otimes_{\mathcal{O}_X}\mathscr{G})
  \xrightarrow{\sim}
  \mathscr{P}_{X/S}^n(\mathscr{F})
  \otimes_{\mathscr{P}_{X/S}^n}\mathscr{P}_{X/S}^n(\mathscr{G})
  \tag{16.7.8.1}
  \label{IV.16.7.8.1}
\]
(Bourbaki, \emph{Alg.}, chap.~II, 3\textsuperscript{e} éd., §~5,
n\textsuperscript{o}~1, prop.~3).

We conclude from this in particular (or we see directly from the definition
(16.7.2.1)) that if $\mathscr{F}$ is endowed with a structure of
(not necessarily associative) $\mathcal{O}_X$-algebra, then
$\mathscr{P}_{X/S}^n(\mathscr{F})$ is canonically endowed with a structure of
$\mathscr{P}_{X/S}^n$-algebra; the latter is associative (resp.
commutative, resp. unital, resp. a Lie algebra) when $\mathscr{F}$ is.
Moreover, the canonical homomorphisms
$\mathscr{P}_{X/S}^m(\mathscr{F})\to\mathscr{P}_{X/S}^n(\mathscr{F})$ for
$n\leq m$ (16.7.7) are then di-homomorphisms of algebras; likewise,
(16.7.5.1) is then a homomorphism of $\mathcal{O}_X$-algebras when
$\mathscr{P}_{X/S}^n(\mathscr{F})$ is endowed with its
$\mathcal{O}_X$-algebra structure coming from its \emph{right}
$\mathcal{O}_X$-module structure.

With the same notation, we also have a canonical homomorphism of
$\mathscr{P}_{X/S}^n$-modules
\[
  \mathscr{P}_{X/S}^n
  (\mathscr{H}\!om_{\mathcal{O}_X}(\mathscr{F},\mathscr{G}))
  \longrightarrow
  \mathscr{H}\!om_{\mathscr{P}_{X/S}^n}
  (\mathscr{P}_{X/S}^n(\mathscr{F}),\mathscr{P}_{X/S}^n(\mathscr{G}))
  \tag{16.7.8.2}
  \label{IV.16.7.8.2}
\]
(Bourbaki, \emph{Alg.}, 3\textsuperscript{e} éd., §~5,
n\textsuperscript{o}~3), which is bijective when $\mathscr{P}_{X/S}^n$ is
an $\mathcal{O}_X$-module that is \emph{locally free of finite type}
(\emph{loc. cit.}, prop.~7).
\end{env}

\begin{env}[16.7.9]
\label{IV.16.7.9}
Suppose we are in the situation described in (16.4.1); then from the
canonical homomorphism $P^n(u)$ (16.4.3.3) we immediately deduce a
canonical homomorphism of $\mathcal{O}_{X'}$-bimodules
\[
  u^*(\mathscr{P}_{X/S}^n(\mathscr{F}))
  \longrightarrow
  \mathscr{P}_{X'/S'}^n(u^*(\mathscr{F})).
  \tag{16.7.9.1}
  \label{IV.16.7.9.1}
\]
We leave it to the reader to extend to this homomorphism the properties
seen in (16.4) for the case $\mathscr{F}=\mathcal{O}_X$.
\end{env}

\begin{remark}[16.7.10]
\label{IV.16.7.10}
The definition of $\mathscr{P}_{X/S}^n(\mathscr{F})$ in the form (16.7.1.2)
still makes sense when $\mathscr{F}$ is an arbitrary sheaf of sets
(the inverse image of a sheaf of sets under $p_2^{(n)}$ being defined
in (0\textsubscript{I}, 3.7.1)); a variant of this definition makes it
possible to define the ``jet scheme'' (relative to $S$) of an arbitrary
$X$-prescheme.
\end{remark}

\subsection{Differential operators.}
\footnote{For a more general formalism, see Exposé VII of [42]
(due to P. Gabriel).}
\EGAReaderTarget{ega4.c04.subsec.16-8}{IV.16.8}

\begin{definition}[16.8.1]\label{IV.16.8.1}
\EGAReaderTarget{ega4.c04.def.16-8-1}{IV.16.8.1}
Let $f=(\psi,\theta):X\to S$ be a morphism of preschemes,
$\mathscr{F}$, $\mathscr{G}$ two $\mathcal{O}_X$-modules, $n$ an integer
$\geq0$. We say that a homomorphism of sheaves of additive groups
$D:\mathscr{F}\to\mathscr{G}$ is a \emph{differential operator of order
$\leq n$} (relative to $S$) if there exists a homomorphism of
$\mathcal{O}_X$-modules
$u:\mathscr{P}_{X/S}^n(\mathscr{F})\to\mathscr{G}$ (where
$\mathscr{P}_{X/S}^n(\mathscr{F})$ is endowed with its left
$\mathcal{O}_X$-module structure (16.7.4)) such that
$D=u\circ d_{X/S,\mathscr{F}}^n$.
\end{definition}

It is clear, by virtue of the existence of the canonical homomorphisms
\[
  \mathscr{P}_{X/S}^m(\mathscr{F})
  \longrightarrow\mathscr{P}_{X/S}^n(\mathscr{F})
\]
\oldpage[IV-4]{40}
for $n\leq m$ (16.7.7), that a differential operator of order $\leq n$ is
also a differential operator of order $\leq m$ for every $m\geq n$. If
$D:\mathscr{F}\to\mathscr{G}$ is a differential operator of order
$\leq n$, then, for every open subset $U$ of $X$,
$D|U:\mathscr{F}|U\to\mathscr{G}|U$ is also a differential operator
of order $\leq n$.

We say that a homomorphism $D:\mathscr{F}\to\mathscr{G}$ of the sheaves of
additive groups underlying $\mathscr{F}$ and $\mathscr{G}$ is a
\emph{differential operator} (relative to $S$) if for every $x\in X$ there
exist an open neighbourhood $U$ of $x$ and an integer $n\geq0$ such that
$D|U:\mathscr{F}|U\to\mathscr{G}|U$ is a differential operator of order
$\leq n$. The \emph{order} of a differential operator
$D:\mathscr{F}\to\mathscr{G}$ is the infimum of the integers $n$ such
that $D$ is a differential operator of order $\leq n$ (and hence $+\infty$
if there are no such integers); this order is always finite if $X$ is
\emph{quasi-compact}. The differential operators of order $0$ are none other
than the homomorphisms of $\mathcal{O}_X$-modules
$\mathscr{F}\to\mathscr{G}$; by convention, every differential operator
of order $<0$ is \emph{zero}. For $n\geq0$, a differential operator is
not in general a homomorphism of $\mathcal{O}_X$-modules, but it is always
a homomorphism of $\psi^*(\mathcal{O}_S)$-modules.

When $\mathscr{F}=\mathcal{O}_X$, a differential operator of order
$\leq1$ from $\mathcal{O}_X$ to $\mathscr{G}$ can be written in exactly one
way in the form $v+D$, where $v:\mathcal{O}_X\to\mathscr{G}$ is an
$\mathcal{O}_X$-homomorphism and $D$ an $S$-\emph{derivation} (16.5.1) of
$\mathcal{O}_X$ into $\mathscr{G}$: this follows from the structure of
$P_{B/A}^1$ (0, 20.4.8).

\begin{env}[16.8.2]
\label{IV.16.8.2}
To describe more explicitly a differential operator of order
$\leq n$, $D:\mathscr{F}\to\mathscr{G}$, it suffices, for every affine open
subset $U$ of $X$ whose image in $S$ is contained in an affine open subset
$V$, to characterize the homomorphism
$D=D_U:\Gamma(U,\mathscr{F})\to\Gamma(U,\mathscr{G})$. If we set
$\Gamma(V,\mathcal{O}_S)=A$, $\Gamma(U,\mathcal{O}_X)=B$, so that $B$
is an $A$-algebra, we have
\[
  \Gamma(U,\mathscr{P}_{X/S}^n)
  =(B\otimes_A B)/\mathfrak{I}^{n+1},
\]
where we have set $\mathfrak{I}=\mathfrak{I}_{B/A}$ for brevity. Let us
further set $M=\Gamma(U,\mathscr{F})$, $N=\Gamma(U,\mathscr{G})$; then the
definition of $D$ means that for each pair $(U,V)$ satisfying the
preceding conditions, the $A$-homomorphism $D:M\to N$ factors as
\[
  M\longrightarrow
  ((B\otimes_A B)/\mathfrak{I}^{n+1})\otimes_B M
  \xrightarrow{v}N
\]
where the first arrow is the canonical homomorphism $t\mapsto1\otimes t$, and
$v$ is a $B$-homomorphism, the $B$-module structure of
$((B\otimes_A B)/\mathfrak{I}^{n+1})\otimes_B M$ coming from the first
factor $B$ (whereas, we recall, in the formation of the tensor product
over $B$, the $B$-module structure of
$(B\otimes_A B)/\mathfrak{I}^{n+1}$ comes from the second factor $B$). Let us
now note that the $B$-module
$((B\otimes_A B)/\mathfrak{I}^{n+1})\otimes_B M$ is isomorphic to
$(B\otimes_A M)/\mathfrak{I}^{n+1}(B\otimes_A M)$, where $B\otimes_A M$ is
regarded as a $(B\otimes_A B)$-module and its $B$-module structure
comes from the homomorphism $b\mapsto b\otimes1$ from $B$ to $B\otimes_A B$.
Let then $D'$ be the $B$-homomorphism from $B\otimes_A M$ to $N$ such that
$D'(b\otimes t)=bD(t)$; the factorization condition on $D$ can also be
expressed by saying that $D'$ must \emph{be zero on the $B$-module}
$\mathfrak{I}^{n+1}(B\otimes_A M)$.
\end{env}

\begin{env}[16.8.3]
\label{IV.16.8.3}
It is clear that the set of differential operators of order $\leq n$ from
$\mathscr{F}$ to $\mathscr{G}$ is an additive group, denoted
$\operatorname{Diff}_{X/S}^n(\mathscr{F},\mathscr{G})$; when
$\mathscr{F}=\mathscr{G}=\mathcal{O}_X$, we also write
$\operatorname{Diff}_{X/S}^n$ instead of
$\operatorname{Diff}_{X/S}^n(\mathcal{O}_X,\mathcal{O}_X)$.

We have seen (16.8.1) that for two open subsets $U\supset V$ of $X$ we have
a canonical restriction homomorphism
\[
  \operatorname{Diff}_{U/S}^n(\mathscr{F}|U,\mathscr{G}|U)
  \longrightarrow
  \operatorname{Diff}_{V/S}^n(\mathscr{F}|V,\mathscr{G}|V)
\]
\oldpage[IV-4]{41}
hence $U\mapsto\operatorname{Diff}_{U/S}^n(\mathscr{F}|U,\mathscr{G}|U)$ is
a presheaf of additive groups; in fact it is even a \emph{sheaf}, because,
for $U$ varying over the open subsets of $X$, the homomorphisms
$u\mapsto u\circ d_{U/S,\mathscr{F}|U}^n$ are \emph{isomorphisms} of
additive groups
\[
  \mathscr{H}\!om_{\mathcal{O}_U}
  (\mathscr{P}_{U/S}^n(\mathscr{F}|U),\mathscr{G}|U)
  \xrightarrow{\sim}
  \operatorname{Diff}_{U/S}^n(\mathscr{F}|U,\mathscr{G}|U),
  \tag{16.8.3.1}
  \label{IV.16.8.3.1}
\]
by virtue of the fact that the image of $\mathscr{F}$ under
$d_{X/S,\mathscr{F}}^n$ generates $\mathscr{P}_{X/S}^n(\mathscr{F})$
(16.7.6). We denote this sheaf by
$\mathscr{D}\!iff_{X/S}^n(\mathscr{F},\mathscr{G})$, and we therefore have:
\end{env}

\begin{proposition}[16.8.4]
\label{IV.16.8.4}
The isomorphisms (16.8.3.1) define an isomorphism of sheaves of
additive groups
\[
  \mathscr{H}\!om_{\mathcal{O}_X}
  (\mathscr{P}_{X/S}^n(\mathscr{F}),\mathscr{G})
  \xrightarrow{\sim}
  \mathscr{D}\!iff_{X/S}^n(\mathscr{F},\mathscr{G}).
  \tag{16.8.4.1}
  \label{IV.16.8.4.1}
\]
\end{proposition}

When $\mathscr{F}=\mathscr{G}=\mathcal{O}_X$, we also write
$\mathscr{D}\!iff_{X/S}^n$ instead of
$\mathscr{D}\!iff_{X/S}^n(\mathcal{O}_X,\mathcal{O}_X)$; it follows from
(16.8.4) that $\mathscr{D}\!iff_{X/S}^n$ is canonically identified with the
\emph{dual} of the $\mathcal{O}_X$-module $\mathscr{P}_{X/S}^n$; accordingly,
we write $\langle t,D\rangle$ instead of $u(t)$ if $t$ is a section of
$\mathscr{P}_{X/S}^n$ over an open subset and $u$ is the homomorphism from
$\mathscr{P}_{X/S}^n$ to $\mathcal{O}_X$ corresponding to $D$.

\begin{env}[16.8.5]
\label{IV.16.8.5}
Since $\mathscr{P}_{X/S}^n(\mathscr{F})$ is endowed with an
$\mathcal{O}_X$-bimodule structure (16.7.4), we canonically deduce from it an
$\mathcal{O}_X$-bimodule structure on
$\mathscr{H}\!om_{\mathcal{O}_X}
(\mathscr{P}_{X/S}^n(\mathscr{F}),\mathscr{G})$, hence also on
$\mathscr{D}\!iff_{X/S}^n(\mathscr{F},\mathscr{G})$ by virtue of (16.8.4.1).
Precisely, to the left $\mathcal{O}_X$-module structure of
$\mathscr{P}_{X/S}^n(\mathscr{F})$ there corresponds, by virtue of the
definition (16.8.1), the left $\mathcal{O}_X$-module structure on
$\mathscr{D}\!iff_{X/S}^n(\mathscr{F},\mathscr{G})$ made explicit as follows:
for every open subset $U$ of $X$, every section
$a\in\Gamma(U,\mathcal{O}_X)$ and every differential operator
$D:\mathscr{F}|U\to\mathscr{G}|U$, $aD$ is the differential operator which
assigns to every section $t\in\Gamma(U,\mathscr{F})$ the section
\[
  (aD)(t)=a(D(t))
  \tag{16.8.5.1}
  \label{IV.16.8.5.1}
\]
of $\Gamma(U,\mathscr{G})$. Likewise, to the right
$\mathcal{O}_X$-module structure of $\mathscr{P}_{X/S}^n(\mathscr{F})$
there corresponds the right $\mathcal{O}_X$-module structure on
$\mathscr{D}\!iff_{X/S}^n(\mathscr{F},\mathscr{G})$ made explicit as follows:
with the same notation as above, $Da$ is the differential operator which
assigns to $t\in\Gamma(U,\mathscr{F})$ the section
\[
  (Da)(t)=D(at).
  \tag{16.8.5.2}
  \label{IV.16.8.5.2}
\]
\end{env}

\begin{proposition}[16.8.6]
\label{IV.16.8.6}
If $f:X\to S$ is a morphism locally of finite presentation,
$\mathscr{F}$ a quasi-coherent $\mathcal{O}_X$-module of finite presentation
and $\mathscr{G}$ a quasi-coherent $\mathcal{O}_X$-module, then
$\mathscr{D}\!iff_{X/S}^n(\mathscr{F},\mathscr{G})$ is a quasi-coherent
$\mathcal{O}_X$-module for either of the structures defined
in (16.8.5).
\end{proposition}

The proposition follows from the fact that, under the hypotheses made,
$\mathscr{P}_{X/S}^n(\mathscr{F})$ is a quasi-coherent $\mathcal{O}_X$-module
of finite presentation (16.7.4), and from (I, 1.3.12).

\begin{env}[16.8.7]
\label{IV.16.8.7}
The set of differential operators from $\mathscr{F}$ to $\mathscr{G}$
(of unspecified order (16.8.1)) is denoted
$\operatorname{Diff}_{X/S}(\mathscr{F},\mathscr{G})$; we see again, as
in (16.8.3), that
$U\mapsto\operatorname{Diff}_{U/S}(\mathscr{F}|U,\mathscr{G}|U)$ is a
sheaf of additive groups, which we will denote by
$\mathscr{D}\!iff_{X/S}(\mathscr{F},\mathscr{G})$. It is immediate that
$\mathscr{D}\!iff_{X/S}(\mathscr{F},\mathscr{G})$ is the union of the
increasing filtered family of its sub-sheaves
$\mathscr{D}\!iff_{X/S}^n(\mathscr{F},\mathscr{G})$; if $X$ is
quasi-compact, $\operatorname{Diff}_{X/S}(\mathscr{F},\mathscr{G})$
\oldpage[IV-4]{42}
is likewise the union of its subgroups
$\operatorname{Diff}_{X/S}^n(\mathscr{F},\mathscr{G})$ (16.8.1). The
$\mathcal{O}_X$-bimodule structures on the
$\mathscr{D}\!iff_{X/S}^n(\mathscr{F},\mathscr{G})$ therefore define an
$\mathcal{O}_X$-bimodule structure on
$\mathscr{D}\!iff_{X/S}(\mathscr{F},\mathscr{G})$, again made explicit by
(16.8.5.1) and (16.8.5.2).

Note that, for $n\leq m$, we have a commutative diagram
\[
  \xymatrix{
    \mathscr{H}\!om_{\mathcal{O}_X}
      (\mathscr{P}_{X/S}^n(\mathscr{F}),\mathscr{G})
      \ar[r]^-{\sim}\ar[d]
      &\mathscr{D}\!iff_{X/S}^n(\mathscr{F},\mathscr{G})\ar[d]\\
    \mathscr{H}\!om_{\mathcal{O}_X}
      (\mathscr{P}_{X/S}^m(\mathscr{F}),\mathscr{G})
      \ar[r]^-{\sim}
      &\mathscr{D}\!iff_{X/S}^m(\mathscr{F},\mathscr{G})
  }
  \tag{16.8.7.1}
  \label{IV.16.8.7.1}
\]
where the horizontal arrows are the isomorphisms (16.8.4.2) and the left
vertical arrow comes from the canonical homomorphism
$\mathscr{P}_{X/S}^m(\mathscr{F})\to
\mathscr{P}_{X/S}^n(\mathscr{F})$ (16.7.7). For every open subset $U$ of $X$,
let us then endow
\[
  \Gamma(U,\mathscr{P}_{X/S}^{\infty}(\mathscr{F}))
  =\varprojlim\Gamma(U,\mathscr{P}_{X/S}^n(\mathscr{F}))
\]
with the projective limit topology of the discrete topologies on the
$\Gamma(U,\mathscr{P}_{X/S}^n(\mathscr{F}))$, which defines
$\Gamma(U,\mathscr{P}_{X/S}^{\infty}(\mathscr{F}))$ as a topological
$\Gamma(U,\mathcal{O}_X)$-bimodule, so that
$\mathscr{P}_{X/S}^{\infty}(\mathscr{F})$ appears as a sheaf with values
in the category of topological commutative groups
(0\textsubscript{I}, 3.2.6). Then (G, II, 1.11) the limit of the inductive
system of sheaves of commutative groups
\[
  (\mathscr{H}\!om_{\mathcal{O}_X}
  (\mathscr{P}_{X/S}^n(\mathscr{F}),\mathscr{G}))
\]
is none other than the sheaf of germs of \emph{continuous} homomorphisms from
$\mathscr{P}_{X/S}^{\infty}(\mathscr{F})$ to $\mathscr{G}$ (the latter
being endowed with the discrete topology): indeed, the continuous
homomorphisms from $\Gamma(U,\mathscr{P}_{X/S}^{\infty}(\mathscr{F}))$ to the
discrete group $\Gamma(U,\mathscr{G})$ correspond bijectively to the
inductive systems of group homomorphisms
$\Gamma(U,\mathscr{P}_{X/S}^n(\mathscr{F}))\to\Gamma(U,\mathscr{G})$. We
can therefore also express (16.8.4) by saying that we have a canonical
isomorphism
\[
  \mathscr{H}\!om.\mathrm{cont}_{\mathcal{O}_X}
  (\mathscr{P}_{X/S}^{\infty}(\mathscr{F}),\mathscr{G})
  \xrightarrow{\sim}
  \mathscr{D}\!iff_{X/S}(\mathscr{F},\mathscr{G})
\]
where the left-hand side denotes the sheaf of germs of continuous
homomorphisms from $\mathscr{P}_{X/S}^{\infty}(\mathscr{F})$ to $\mathscr{G}$.
\end{env}

\begin{proposition}[16.8.8]
\label{IV.16.8.8}
Let $\mathscr{F}$, $\mathscr{G}$ be two $\mathcal{O}_X$-modules,
$D:\mathscr{F}\to\mathscr{G}$ a homomorphism of
$\psi^*(\mathcal{O}_S)$-modules, $n$ an integer $\geq0$. The following
conditions are equivalent:
\begin{enumerate}
  \item[(a)] $D$ is a differential operator of order $\leq n$.
  \item[(b)] For every section $a$ of $\mathcal{O}_X$ over an open subset
  $U$, the homomorphism $D_a:\mathscr{F}|U\to\mathscr{G}|U$ such that, for
  every section $t$ of $\mathscr{F}$ over an open subset $V\subset U$, we
  have
  \[
    D_a(t)=D(at)-aD(t)
    \tag{16.8.8.1}
    \label{IV.16.8.8.1}
  \]
  is a differential operator of order $\leq n-1$.
  \item[(c)] For every open subset $U$ of $X$, every family
  $(a_i)_{1\leq i\leq n+1}$ of $n+1$ sections of $\mathcal{O}_X$ over
  $U$ and every section $t$ of $\mathscr{F}$ over $U$, we have
  the identity
  \[
    \sum_{H\subset I_{n+1}}(-1)^{\operatorname{Card}(H)}
    \left(\prod_{i\in H}a_i\right)
    D\!\left(\left(\prod_{i\notin H}a_i\right)t\right)=0
    \tag{16.8.8.2}
    \label{IV.16.8.8.2}
  \]
  (where $I_{n+1}$ is the interval $1\leq i\leq n+1$ of $\mathbb{N}$).
\end{enumerate}
\oldpage[IV-4]{43}
\end{proposition}

Let us first prove the equivalence of a) and c). By definition, to prove that
$D$ is a differential operator of order $\leq n$, it suffices to show that
this is so for the restriction
$D|U:\mathscr{F}|U\to\mathscr{G}|U$ to every affine open subset $U$ of $X$,
and on the other hand property c) holds for every open subset $U$ of $X$ if
it holds for every affine open subset. We may therefore restrict to the case
where $S=\operatorname{Spec}(A)$ and $X=\operatorname{Spec}(B)$ are affine.
By virtue of (16.8.2) (whose notation we keep), condition a) then means
that the $A$-homomorphism $D':B\otimes_A M\to N$ such that
$D'(b\otimes t)=bD(t)$ vanishes on
$\mathfrak{I}^{n+1}(B\otimes_A M)$, which, by virtue of (0, 20.4.4), amounts
to saying that $D'$ vanishes on all elements of the form
\[
  \left(\prod_{i=1}^{n+1}(a_i\otimes1-1\otimes a_i)\right).(1\otimes t)
\]
where $a_i\in B$ and $t\in M$. Now this element can be written
\[
  \sum_{H\subset I_{n+1}}
  \left(\prod_{i\in H}a_i\right)\otimes
  \left(\left(\prod_{i\notin H}a_i\right)t\right),
\]
and the value of $D'$ at this element is none other than the left-hand side
of (16.8.8.2), which proves the equivalence of a) and c).

Let us now prove the equivalence of b) and c). We argue by induction on
$n$, the assertion being trivial for $n=0$. Writing $a_{n+1}$ instead of $a$
in condition b), we see, by virtue of the induction hypothesis, that
condition b) means that for every family $(a_i)_{1\leq i\leq n}$ of $n$
sections of $\mathcal{O}_X$ over $U$ and every section $t$ of
$\mathscr{F}$ over $U$, we have
\[
  \sum_{H'\subset I_n}(-1)^{\operatorname{Card}(H')}
  \left(\prod_{i\in H'}a_i\right)
  D_{a_{n+1}}\!\left(\left(\prod_{i\notin H'}a_i\right)t\right)=0.
\]
But if in this relation we replace $D_{a_{n+1}}$ by its definition
(16.8.8.1), we immediately find that we obtain, up to sign, the left-hand
side of (16.8.8.2); whence the conclusion.

\begin{proposition}[16.8.9]
\label{IV.16.8.9}
If $D:\mathscr{F}\to\mathscr{G}$ is a differential operator of order
$\leq n$, and $D':\mathscr{G}\to\mathscr{H}$ a differential operator of order
$\leq n'$, then $D'\circ D:\mathscr{F}\to\mathscr{H}$ is a differential
operator of order $\leq n+n'$.
\end{proposition}

By hypothesis, we can write $D=u\circ d_{X/S,\mathscr{F}}^n$ and
$D'=v\circ d_{X/S,\mathscr{G}}^{n'}$, where
$u:\mathscr{P}_{X/S}^n\otimes_{\mathcal{O}_X}\mathscr{F}\to\mathscr{G}$ and
$v:\mathscr{P}_{X/S}^{n'}\otimes_{\mathcal{O}_X}\mathscr{G}\to\mathscr{H}$
are $\mathcal{O}_X$-homomorphisms. Everything comes down to showing that
the composite homomorphism of sheaves of additive groups
\[
  \xymatrix{
    \mathscr{F}\ar[r]^-{d_{X/S,\mathscr{F}}^n}
      &\mathscr{P}_{X/S}^n\otimes_{\mathcal{O}_X}\mathscr{F}
        \ar[r]^-u
      &\mathscr{G}\ar[r]^-{d_{X/S,\mathscr{G}}^{n'}}
      &\mathscr{P}_{X/S}^{n'}\otimes_{\mathcal{O}_X}\mathscr{G}
  }
\]
factors as
\[
  \xymatrix{
    \mathscr{F}\ar[r]^-{d_{X/S,\mathscr{F}}^{n+n'}}
      &\mathscr{P}_{X/S}^{n+n'}\otimes_{\mathcal{O}_X}\mathscr{F}
        \ar[r]^-w
      &\mathscr{P}_{X/S}^{n'}\otimes_{\mathcal{O}_X}\mathscr{G}
  }
\]
where $w$ is an $\mathcal{O}_X$-homomorphism. It will suffice to prove the
following

\begin{lemma}[16.8.9.1]
\label{IV.16.8.9.1}
There exists one and only one $\mathcal{O}_X$-homomorphism
\[
  \delta:\mathscr{P}_{X/S}^{n+n'}
  \longrightarrow
  \mathscr{P}_{X/S}^{n'}(\mathscr{P}_{X/S}^n)
  =\mathscr{P}_{X/S}^{n'}\otimes_{\mathcal{O}_X}\mathscr{P}_{X/S}^n
  \tag{16.8.9.2}
  \label{IV.16.8.9.2}
\]
\oldpage[IV-4]{44}
making the diagram
\[
  \xymatrix{
    \mathcal{O}_X\ar[r]^-{d_{X/S}^{n+n'}}\ar[d]_-{d_{X/S}^n}
      &\mathscr{P}_{X/S}^{n+n'}\ar[d]^-\delta\\
    \mathscr{P}_{X/S}^n
      \ar[r]_-{d_{X/S,\mathscr{P}_{X/S}^n}^{n'}}
      &\mathscr{P}_{X/S}^{n'}(\mathscr{P}_{X/S}^n)
  }
  \tag{16.8.9.3}
  \label{IV.16.8.9.3}
\]
commutative.
\end{lemma}

Indeed, we will then have a commutative diagram deduced from (16.8.9.3) by
tensoring with $\mathscr{F}$
\[
  \xymatrix{
    \mathscr{F}\ar[r]^-{d_{X/S,\mathscr{F}}^{n+n'}}
      \ar[d]_-{d_{X/S,\mathscr{F}}^n}
      &\mathscr{P}_{X/S}^{n+n'}(\mathscr{F})
        \ar[d]^-{\delta\otimes1}\\
    \mathscr{P}_{X/S}^n(\mathscr{F})
      \ar[r]_-{d_{X/S,\mathscr{P}_{X/S}^n(\mathscr{F})}^{n'}}
      &\mathscr{P}_{X/S}^{n'}(\mathscr{P}_{X/S}^n(\mathscr{F}))
  }
\]
and, on the other hand, we verify immediately from the definition (16.7.5)
that the diagram
\[
  \xymatrix{
    \mathscr{P}_{X/S}^n(\mathscr{F})\ar[r]^-u
      \ar[d]_-{d_{X/S,\mathscr{P}_{X/S}^n(\mathscr{F})}^{n'}}
      &\mathscr{G}\ar[d]^-{d_{X/S,\mathscr{G}}^{n'}}\\
    \mathscr{P}_{X/S}^{n'}(\mathscr{P}_{X/S}^n(\mathscr{F}))
      \ar[r]_-{1\otimes u}
      &\mathscr{P}_{X/S}^{n'}(\mathscr{G})
  }
\]
is commutative. We will therefore answer the question by taking for $w$ the
composite $\mathcal{O}_X$-homomorphism
\[
  \xymatrix{
    \mathscr{P}_{X/S}^{n+n'}(\mathscr{F})
      \ar[r]^-{\delta\otimes1}
      &\mathscr{P}_{X/S}^{n'}(\mathscr{P}_{X/S}^n(\mathscr{F}))
        \ar[r]^-{1\otimes u}
      &\mathscr{P}_{X/S}^{n'}(\mathscr{G}).
  }
\]

It remains to prove Lemma (16.8.9.1). Taking into account (16.7.6), which
proves the uniqueness of $\delta$, we are reduced to the case where
$S=\operatorname{Spec}(A)$ and $X=\operatorname{Spec}(B)$ are affine; setting
$\mathfrak{I}=\mathfrak{I}_{B/A}$, we have to define a canonical
homomorphism of $B$-modules
\[
  \varphi:(B\otimes_A B)/\mathfrak{I}^{n+n'+1}
  \longrightarrow
  ((B\otimes_A B)/\mathfrak{I}^{n'+1})
  \otimes_B((B\otimes_A B)/\mathfrak{I}^{n+1})
\]
the $B$-module structures of the two sides coming from the first factor
$B$; recall that in the tensor product on the right-hand side,
$(B\otimes_A B)/\mathfrak{I}^{n'+1}$ must be regarded
\oldpage[IV-4]{45}
as a right $B$-module via its second factor $B$, and
$(B\otimes_A B)/\mathfrak{I}^{n+1}$ as a left $B$-module via its first
factor $B$ (16.7.2). It amounts to the same to define a homomorphism of
$B$-modules
\[
  \varphi_0:B\otimes_A B\longrightarrow
  ((B\otimes_A B)/\mathfrak{I}^{n'+1})
  \otimes_B((B\otimes_A B)/\mathfrak{I}^{n+1})
\]
and to prove that it vanishes on $\mathfrak{I}^{n+n'+1}$. Now such a
homomorphism is immediately defined by the condition that
\[
  \varphi_0(b\otimes b')
  =\pi_{n'}(b\otimes1)\otimes\pi_n(1\otimes b')
  \qquad\text{for $b,b'$ in $B$}
\]
with the notation of (16.3.7). Moreover, it is immediate that $\varphi_0$ is
a homomorphism of \emph{rings}. Now we can write
\[
  \begin{aligned}
  \varphi_0(b\otimes1-1\otimes b)
  ={}&\pi_{n'}(b\otimes1-1\otimes b)\otimes\pi_n(1\otimes1)\\
  &+\pi_{n'}(1\otimes b)\otimes\pi_n(1\otimes1)
  -\pi_{n'}(1\otimes1)\otimes\pi_n(1\otimes b)
  \end{aligned}
\]
and we have
\[
  \begin{aligned}
  \pi_{n'}(1\otimes b)\otimes\pi_n(1\otimes1)
  &=\pi_{n'}(1\otimes1)b\otimes\pi_n(1\otimes1)\\
  &=\pi_{n'}(1\otimes1)\otimes b\pi_n(1\otimes1)\\
  &=\pi_{n'}(1\otimes1)\otimes\pi_n(b\otimes1)
  \end{aligned}
\]
whence finally
\[
  \begin{aligned}
  \varphi_0(b\otimes1-1\otimes b)
  ={}&\pi_{n'}(b\otimes1-1\otimes b)\otimes\pi_n(1\otimes1)\\
  &+\pi_{n'}(1\otimes1)\otimes\pi_n(b\otimes1-1\otimes b).
  \end{aligned}
  \tag{16.8.9.4}
  \label{IV.16.8.9.4}
\]
A product of $n+n'+1$ terms of the form (16.8.9.4) is therefore
necessarily zero, since this is so for a product of $n+1$ terms of the
form $\pi_n(b\otimes1-1\otimes b)$ and for a product of $n'+1$ terms of the
form $\pi_{n'}(b\otimes1-1\otimes b)$. The conclusion therefore follows from
(0, 20.4.4).

\begin{corollary}[16.8.10]
\label{IV.16.8.10}
The sheaf $\mathscr{D}\!iff_{X/S}(\mathcal{O}_X,\mathcal{O}_X)$ (also denoted
$\mathscr{D}\!iff_{X/S}$) is canonically endowed with a structure of sheaf
of rings, the $\mathscr{D}\!iff_{X/S}^n$ forming an increasing filtration
compatible with this structure.
\end{corollary}

In particular, $\mathscr{D}\!iff_{X/S}^0$ is a sheaf of subrings of
$\mathscr{D}\!iff_{X/S}$, which is canonically identified with $\mathcal{O}_X$
(16.8.1). The formulas (16.8.5.1) and (16.8.5.2) show that the
$\mathcal{O}_X$-bimodule structure of $\mathscr{D}\!iff_{X/S}$ comes from
left and right multiplication by the sections of $\mathcal{O}_X$,
regarded as a sheaf of subrings of $\mathscr{D}\!iff_{X/S}$.

\begin{remarks}[16.8.11]
\label{IV.16.8.11}
\begin{enumerate}
  \item[(i)] Suppose that
  $\mathscr{F}=\bigoplus_{\lambda\in L}\mathscr{F}_\lambda$; then it is
  clear (16.7.2.1) that
  $\mathscr{P}_{X/S}^n(\mathscr{F})=
  \bigoplus_{\lambda\in L}\mathscr{P}_{X/S}^n(\mathscr{F}_\lambda)$; since the
  functor $\mathscr{F}\mapsto\Gamma(U,\mathscr{F})$ commutes with the formation
  of arbitrary direct sums, $d_{X/S,\mathscr{F}}^n$ is
  the homomorphism whose restriction to each $\mathscr{F}_\lambda$ is
  \[
    d_{X/S,\mathscr{F}_\lambda}^n:\mathscr{F}_\lambda
    \longrightarrow\mathscr{P}_{X/S}^n(\mathscr{F}_\lambda);
  \]
  we conclude immediately that we have
  \[
    \operatorname{Diff}_{X/S}^n(\mathscr{F},\mathscr{G})
    =\prod_{\lambda\in L}
    \operatorname{Diff}_{X/S}^n(\mathscr{F}_\lambda,\mathscr{G}),
  \]
  and consequently also (0\textsubscript{I}, 3.2.6)
  \[
    \mathscr{D}\!iff_{X/S}^n(\mathscr{F},\mathscr{G})
    =\prod_{\lambda\in L}
    \mathscr{D}\!iff_{X/S}^n(\mathscr{F}_\lambda,\mathscr{G}).
  \]
  On the other hand, if $\mathscr{G}=\prod_{\mu\in M}\mathscr{G}_\mu$
  (0\textsubscript{I}, 3.2.6), we have
  \[
    \operatorname{Hom}_{\mathcal{O}_X}
    (\mathscr{P}_{X/S}^n(\mathscr{F}),\mathscr{G})
    =\prod_{\mu\in M}\operatorname{Hom}_{\mathcal{O}_X}
    (\mathscr{P}_{X/S}^n(\mathscr{F}),\mathscr{G}_\mu),
  \]
  \oldpage[IV-4]{46}
  each homomorphism $u$ from $\mathscr{P}_{X/S}^n(\mathscr{F})$ to
  $\mathscr{G}$ corresponding bijectively to the family of its composites
  \[
    u_\mu:\mathscr{P}_{X/S}^n(\mathscr{F})
    \longrightarrow\mathscr{G}\longrightarrow\mathscr{G}_\mu.
  \]
  We therefore have
  \[
    \operatorname{Diff}_{X/S}^n(\mathscr{F},\mathscr{G})
    =\prod_{\mu\in M}
    \operatorname{Diff}_{X/S}^n(\mathscr{F},\mathscr{G}_\mu),
  \]
  and consequently also
  \[
    \mathscr{D}\!iff_{X/S}^n(\mathscr{F},\mathscr{G})
    =\prod_{\mu\in M}
    \mathscr{D}\!iff_{X/S}^n(\mathscr{F},\mathscr{G}_\mu).
  \]
  \item[(ii)] Until now, differential operators
  $\mathscr{F}\to\mathscr{G}$ have hardly been encountered except when
  $\mathscr{F}$ and $\mathscr{G}$ are locally free $\mathcal{O}_X$-modules of
  finite rank, in which case their structure reduces locally, by virtue
  of (i), to that of the sheaf $\mathscr{D}\!iff_{X/S}$; the latter will be
  studied later (16.11) in a particular case.
\end{enumerate}
\end{remarks}


\subsection{Regular and quasi-regular immersions.}
\EGAReaderTarget{ega4.c04.subsec.16-9}{IV.16.9}

\begin{definition}[16.9.1]\label{IV.16.9.1}
\EGAReaderTarget{ega4.c04.def.16-9-1}{IV.16.9.1}
Let $X$ be a ringed space. We say that an ideal $\mathscr{I}$ of
$\mathcal{O}_X$ is \emph{regular} (resp. \emph{quasi-regular}) if, for
every point $x\in\operatorname{Supp}(\mathcal{O}_X/\mathscr{I})$, there exist
an open neighbourhood $U$ of $x$ in $X$ and a regular (0, 15.2.2) (resp.
quasi-regular (0, 15.2.2)) sequence of elements of $\Gamma(U,\mathcal{O}_X)$
which generates $\mathscr{I}|U$.
\end{definition}

We will say that a regular (resp. quasi-regular) sequence of sections of
$\mathcal{O}_X$ over $U$ which generates $\mathscr{I}|U$ is a
\emph{regular} (resp. \emph{quasi-regular}) \emph{system of generators}
of $\mathscr{I}|U$.

\begin{definition}[16.9.2]\label{IV.16.9.2}
\EGAReaderTarget{ega4.c04.def.16-9-2}{IV.16.9.2}
Let $j:Y\to X$ be an immersion of preschemes and let $U$ be an open subset of
$X$ such that $j(Y)\subset U$ and $j$ is a closed immersion of $Y$ into $U$.
We say that $j$ is \emph{regular} (resp. \emph{quasi-regular}) if the
closed sub-prescheme $j(Y)$ of $U$ associated to $j$ is defined by a regular
(resp. quasi-regular) ideal of $\mathcal{O}_U$ (a condition independent of
the open subset $U$ chosen).
\end{definition}

We say that a sub-prescheme $Y$ of a prescheme $X$ is \emph{regularly
immersed} (resp. \emph{quasi-regularly immersed}) if the canonical injection
$j:Y\to X$ is a regular (resp. quasi-regular) immersion. If $Y$ is a
closed sub-prescheme of $X$ and $\mathscr{I}$ the ideal of $\mathcal{O}_X$
that defines $Y$, it amounts to the same to say that $\mathscr{I}$ is
\emph{regular} (resp. \emph{quasi-regular}).

For example, if $A$ is an \emph{integral} ring and $f$ an element $\neq0$ of
$A$, the closed sub-prescheme $V(f)$ of $\operatorname{Spec}(A)$ (isomorphic to
$\operatorname{Spec}(A/fA)$) is \emph{regularly immersed} in
$\operatorname{Spec}(A)$.

Every regular ideal is quasi-regular (0, 15.2.2); every regular immersion
is quasi-regular (cf. (16.9.11) for a converse).

\begin{proposition}[16.9.3]
\label{IV.16.9.3}
Let $X$ be a ringed space, $\mathscr{I}$ an ideal of $\mathcal{O}_X$,
$(f_i)_{1\leq i\leq m}$ a finite sequence of sections of $\mathcal{O}_X$
over $X$ generating $\mathscr{I}$. For $(f_i)$ to be a quasi-regular
sequence (0, 15.2.2), it is necessary and sufficient that the following
conditions be satisfied:
\begin{enumerate}
  \item[(i)] The canonical images of the $f_i$ in
  $\mathscr{I}/\mathscr{I}^2$ form a basis of this
  $\mathcal{O}_X/\mathscr{I}$-module.
  \item[(ii)] The canonical surjective homomorphism (16.1.2.2)
  \[
    \mathbf{S}_{\mathcal{O}_X/\mathscr{I}}^{\bullet}
    (\mathscr{I}/\mathscr{I}^2)
    \longrightarrow\mathscr{G}r_{\mathscr{I}}^{\bullet}(\mathcal{O}_X)
  \]
  is bijective.
\end{enumerate}
Moreover, if this is so, every sequence $(f'_i)_{1\leq i\leq n}$ of $n$
sections of $\mathscr{I}$ over $X$ which generates $\mathscr{I}$ is
quasi-regular.
\end{proposition}

\oldpage[IV-4]{47}

The two conditions of the statement merely translate the definition given in
(0, 15.2.2), taking into account the definition of the canonical
homomorphisms (0, 15.2.1.1). The last assertion follows from the fact that,
if a module $M$ over a commutative ring $A$ has a basis of $n$ elements,
every system of generators of $M$ having $n$ elements is a basis of $M$
(Bourbaki, \emph{Alg. comm.}, chap.~II, §~3, cor.~5 of th.~1).

\begin{corollary}[16.9.4]
\label{IV.16.9.4}
Let $X$ be a ringed space in local rings, $\mathscr{I}$ an ideal of
$\mathcal{O}_X$. For $\mathscr{I}$ to be quasi-regular, it is necessary and
sufficient that it satisfy the following conditions:
\begin{enumerate}
  \item[(i)] $\mathscr{I}$ is of finite type.
  \item[(ii)] $\mathscr{I}/\mathscr{I}^2$ is a locally free
  $(\mathcal{O}_X/\mathscr{I})$-module.
  \item[(iii)] The canonical homomorphism
  \[
    \mathbf{S}_{\mathcal{O}_X/\mathscr{I}}^{\bullet}
    (\mathscr{I}/\mathscr{I}^2)
    \longrightarrow\mathscr{G}r_{\mathscr{I}}^{\bullet}(\mathcal{O}_X)
    \tag{16.9.4.1}
    \label{IV.16.9.4.1}
  \]
  is bijective.
\end{enumerate}
\end{corollary}

The necessity of the conditions follows immediately from (16.9.3). To see that
these conditions are sufficient, we must show, by virtue of (16.9.3), that if,
at a point $x\in\operatorname{Supp}(\mathcal{O}_X/\mathscr{I})$, there exist
an open neighbourhood $U$ of $x$ in $X$ and $n$ sections $f_i$
($1\leq i\leq n$) of $\mathscr{I}$ over $U$ whose canonical images in
$\mathscr{I}/\mathscr{I}^2$ form a basis of
$(\mathscr{I}/\mathscr{I}^2)|U$ over
$(\mathcal{O}_X/\mathscr{I})|U$, then there is an open neighbourhood
$V\subset U$ of $x$ such that the $f_i|V$ generate $\mathscr{I}|V$. Now, by
hypothesis, we have $\mathscr{I}_x\neq\mathcal{O}_x$, hence $\mathscr{I}_x$
is contained in the maximal ideal of $\mathcal{O}_x$; since $\mathscr{I}_x$
is an $\mathcal{O}_x$-module of finite type and the classes of the $(f_i)_x$
in $\mathscr{I}_x/\mathscr{I}_x^2$ generate this
$(\mathcal{O}_x/\mathscr{I}_x)$-module, Nakayama's lemma shows that the
$(f_i)_x$ generate $\mathscr{I}_x$. Since $\mathscr{I}$ is of finite type,
we conclude by (0\textsubscript{I}, 5.2.2).

\begin{corollary}[16.9.5]
\label{IV.16.9.5}
Let $X$ be a ringed space in local rings, $\mathscr{I}$ a quasi-regular
ideal of $\mathcal{O}_X$, $(f_i)_{1\leq i\leq n}$ a sequence of
sections of $\mathscr{I}$ over $X$,
$x$ a point of $\operatorname{Supp}(\mathcal{O}_X/\mathscr{I})$. The
following conditions are equivalent:
\begin{enumerate}
  \item[(a)] There exists an open neighbourhood $U$ of $x$ in $X$ such that
  the $f_i|U$ form a quasi-regular sequence of elements of
  $\Gamma(U,\mathcal{O}_X)$ generating $\mathscr{I}|U$.
  \item[(b)] The $(f_i)_x$ form a system of generators of
  $\mathscr{I}_x$ whose number of elements is the smallest possible.
  \item[(b')] The $(f_i)_x$ form a minimal system of generators of
  $\mathscr{I}_x$.
  \item[(c)] If $\bar f_i$ is the canonical image of $f_i$ in
  $\Gamma(X,\mathscr{I}/\mathscr{I}^2)$, the $(\bar f_i)_x$ form a basis
  of the $(\mathcal{O}_x/\mathscr{I}_x)$-module
  $\mathscr{I}_x/\mathscr{I}_x^2$.
\end{enumerate}
\end{corollary}

By hypothesis, $\mathcal{O}_x$ is a local ring and $\mathscr{I}_x$ an ideal
of finite type of $\mathcal{O}_x$ contained in the maximal ideal of
$\mathcal{O}_x$; the equivalence of b), b') and c) therefore follows from
Nakayama's lemma (Bourbaki, \emph{Alg. comm.}, chap.~II, §~3,
n\textsuperscript{o}~2, prop.~5). It is clear that a) implies c) by virtue
of (16.9.3); on the other hand, it follows from (0\textsubscript{I}, 5.2.2)
that if condition c) is satisfied (hence also b)), there is a neighbourhood
$U$ of $x$ in $X$ such that $(\mathscr{I}/\mathscr{I}^2)|U$ has constant rank
equal to $n$ and the $f_i|U$ generate $\mathscr{I}|U$; it therefore suffices
to apply in $U$ the last assertion of (16.9.3).

\begin{remarks}[16.9.6]
\label{IV.16.9.6}
\begin{enumerate}
  \item[(i)] Under the general hypotheses of (16.9.5), it is not sufficient
  that the $(\bar f_i)_y$ form a basis of the
  $(\mathcal{O}_y/\mathscr{I}_y)$-module
  $\mathscr{I}_y/\mathscr{I}_y^2$ for every $y\in X$ in order that the
  sequence $(f_i)$ generate $\mathscr{I}$. An example of this is obtained by
\oldpage[IV-4]{48}
  taking $X=\operatorname{Spec}(A)$, where $A$ is a Dedekind ring, and
  $\mathscr{I}=\widetilde{\mathfrak{I}}$, where $\mathfrak{I}$ is a
  non-principal prime ideal of $A$; indeed, we then have
  $\mathscr{I}_y/\mathscr{I}_y^2=0$ at every point $y$ distinct from the point
  $x\in X$ corresponding to $\mathfrak{I}$, and
  $\mathscr{I}_x/\mathscr{I}_x^2$ is of rank $1$ over the field
  $\mathcal{O}_x/\mathscr{I}_x$; moreover, $\mathscr{I}$ is evidently a
  regular ideal.
  \item[(ii)] In (16.9.5), ``quasi-regular'' cannot be replaced by
  ``regular'', even when $X$ is a prescheme (cf. (16.9.12)).
  Indeed, let $B$ denote the ring of germs of infinitely
  differentiable functions at the point $0$ in $\mathbb{R}$; it has a maximal
  ideal $\mathfrak{m}$ generated by the germ $t$ of the identity map of
  $\mathbb{R}$ at the point $0$, and the intersection $\mathfrak{n}$ of the
  $\mathfrak{m}^k$ for $k>0$ is not reduced to $0$. Now let $A$ be
  the quotient ring $B[T]/\mathfrak{n}TB[T]$, and let $f_1,f_2$ be the
  canonical images in $A$ of the elements $t$ and $T$ of $B[T]$. The sequence
  $(f_1,f_2)$ is regular in $A$: indeed, $f_1$ is not a divisor of
  $0$ in $A$, because the relation $tP[T]\in\mathfrak{n}TB[T]$, for a
  polynomial $P\in B[T]$, implies that the products of $t$ by the coefficients
  of $P$ belong to the ideal $\mathfrak{n}$, and it follows immediately that
  these coefficients are themselves in $\mathfrak{n}$, hence
  $P[T]\in\mathfrak{n}TB[T]$. Since $B/tB$ is isomorphic to $\mathbb{R}$,
  $A/f_1A$ is isomorphic to the polynomial ring $\mathbb{R}[T]$, hence
  integral, and the image of $f_2$ in $A/f_1A$, being equal to $T$, is
  therefore not a divisor of $0$, which proves our assertion. However, $f_2$
  is a divisor of $0$ in $A$, because for every nonzero element
  $x\in\mathfrak{n}$, the image of $x$ in $A$ is $\ne0$, but the image of $xT$
  is zero. We conclude that the sequence $(f_2,f_1)$ is not regular in $A$;
  on the other hand, the ideal $\mathfrak{I}=f_1A+f_2A$ is distinct from $A$,
  hence conditions b), b') and c) of (16.9.5) do not imply condition a)
  when ``quasi-regular'' is replaced in it by ``regular''.
\end{enumerate}
\end{remarks}

\begin{env}[16.9.7]\label{IV.16.9.7}
\EGAReaderTarget{ega4.c04.statement.16-9-7}{IV.16.9.7}
If $X=\operatorname{Spec}(A)$ is an affine scheme, we will say that an ideal
$\mathfrak{I}$ of $A$ is regular (resp. quasi-regular) if the ideal
$\mathscr{I}=\widetilde{\mathfrak{I}}$ of $\mathcal{O}_X$ is regular
(resp. quasi-regular); note that this notion is local and in no way implies
the existence of a system of generators of $\mathfrak{I}$ forming a regular
(resp. quasi-regular) sequence in $A$, as the example (16.9.6, (i)) shows;
this is however the case when $A$ is local (16.9.5).

Proposition (16.9.4) translates in terms of quasi-regular immersions
as follows:
\end{env}

\begin{proposition}[16.9.8]
\label{IV.16.9.8}
Let $j:Y\to X$ be a morphism of preschemes; for $j$ to be a quasi-regular
immersion, it is necessary and sufficient that $j$ satisfy the following
conditions:
\begin{enumerate}
  \item[(i)] $j$ is an immersion locally of finite presentation.
  \item[(ii)] The conormal sheaf
  $\operatorname{Gr}^1(j)=\mathcal{N}_{Y/X}$ (16.1.2) is a locally free
  $\mathcal{O}_Y$-module.
  \item[(iii)] The canonical homomorphism
  \[
    \operatorname{S}_{\mathcal{O}_Y}^{\bullet}
      (\operatorname{Gr}^1(j))\longrightarrow\operatorname{Gr}^{\bullet}(j)
  \]
  (16.1.2.2) is bijective.
\end{enumerate}
\end{proposition}

The question being local on $Y$, we may restrict to the case where $j$ is
the canonical injection of a closed sub-prescheme $Y$ of $X$, and then the
translation of (16.9.4) into (16.9.8) follows from the explicit description
of $\operatorname{Gr}^1(j)$ and $\operatorname{Gr}^{\bullet}(j)$ in terms of
the ideal $\mathscr{I}$ of $\mathcal{O}_X$ defining the sub-prescheme $Y$
(16.1.3, (ii)).

\begin{corollary}[16.9.9]
\label{IV.16.9.9}
Let $Y$ be a prescheme, $X$ a $Y$-prescheme, $j:Y\to X$ a $Y$-section
of $X$, so that the $n$-th normal invariant $\mathscr{A}^{(n)}$ of $j$
(16.1.2) is an augmented $\mathcal{O}_Y$-algebra (16.1.7); set
$\mathscr{A}^{(\infty)}=\varprojlim\mathscr{A}^{(n)}$. For $j$ to be a
quasi-regular immersion, it is necessary and sufficient that $j$ be locally
of finite presentation and that every $y\in Y$ have an affine open
neighbourhood $U$ with ring $C$ such that $\mathscr{A}^{(\infty)}|U$ is
isomorphic, as an augmented topological $\mathcal{O}_U$-algebra, to
$\mathcal{O}_U[[T_1,\ldots,T_n]]$.
\end{corollary}

We may restrict to the case where $j$ is a closed immersion by restricting
to a sufficiently small neighbourhood of $y$ (see the argument of (16.4.11)),
and then $\mathcal{O}_Y$ is identified with a quotient algebra
$\mathcal{O}_X/\mathscr{I}$ and the canonical surjective homomorphism
$\mathcal{O}_X\to\mathcal{O}_Y$ has a
\oldpage[IV-4]{49}
right inverse (16.1.7). We may therefore suppose
$X=\operatorname{Spec}(B)$ and $Y=\operatorname{Spec}(A)$ affine, $B$ being
an augmented $A$-algebra and the augmentation ideal $\mathfrak{I}$ being of
finite type. Since $\mathscr{A}^{(n)}$ is then identified with
$(B/\mathfrak{I}^{n+1})^{\sim}$, the corollary follows from the equivalence
of b) and c) in (0, 19.5.4), since $B/\mathfrak{I}=A$.

Note that, in the affine case under consideration, the fact that $j$ is a
quasi-regular immersion is also equivalent, by virtue of (0, 19.5.4), to
saying that $B$ is a formally smooth $A$-algebra for the
$\mathfrak{I}$-preadic topology.

Note also that the condition that $j$ be an immersion locally of finite
presentation is always satisfied when the morphism $X\to Y$ is locally of
finite type (1.4.3, (v)).

\begin{proposition}[16.9.10]\label{IV.16.9.10}
\EGAReaderTarget{ega4.c04.prop.16-9-10}{IV.16.9.10}
Let $X$ be a locally Noetherian prescheme, $Y$ a sub-prescheme of $X$,
$j:Y\to X$ the canonical injection, $y$ a point of $Y$.
\begin{enumerate}
  \item[(i)] For there to exist an open neighbourhood $U$ of $y$ in $X$ such
  that the restriction $Y\cap U\to U$ of $j$ is a regular immersion, it is
  necessary and sufficient that the kernel $\mathscr{I}_y$ of the surjective
  homomorphism $\mathcal{O}_{X,y}\to\mathcal{O}_{Y,y}$ be generated by a
  regular sequence of elements of $\mathcal{O}_{X,y}$.
  \item[(ii)] For the immersion $j$ to be regular, it is necessary and
  sufficient that it be quasi-regular.
\end{enumerate}
\end{proposition}

\begin{enumerate}
  \item[(i)] We may restrict to the case where $Y$ is a closed sub-prescheme
  of $X$ defined by a \emph{coherent} ideal $\mathscr{I}$ of $\mathcal{O}_X$.
  The condition is evidently necessary. Conversely, if $\mathscr{I}_y$
  is generated by a regular sequence $(s_i)_y$, where the $s_i$ are
  sections of $\mathscr{I}$ over an open neighbourhood $U$ of $y$ in
  $X$, we may suppose that the $s_i$ generate $\mathscr{I}|U$
  (0\textsubscript{I}, 5.2.2) and form a regular sequence (0, 15.2.4),
  whence the assertion.
  \item[(ii)] The fact that a quasi-regular immersion is regular follows
  from (i) and from the coincidence of quasi-regular sequences and regular
  sequences in $\mathcal{O}_{X,y}$ formed of elements of the maximal ideal
  (0, 15.1.11).
\end{enumerate}

When (without a Noetherian hypothesis on $X$) the kernel $\mathscr{I}_y$ of
$\mathcal{O}_{X,y}\to\mathcal{O}_{Y,y}$ is generated by a regular sequence
of elements of $\mathcal{O}_{X,y}$, we say that the immersion $j$ is
\emph{regular at the point $y$}.

\begin{corollary}[16.9.11]
\label{IV.16.9.11}
Let $X$ be a locally Noetherian prescheme; then every quasi-regular ideal of
$\mathcal{O}_X$ is regular.
\end{corollary}

\begin{remarks}[16.9.12]
\label{IV.16.9.12}
\begin{enumerate}
  \item[(i)] Note that a regular immersion is not in general a flat
  morphism, nor \emph{a fortiori} a \emph{regular} morphism in the sense of
  (6.8.1).
  \item[(ii)] Let $A$ be a Noetherian local ring; it follows immediately from
  (16.9.4) and (0, 17.1.1) that for $A$ to be \emph{regular}, it is
  necessary and sufficient that its maximal ideal $\mathfrak{m}$ be
  \emph{quasi-regular} (or \emph{regular}, which amounts to the same since $A$
  is Noetherian). For a Noetherian affine scheme $X$ to be \emph{regular}, it
  is necessary and sufficient that for every closed point $x\in X$ the
  canonical injection $\operatorname{Spec}(k(x))\to X$ be a \emph{regular}
  immersion.
\end{enumerate}
\end{remarks}

\begin{proposition}[16.9.13]
\label{IV.16.9.13}
Let $X$ be a locally Noetherian prescheme, $Y$ a sub-prescheme of $X$,
$Y'$ a sub-prescheme of $Y$, such that the canonical injection $j:Y'\to Y$ is
regular. Then the sequence of $\mathcal{O}_{Y'}$-modules
\[
  \xymatrix{
    0\ar[r] & j^*(\mathcal{N}_{Y/X})\ar[r]
      & \mathcal{N}_{Y'/X}\ar[r] & \mathcal{N}_{Y'/Y}\ar[r] & 0
  }
  \tag{16.9.13.1}
  \label{IV.16.9.13.1}
\]
\oldpage[IV-4]{50}
is exact; moreover, for every $x\in X$, there exists an open neighbourhood $U$
of $x$ such that the restrictions to $U$ of the homomorphisms of (16.9.13.1)
form a split exact sequence.
\end{proposition}

Let us first prove the following lemma:

\begin{lemma}[16.9.13.2]
\label{IV.16.9.13.2}
Let $A$ be a ring, $\mathfrak{I}$ an ideal of $A$,
$A'=A/\mathfrak{I}$, $(f_i)_{1\leq i\leq r}$ a sequence of elements of $A$
that is $A'$-regular, $\mathfrak{K}=\sum_i f_iA$,
$\mathfrak{L}=\mathfrak{I}+\mathfrak{K}$,
$\mathfrak{K}'=\sum_i f_iA'$, so that $C=A/\mathfrak{L}$ is isomorphic to
$A'/\mathfrak{K}'$. Then for every integer $n>0$ and every integer $N\geq n$,
we have the relation
\[
  \mathfrak{I}\cap\mathfrak{K}^n
  =\mathfrak{I}\mathfrak{K}^n+(\mathfrak{I}\cap\mathfrak{K}^N).
  \tag{16.9.13.3}
  \label{IV.16.9.13.3}
\]
\end{lemma}

It evidently suffices to prove that every element of the left-hand side is
contained in the right-hand side, and, by induction on $n$, we are reduced
to the case where $N=n+1$. An element of the left-hand side of (16.9.13.3),
being in $\mathfrak{K}^n$, can be written $P(f_1,\ldots,f_r)$, where
$P\in A[T_1,\ldots,T_r]$ is homogeneous of degree $n$. If $f_i'$ is the
canonical image of $f_i$ in $A'$, the hypothesis
$P(f_1,\ldots,f_r)\in\mathfrak{I}$ means that
$P(f_1',\ldots,f_r')=0$. But
$P(f_1',\ldots,f_r')\in\mathfrak{K}'^n$, hence the canonical image of
$P(f_1',\ldots,f_r')$ in
$\mathfrak{K}'^n/\mathfrak{K}'^{n+1}$ is zero. Now the hypothesis that the
sequence $(f_i)$ is $A'$-regular implies that the canonical homomorphism
\[
  \mathbf{S}_C^n(\mathfrak{K}'/\mathfrak{K}'^2)
  \longrightarrow\mathfrak{K}'^n/\mathfrak{K}'^{n+1}
\]
is bijective (0, 15.1.9); we conclude that the coefficients of $P$ belong
to $\mathfrak{L}=\mathfrak{I}+\mathfrak{K}$. It follows immediately that we
have $P(f_1,\ldots,f_r)\in\mathfrak{I}\mathfrak{K}^n+\mathfrak{K}^{n+1}$, and
since $P(f_1,\ldots,f_r)\in\mathfrak{I}$, we finally have
\[
  P(f_1,\ldots,f_r)
  \in\mathfrak{I}\mathfrak{K}^n
    +(\mathfrak{I}\cap\mathfrak{K}^{n+1}),
\]
which proves the lemma.

Taking the quotients of both sides of (16.9.13.3) by
$\mathfrak{I}\mathfrak{K}^n$, we see that the relations (16.9.13.3) for
$N\geq n$ imply
\[
  (\mathfrak{I}\cap\mathfrak{K}^n)/\mathfrak{I}\mathfrak{K}^n
  \subset\bigcap_{N\geq n}\mathfrak{K}^N
    \cdot(A/(\mathfrak{I}\mathfrak{K}^n)).
  \tag{16.9.13.4}
  \label{IV.16.9.13.4}
\]
From this we deduce the following

\begin{corollary}[16.9.13.5]
\label{IV.16.9.13.5}
Suppose that the hypotheses of (16.9.13.2) are satisfied, and suppose
moreover that the ring $A$ is Noetherian and that $\mathfrak{K}$ is contained
in the radical of $A$. Then for every integer $n>0$ we have
\[
  \mathfrak{I}\cap\mathfrak{K}^n=\mathfrak{I}\mathfrak{K}^n.
  \tag{16.9.13.6}
  \label{IV.16.9.13.6}
\]
\end{corollary}

Indeed, the right-hand side of (16.9.13.4) is then zero, since
$A/\mathfrak{I}\mathfrak{K}^n$ is an $A$-module of finite type (Bourbaki,
\emph{Alg. comm.}, chap.~III, §~3, n\textsuperscript{o}~3, prop.~6).

Let us take in particular $n=2$ in (16.9.13.6), and note that we have
$\mathfrak{L}^2=\mathfrak{I}^2+\mathfrak{I}\mathfrak{K}+\mathfrak{K}^2
=\mathfrak{I}\mathfrak{L}+\mathfrak{K}^2$; since
$\mathfrak{I}\mathfrak{L}\subset\mathfrak{L}^2$, we deduce that
\[
  \mathfrak{I}\cap\mathfrak{L}^2
  =\mathfrak{I}\mathfrak{L}+(\mathfrak{I}\cap\mathfrak{K}^2)
  =\mathfrak{I}\mathfrak{L}+\mathfrak{I}\mathfrak{K}^2
  =\mathfrak{I}\mathfrak{L},
\]
in other words
\[
  \mathfrak{I}\cap\mathfrak{L}^2=\mathfrak{I}\mathfrak{L},
  \tag{16.9.13.7}
  \label{IV.16.9.13.7}
\]
which can also be expressed by saying that the canonical homomorphism
\[
  \mathfrak{I}/\mathfrak{I}\mathfrak{L}
  \longrightarrow(\mathfrak{I}+\mathfrak{L}^2)/\mathfrak{L}^2
\]
is bijective.

These lemmas having been proved, let us prove the first assertion of
(16.9.13): it evidently suffices
\oldpage[IV-4]{51}
to prove that the sequence of stalks of the sheaves appearing in
(16.9.13.1) at a point $x\in Y'$ is exact. Now, if we set
$A=\mathcal{O}_{X,x}$, we can write
$\mathcal{O}_{Y,x}=A'=A/\mathfrak{I}$, where $\mathfrak{I}$ is an ideal
contained in the maximal ideal of $A$, and then
$\mathcal{O}_{Y',x}=A'/\mathfrak{K}'$, where $\mathfrak{K}'$ is generated by
an $A'$-regular sequence of elements of $A'$, which are themselves the images
of the elements of an $A'$-regular sequence of elements of $A$ belonging to
the maximal ideal of $A$. If $\mathfrak{K}$ is the ideal generated by the
latter and $\mathfrak{L}=\mathfrak{I}+\mathfrak{K}$, we have
$\mathcal{O}_{Y',x}=A/\mathfrak{L}$, and since we are in the situation of
(16.9.13.5), the canonical homomorphism
\[
  \mathfrak{I}/\mathfrak{I}\mathfrak{L}
  \longrightarrow(\mathfrak{I}+\mathfrak{L}^2)/\mathfrak{L}^2
\]
is bijective. But this shows that the sequence
\[
  0\longrightarrow\mathfrak{I}/\mathfrak{I}\mathfrak{L}
  \longrightarrow\mathfrak{L}/\mathfrak{L}^2
  \longrightarrow(\mathfrak{L}/\mathfrak{I})/
    (\mathfrak{L}/\mathfrak{I})^2\longrightarrow0
\]
is exact (see the proof of (16.2.7)), and the modules appearing in this
sequence are precisely the stalks at $x$ of the sheaves of (16.9.13.1).
The second assertion follows from the fact that $\mathcal{N}_{Y'/Y}$ is a
locally \emph{free} $\mathcal{O}_{Y'}$-module (16.9.8) and from Bourbaki,
\emph{Alg.}, chap.~II, 3\textsuperscript{e} éd., §~1, n\textsuperscript{o}~11,
prop.~21.

\subsection{Differentially smooth morphisms}
\EGAReaderTarget{ega4.c04.subsec.16-10}{IV.16.10}

\begin{definition}[16.10.1]\label{IV.16.10.1}
\EGAReaderTarget{ega4.c04.def.16-10-1}{IV.16.10.1}
We say that a morphism of preschemes $f:X\to S$ is
\emph{differentially smooth} (or that $X$ is \emph{differentially
smooth over $S$}) if it satisfies the following conditions:
\begin{enumerate}
  \item[(i)] $\Omega_{X/S}^1$ is a locally projective
  $\mathcal{O}_X$-module, that is, every point of $X$ has an affine open
  neighbourhood $U$ such that $\Gamma(U,\Omega_{X/S}^1)$ is a projective
  $\Gamma(U,\mathcal{O}_X)$-module (not necessarily of finite
  type).
  \item[(ii)] The canonical homomorphism (16.3.1.1)
  \[
    \mathbf{S}_{\mathcal{O}_X}^{\bullet}(\Omega_{X/S}^1)
    \longrightarrow\mathscr{G}r_{\bullet}(\mathscr{P}_{X/S})
  \]
  is bijective.
\end{enumerate}
In particular, if $\Omega_{X/S}^1$ is locally free of finite rank, the
$\mathscr{P}_{X/S}^n$ are locally free $\mathcal{O}_X$-modules of finite
rank (being extensions of such modules).
\end{definition}

We say that $f$ is \emph{differentially smooth at a point $x\in X$} (or
that $X$ is differentially smooth over $S$ at the point $x$) if there exists an
open neighbourhood $U$ of $x$ in $X$ such that $f|U$ is differentially
smooth.

We will see later (17.12.4) that a smooth morphism is differentially
smooth, which justifies the terminology; but the converse is false; indeed,
a \emph{monomorphism} $f:X\to S$ is differentially smooth,
since $\Omega_{X/S}^1=0$ by virtue of (I, 5.3.8), and consequently the
surjective homomorphism (16.3.1.1) is evidently bijective; but a
monomorphism is not even necessarily flat, nor \emph{a fortiori} smooth.
Let us restrict ourselves here to noting the following proposition:

\begin{proposition}[16.10.2]
\label{IV.16.10.2}
Let $A$ be a ring, $B$ an $A$-algebra that is formally smooth for the
discrete topologies (0, 19.3.1). Then $\operatorname{Spec}(B)$ is
differentially smooth over $\operatorname{Spec}(A)$.
\end{proposition}

Indeed, $B\otimes_A B$ is then (for the discrete topologies) a formally
smooth $B$-algebra (for either of the canonical
homomorphisms $b\mapsto b\otimes1$, $b\mapsto1\otimes b$ from $B$
\oldpage[IV-4]{52}
to $B\otimes_A B$) (0, 19.3.5, (iii)); hence $B\otimes_A B$ is also a
formally smooth $A$-algebra for the discrete topologies (0, 19.3.5,
(ii)). Setting $\mathfrak{I}=\mathfrak{I}_{B/A}$, it follows that
$B\otimes_A B$ is also a formally smooth $A$-algebra for the
$\mathfrak{I}$-preadic topology (0, 19.3.8); since by hypothesis
$B=(B\otimes_A B)/\mathfrak{I}$ is a formally smooth $A$-algebra for
the discrete topologies, the proposition follows from the equivalence of a)
and b) in (0, 19.5.4).

\begin{proposition}[16.10.3]
\label{IV.16.10.3}
For a morphism $f:X\to S$ to be differentially smooth, it is necessary and
sufficient that for every $x\in X$ there exist an affine open neighbourhood
$U$ of $x$, with ring $A$, such that $\Gamma(U,\mathscr{P}_{X/S}^{\infty})$ is
an augmented topological $A$-algebra isomorphic to the completed algebra
$\widehat B$, where $B=\mathbf{S}_A(V)$, $V$ being a projective $A$-module and
$B$ being endowed with the $B^+$-preadic topology (where $B^+$ is the
augmentation ideal). If $\Omega_{X/S}^1$ is locally free of finite rank,
$\widehat B$ can be replaced by the formal power series algebra
$A[[T_1,\ldots,T_n]]$.
\end{proposition}

The notion of differentially smooth morphism being evidently local on
$X$, we may restrict to the case where $S=\operatorname{Spec}(B)$,
$X=\operatorname{Spec}(C)$. Let us regard $C\otimes_B C$ as a $C$-algebra
(via the first factor); set $\mathfrak{I}=\mathfrak{I}_{C/B}$ and
endow $C\otimes_B C$ with the $\mathfrak{I}$-preadic topology; we can
apply to the topological $C$-algebra $C\otimes_B C$ and to the ideal
$\mathfrak{I}$ of $C\otimes_B C$ the equivalence of b) and c) in
(0, 19.5.4), since $(C\otimes_B C)/\mathfrak{I}=C$ is evidently a formally
smooth $C$-algebra for the discrete topologies. The topology on
$\Gamma(U,\mathscr{P}_{X/S}^{\infty})$ is evidently the projective limit
topology of this ring (16.1.11).

Note that the integer $n$ in the statement of (16.10.3) is the \emph{rank}
of $\Omega_{X/S}^1$ at the point $x$. We will see later (17.13.5) that when
$f$ is differentially smooth and locally of finite type, $n$ is also
equal to the dimension of the fibre $f^{-1}(f(x))$ at the point $x$.

\begin{proposition}[16.10.4]
\label{IV.16.10.4}
Let $f:X\to S$, $g:S'\to S$ be two morphisms, and set
$X'=X\times_S S'$, $f'=f_{(S')}:X'\to S'$.
\begin{enumerate}
  \item[(i)] If $f$ is differentially smooth, so is $f'$.
  \item[(ii)] Conversely, if $g$ is faithfully flat and quasi-compact, and if
  $f'$ is differentially smooth and $\Omega_{X'/S'}^1$ is an
  $\mathcal{O}_{X'}$-module of finite type, then $f$ is differentially smooth
  and $\Omega_{X/S}^1$ is an $\mathcal{O}_X$-module of finite type.
\end{enumerate}
\end{proposition}

Indeed, if $f$ is differentially smooth, the
$\mathscr{G}r_n(\mathscr{P}_{X/S})$ are \emph{flat}
$\mathcal{O}_X$-modules; consequently (16.4.6), the homomorphism
\[
  \mathscr{G}r_n(\mathscr{P}_{X/S})\otimes_{\mathcal{O}_X}\mathcal{O}_{X'}
  \longrightarrow\mathscr{G}r_n(\mathscr{P}_{X'/S'})
\]
is bijective for every $n$, and by virtue of the commutativity of the diagram
(16.2.1.3), it follows from the definition (16.10.1) that $f'$ is
differentially smooth. On the other hand, if $g$ is faithfully flat and
quasi-compact, it again follows from (16.4.6) that
\[
  \mathscr{G}r_n(\mathscr{P}_{X/S})\otimes_{\mathcal{O}_X}\mathcal{O}_{X'}
  \longrightarrow\mathscr{G}r_n(\mathscr{P}_{X'/S'})
\]
is bijective for every $n$. Suppose then that $f'$ is differentially smooth
and $\Omega_{X'/S'}^1$ is of finite rank. Since the canonical projection
$X'\to X$ is a faithfully flat and quasi-compact morphism, it follows first
from (2.5.2) that $\Omega_{X/S}^1$ is a locally free $\mathcal{O}_X$-module of
finite rank, and then from (2.2.7) that the canonical homomorphism (16.3.1.1)
is bijective, hence $f$ is differentially smooth.

\begin{proposition}[16.10.5]
\label{IV.16.10.5}
For a morphism $f:X\to S$ locally of finite type to be
differentially smooth, it is necessary and sufficient that the diagonal
immersion $\Delta_f:X\to X\times_S X$ be quasi-regular.
\end{proposition}

\oldpage[IV-4]{53}

The question being local, we may restrict to the case where $S$ and $X$ are
affine, and consequently the diagonal sub-prescheme of $X\times_S X$ is
closed. The hypothesis that $f$ is locally of finite type implies that
$\Delta_f$ is locally of finite presentation (I, 4.3.1), hence the diagonal
sub-prescheme of $X\times_S X$ is defined by an ideal of finite type
$\mathscr{I}$, and $\Omega_{X/S}^1=\mathscr{I}/\mathscr{I}^2$ is an
$\mathcal{O}_X$-module of finite type. The proposition then follows
immediately from comparing the conditions in (16.10.1) and (16.9.4).

\begin{remark}[16.10.6]
\label{IV.16.10.6}
Let $f:X\to S$ be a morphism such that the $\mathcal{O}_X$-module
$\Omega_{X/S}^1$ is locally free of finite rank. It then follows from
(0, 20.4.7) that every $x\in X$ has an open neighbourhood $U$ such that there
exists a finite family $(z_\lambda)_{\lambda\in L}$ of sections of
$\mathcal{O}_X$ over $U$ for which $(dz_\lambda)_{\lambda\in L}$
forms a \emph{basis} of the $\Gamma(U,\mathcal{O}_X)$-module
$\Gamma(U,\Omega_{X/S}^1)$.
\end{remark}

\subsection{\texorpdfstring{Differential operators on a differentially smooth
$S$-prescheme}{Differential operators on a differentially smooth S-prescheme}}
\EGAReaderTarget{ega4.c04.subsec.16-11}{IV.16.11}

\begin{env}[16.11.1]
\label{IV.16.11.1}
Let $f:X\to S$ be a morphism, $U$ an open subset of $X$,
$(z_\lambda)_{\lambda\in L}$ a family of sections of $\mathcal{O}_X$
over $U$ such that the $dz_\lambda$ form a system of generators
of $\Omega_{X/S}^1|U=\Omega_{U/S}^1$. Let $m$ be an integer or the symbol
$\infty$, and set, for every $\lambda$,
\[
  \zeta_\lambda=\delta z_\lambda=d^m z_\lambda-z_\lambda
  \in\Gamma(U,\mathscr{P}_{X/S}^m).
  \tag{16.11.1.1}
  \label{IV.16.11.1.1}
\]

We will moreover use the usual notation of analysis; for every
$\mathbf{p}=(p_\lambda)\in\mathbb{N}^{(L)}$ (with $p_\lambda=0$ except for
finitely many indices), we set
\[
  |\mathbf{p}|=\sum_\lambda p_\lambda,
  \qquad \mathbf{p}!=\prod_\lambda(p_\lambda!)
  \tag{16.11.1.2}
  \label{IV.16.11.1.2}
\]
\[
  \binom{\mathbf{p}}{\mathbf{q}}
  =\mathbf{p}!/(\mathbf{q}!(\mathbf{p}-\mathbf{q})!)
  \quad\text{for $\mathbf{p},\mathbf{q}$ in $\mathbb{N}^{(L)}$,
  $\mathbf{q}\leq\mathbf{p}$}
  \tag{16.11.1.3}
  \label{IV.16.11.1.3}
\]
and we adopt the convention that $\binom{\mathbf{p}}{\mathbf{q}}=0$ if
$\mathbf{q}\nleq\mathbf{p}$,
\[
  \mathbf{z}^{\mathbf{p}}=\prod_\lambda(z_\lambda)^{p_\lambda},
  \qquad
  \boldsymbol{\zeta}^{\mathbf{p}}=\prod_\lambda(\zeta_\lambda)^{p_\lambda}.
  \tag{16.11.1.4}
  \label{IV.16.11.1.4}
\]

We will therefore have, with this notation,
\[
  d^m(\mathbf{z}^{\mathbf{p}})
  =(d^m(\mathbf{z}))^{\mathbf{p}}
  =(\boldsymbol{\zeta}+\mathbf{z})^{\mathbf{p}}
  =\sum_{\mathbf{q}\leq\mathbf{p}}
    \binom{\mathbf{p}}{\mathbf{q}}
    \mathbf{z}^{\mathbf{p}-\mathbf{q}}\boldsymbol{\zeta}^{\mathbf{q}}
  \tag{16.11.1.5}
  \label{IV.16.11.1.5}
\]
\[
  \boldsymbol{\zeta}^{\mathbf{p}}
  =(d^m\mathbf{z}-\mathbf{z})^{\mathbf{p}}
  =\sum_{\mathbf{q}\leq\mathbf{p}}
    (-1)^{|\mathbf{p}-\mathbf{q}|}
    \binom{\mathbf{p}}{\mathbf{q}}
    \mathbf{z}^{\mathbf{p}-\mathbf{q}}d^m(\mathbf{z}^{\mathbf{q}}).
  \tag{16.11.1.6}
  \label{IV.16.11.1.6}
\]

Since the $dz_\lambda$ generate $\Omega_{X/S}^1$ and are the images of the
$\delta z_\lambda$, and since the canonical homomorphism (16.3.1.1) is
surjective, we conclude that for finite $m$, the $\delta z_\lambda$ generate
the $\mathcal{O}_U$-algebra $\mathscr{P}_{U/S}^m$ (Bourbaki,
\emph{Alg. comm.}, chap.~III, §~2, n\textsuperscript{o}~8, cor.~2 of th.~1).
Hence the $\boldsymbol{\zeta}^{\mathbf{p}}$ (for $|\mathbf{p}|\leq m$)
generate the $\mathcal{O}_U$-module $\mathscr{P}_{U/S}^m$. A differential
operator $D\in\operatorname{Diff}_{U/S}^m$ is consequently entirely
determined by the values of the
$\langle\boldsymbol{\zeta}^{\mathbf{p}},D\rangle$ for
$|\mathbf{p}|\leq m$, or, what amounts to the same by (16.11.1.5) and
(16.11.1.6), by the values
\oldpage[IV-4]{54}
of the $\langle d^m(\mathbf{z}^{\mathbf{p}}),D\rangle=D(\mathbf{z}^{\mathbf{p}})$
for $|\mathbf{p}|\leq m$; precisely, it follows from (16.11.1.5) that
we have
\[
  D(\mathbf{z}^{\mathbf{p}})
  =\langle d^m(\mathbf{z}^{\mathbf{p}}),D\rangle
  =\sum_{\mathbf{q}\leq\mathbf{p}}
    \binom{\mathbf{p}}{\mathbf{q}}
    \langle\boldsymbol{\zeta}^{\mathbf{q}},D\rangle
    \mathbf{z}^{\mathbf{p}-\mathbf{q}}.
  \tag{16.11.1.7}
  \label{IV.16.11.1.7}
\]
\end{env}

\begin{theorem}[16.11.2]
\label{IV.16.11.2}
Let $f:X\to S$ be a morphism, $U$ an open subset of $X$,
$(z_\lambda)_{\lambda\in L}$ a family of sections of $\mathcal{O}_X$
over $U$ such that the family $(dz_\lambda)_{\lambda\in L}$ generates
$\Omega_{X/S}^1|U=\Omega_{U/S}^1$. The following conditions are
equivalent:
\begin{enumerate}
  \item[a)] $f|U$ is differentially smooth and $(dz_\lambda)$ is a basis
  of the $\mathcal{O}_U$-module $\Omega_{U/S}^1$.
  \item[b)] There exists a family
  $(D_{\mathbf{p}})_{\mathbf{p}\in\mathbb{N}^{(L)}}$ of differential
  operators from $\mathcal{O}_U$ to itself satisfying the conditions
  \[
    D_{\mathbf{p}}(\mathbf{z}^{\mathbf{q}})
    =\binom{\mathbf{q}}{\mathbf{p}}
      \mathbf{z}^{\mathbf{q}-\mathbf{p}}
    \quad(\mathbf{p},\mathbf{q}\text{ in }\mathbb{N}^{(L)}).
    \tag{16.11.2.1}
    \label{IV.16.11.2.1}
  \]
\end{enumerate}
Moreover, when these conditions are satisfied, the family
$(D_{\mathbf{p}})$ is uniquely determined by the conditions
(16.11.2.1) and satisfies the relations
\[
  D_{\mathbf{p}}\circ D_{\mathbf{q}}
  =D_{\mathbf{q}}\circ D_{\mathbf{p}}
  =\frac{(\mathbf{p}+\mathbf{q})!}{\mathbf{p}!\mathbf{q}!}
    D_{\mathbf{p}+\mathbf{q}}
  \quad(\mathbf{p},\mathbf{q}\text{ in }\mathbb{N}^{(L)}).
  \tag{16.11.2.2}
  \label{IV.16.11.2.2}
\]
Finally, if $L$ is finite, then for every integer $m$, the $D_{\mathbf{p}}$
such that $|\mathbf{p}|\leq m$ form a basis of the $\mathcal{O}_U$-module
$\operatorname{Diff}_{U/S}^m$; in other words, every differential operator
of order $\leq m$ on $U$ can be written in exactly one way in the form
\[
  D=\sum_{|\mathbf{p}|\leq m}a_{\mathbf{p}}D_{\mathbf{p}}
\]
where the $a_{\mathbf{p}}$ are sections of $\mathcal{O}_X$ over
$U$.
\end{theorem}

Let us first note that, by virtue of (16.11.1.6) and (16.11.1.5), we verify
immediately that the conditions (16.11.2.1) are equivalent to
\[
  \langle\boldsymbol{\zeta}^{\mathbf{p}},D_{\mathbf{q}}\rangle
  =\delta_{\mathbf{p}\mathbf{q}}
  \quad\text{(Kronecker delta)}.
  \tag{16.11.2.3}
  \label{IV.16.11.2.3}
\]
The existence of the family $(D_{\mathbf{p}})$ satisfying these relations
therefore implies first (taking $|\mathbf{p}|=1$) that the $dz_\lambda$ are
linearly independent, hence form a \emph{basis} of the
$\mathcal{O}_U$-module $\Omega_{U/S}^1$. Then, for every integer $m\geq1$, we
deduce in the same way from (16.11.2.3) that the
$\boldsymbol{\zeta}^{\mathbf{p}}$ such that $|\mathbf{p}|\leq m$ are linearly
independent; consequently the canonical homomorphism (16.3.1.1) is injective,
hence bijective, and this proves that b) implies a). The converse follows
immediately from the definition (16.10.1), the fact that the
$\boldsymbol{\zeta}^{\mathbf{p}}$ form a basis of $\mathscr{P}_{U/S}^m$ for
$|\mathbf{p}|\leq m$ implying the existence and uniqueness of a family of
homomorphisms $u_{\mathbf{q},m}:\mathscr{P}_{U/S}^m\to\mathcal{O}_U$
($|\mathbf{q}|\leq m$) such that
$\langle\boldsymbol{\zeta}^{\mathbf{p}},u_{\mathbf{q},m}\rangle
=\delta_{\mathbf{p}\mathbf{q}}$ for $|\mathbf{p}|\leq m$,
$|\mathbf{q}|\leq m$. For a given value of $\mathbf{q}$, the differential
operators corresponding to the $u_{\mathbf{q},m}$ for
$m\geq|\mathbf{q}|$ are identified with one and the same operator
$D_{\mathbf{q}}$. This therefore proves that a) implies b), and moreover that
the family $(D_{\mathbf{p}})$ is uniquely determined and that, if $L$ is
finite, for $|\mathbf{p}|\leq m$, the $D_{\mathbf{p}}$ form a basis of the
dual $\operatorname{Diff}_{U/S}^m$ of $\mathscr{P}_{U/S}^m$. Finally, the
relations (16.11.2.2) follow immediately from the expression of the values of
the three operators under consideration at the $\mathbf{z}^{\mathbf{r}}$, and
from the fact that the $\boldsymbol{\zeta}^{\mathbf{r}}$ for
$|\mathbf{r}|\leq m$ generate $\mathscr{P}_{U/S}^m$.

\oldpage[IV-4]{55}

\begin{remarks}[16.11.3]\label{IV.16.11.3}
\EGAReaderTarget{ega4.c04.statement.16-11-3}{IV.16.11.3}
\begin{enumerate}
  \item[(i)] The fact that the $D_{\mathbf{p}}$ pairwise commute by
  virtue of (16.11.2.2) does not of course imply that the
  $\mathcal{O}_U$-algebra $\operatorname{Diff}_{U/S}$ is commutative, since the
  $D_{\mathbf{p}}$ commute with multiplication by the sections of
  $\mathcal{O}_U$ only if $n=0$.
  \item[(ii)] The indices $\mathbf{p}$ such that $|\mathbf{p}|=1$ are the
  $\boldsymbol{\varepsilon}_\lambda
  =(\varepsilon_{\lambda\mu})_{\mu\in L}$ with
  $\varepsilon_{\lambda\mu}=0$ if $\mu\neq\lambda$ and
  $\varepsilon_{\lambda\lambda}=1$; when $L$ is finite, the operators
  $D_{\boldsymbol{\varepsilon}_\lambda}$ are none other than the
  $S$-derivations $D_i$ introduced in (16.5.7). Note that in general
  (and contrary to what happens in classical Analysis), it is not
  true that a differential operator of arbitrary order can be written as a
  linear combination of powers of the $D_i$ (cf. (16.12)).
  \item[(iii)] For every integer $r\geq1$, we can define the notion of a
  morphism \emph{differentially smooth up to order $r$} by replacing
  condition (ii) in (16.10.1) by the condition that the homomorphisms
  \[
    \mathbf{S}_{\mathcal{O}_X}^m(\Omega_{X/S}^1)
    \longrightarrow\mathscr{G}r_m(\mathscr{P}_{X/S})
  \]
  be bijective \emph{for every $m\leq r$}. The argument of (16.11.2)
  then proves that if, in condition a), we replace ``differentially
  smooth'' by ``differentially smooth up to order $r$'', this condition
  is equivalent to condition b) in which we restrict to the
  $\mathbf{p}\in\mathbb{N}^{(L)}$, $\mathbf{q}\in\mathbb{N}^{(L)}$ such that
  $|\mathbf{p}|\leq r$, $|\mathbf{q}|\leq r$.
\end{enumerate}
\end{remarks}

\subsection{The case of characteristic zero: Jacobian criterion for
differentially smooth morphisms}
\EGAReaderTarget{ega4.c04.subsec.16-12}{IV.16.12}

\begin{env}[16.12.1]\label{IV.16.12.1}
\EGAReaderTarget{ega4.c04.statement.16-12-1}{IV.16.12.1}
We say that a prescheme $X$ is \emph{of characteristic $p$} ($p$ equal to $0$
or to a prime number) if, for every affine open subset $U$ of $X$, the ring
$\Gamma(U,\mathcal{O}_X)$ is of characteristic $p$ (0, 21.1.1). It follows
from (0, 21.1.3) that for $X$ to be of characteristic $0$, it is necessary and
sufficient that for every \emph{closed} point $x$ of $X$ the residue field
$k(x)$ be of characteristic $0$, or also that $X$ can be endowed with a
structure of $\mathbb{Q}$-prescheme (necessarily unique).
\end{env}

\begin{theorem}[16.12.2]
\label{IV.16.12.2}
Let $X$ be a prescheme of characteristic $0$, $f:X\to S$ a morphism. If
$\Omega_{X/S}^1$ is a locally free $\mathcal{O}_X$-module (not
necessarily of finite type), then $f$ is differentially smooth.
\end{theorem}

The question being local on $X$, we may suppose that there exists a family
$(z_\lambda)$ of sections of $\mathcal{O}_X$ over $X$ such that
$(dz_\lambda)$ is a basis of the $\mathcal{O}_X$-module $\Omega_{X/S}^1$.
Applying the criterion (16.11.2), it suffices to remark that the operators
\[
  D_{\mathbf{p}}=(\mathbf{p}!)^{-1}
  \prod_\lambda D_\lambda^{p_\lambda}
\]
(where the $D_\lambda$ are the coordinate forms corresponding to the basis
$(dz_\lambda)$) satisfy the relations (16.11.2.1), which is a
consequence of the fact that the $D_\lambda$ are derivations.

\begin{env}[16.12.3]
\label{IV.16.12.3}
The preceding theorem is no longer true when we drop the hypothesis that
$X$ is of characteristic $0$. For example, if
$S=\operatorname{Spec}(k)$, where $k$ is a field of characteristic $p>0$,
and $X=\operatorname{Spec}(K)$ where $K=k(\alpha)$ with
$\alpha\notin k$, $\alpha^p\in k$, we verify immediately
\oldpage[IV-4]{56}
that $\Omega_{X/S}^1$ is of rank $1$, and that the morphism $X\to S$ is
differentially smooth up to order $p-1$ (16.11.3, (iii)), but not
up to order $p$. However, the proof of (16.12.2) shows that if
$\Omega_{X/S}^1$ is locally free, and if $n!1_{\mathcal{O}_X}$ is
invertible in $\Gamma(X,\mathcal{O}_X)$, then $X$ is differentially
smooth over $S$ up to order $n$.
\end{env}

